\documentclass[9pt]{extarticle}
\usepackage[a4paper,margin=1.8cm]{geometry}
\usepackage[utf8]{inputenc}
\usepackage[T1]{fontenc}
\usepackage{lmodern}
\usepackage{xcolor}
\usepackage{listings}
\usepackage[hidelinks]{hyperref}
\definecolor{cmt}{RGB}{96,139,84}
\definecolor{kw}{RGB}{0,0,170}
\definecolor{str}{RGB}{163,21,21}
\definecolor{ln}{RGB}{150,150,150}
\lstset{
  basicstyle=\ttfamily\scriptsize,
  breaklines=true,breakatwhitespace=false,
  columns=fullflexible,keepspaces=true,
  showstringspaces=false,upquote=true,tabsize=4,
  commentstyle=\color{cmt},keywordstyle=\color{kw},stringstyle=\color{str},
  numbers=left,numberstyle=\tiny\color{ln},numbersep=6pt,xleftmargin=2.2em,
  frame=lines,framesep=4pt,
}
\title{\textbf{C\'odigo Final}\\[2mm]\large Modelo Novo-Keynesiano de Pol\'itica
Monet\'aria para o Chile\\\normalsize Todo o c\'odigo (Python + Dynare) em um
\'unico documento}
\author{Daniel Colli e Jo\~ao Zangrandi}
\date{16/06/2026}
\begin{document}
\maketitle
{\small Este documento re\'une \textbf{todo} o c\'odigo do projeto: 49
arquivos (Python e Dynare \texttt{.mod}), agrupados na ordem do \emph{pipeline}.
Caracteres n\~ao-ASCII do c\'odigo (letras gregas, acentos) foram transliterados
para ASCII na listagem para garantir a compila\c{c}\~ao; a l\'ogica \'e id\^entica
\`a dos arquivos-fonte do reposit\'orio.\par}
\tableofcontents
\clearpage

\part*{Infraestrutura compartilhada}\addcontentsline{toc}{part}{Infraestrutura compartilhada}

\section{python/common.py}
\begin{lstlisting}[language=Python]
"""Shared paths and lightweight model-based fallback calculations."""

from __future__ import annotations

import math
import os
import shutil
from pathlib import Path

import numpy as np
import pandas as pd

ROOT = Path(__file__).resolve().parents[1]
DATA_RAW = ROOT / "data" / "raw"
DATA_CLEAN = ROOT / "data" / "clean"
TABLES = ROOT / "outputs" / "tables"
FIGURES = ROOT / "outputs" / "figures"
LOGS = ROOT / "outputs" / "logs"
DYNARE_OUTPUTS = ROOT / "outputs" / "dynare"

# Shock standard deviations for the Chile baseline. They are calibrated to an
# explicitly illustrative FEVD target at the estimated interest-rate persistence.
# The target is a modelling choice, not an empirical decomposition published by BCCh.
SHOCK_STD = {"e_x": 0.00500, "e_pi": 0.00277, "e_i": 0.00068}
INITIAL_SHOCK_STD = {"e_x": 0.0050, "e_pi": 0.0025, "e_i": 0.0015}
COURSE_BENCHMARK_SHOCK_STD_FALLBACK = {
    "e_x": 0.00500,
    "e_pi": 0.00331,
    "e_i": 0.00080,
}
BASELINE = {
    "rstar_annual": 0.03,
    "sigma": 1.0,
    "kappa": 0.10,
    "rho_i": 0.80,
    "phi_pi": 1.75,
    "phi_x": 0.50,
}


def ensure_directories() -> None:
    for path in (DATA_RAW, DATA_CLEAN, TABLES, FIGURES, LOGS, DYNARE_OUTPUTS):
        path.mkdir(parents=True, exist_ok=True)


def quarterly_rstar(annual_rate: float) -> float:
    return (1.0 + annual_rate) ** 0.25 - 1.0


def kalman_gap(log_level) -> tuple[np.ndarray, np.ndarray]:
    """Trend-cycle decomposition by an Unobserved-Components (Kalman) model.

    Replaces the HP filter everywhere in the project. The series (a log level) is
    split into a smooth stochastic trend and a stochastic, damped cycle; the cycle
    is the output-gap analogue. A smooth (integrated random-walk) trend keeps sharp
    shocks such as COVID-2020 in the CYCLE rather than the trend. Returns
    (cycle, trend) as numpy arrays of the same length as the input.
    """
    from statsmodels.tsa.statespace.structural import UnobservedComponents

    raw = np.asarray(log_level, dtype=float)
    y = raw * 100.0  # log points: far better-conditioned for the optimiser
    # An AR(2) cycle on a smooth (I(2)) trend is the classic Clark/Watson gap; the
    # cycle is constrained enough to capture the business cycle instead of letting
    # a flexible trend absorb it (which collapses the gap to ~0).
    attempts = (
        ("autoregressive", dict(level="smooth trend", autoregressive=2)),
        ("cycle", dict(level="smooth trend", cycle=True, stochastic_cycle=True,
                       damped_cycle=True, cycle_period_bounds=(6, 40))),
        ("autoregressive", dict(level="local linear trend", autoregressive=2)),
    )
    last_error: Exception | None = None
    for component, spec in attempts:
        try:
            result = UnobservedComponents(y, **spec).fit(disp=False, maxiter=3000)
            cycle = np.asarray(getattr(result, component).smoothed, dtype=float) / 100.0
            try:
                trend = np.asarray(result.level.smoothed, dtype=float) / 100.0
            except Exception:  # noqa: BLE001
                trend = raw - cycle
            if cycle.shape == raw.shape and np.all(np.isfinite(cycle)) and np.nanstd(cycle) > 1e-3:
                return cycle, trend
        except Exception as exc:  # noqa: BLE001
            last_error = exc
    raise RuntimeError(f"kalman_gap: no UnobservedComponents spec converged ({last_error}).")


def complete_parameters(overrides: dict | None = None) -> dict:
    params = dict(BASELINE)
    if overrides:
        params.update(overrides)
    params["rstar"] = quarterly_rstar(float(params["rstar_annual"]))
    params["beta"] = 1.0 / (1.0 + params["rstar"])
    return params


def calibration_rho() -> float:
    estimate_file = TABLES / "rhoi_estimate.csv"
    if not estimate_file.exists():
        return BASELINE["rho_i"]
    table = pd.read_csv(estimate_file)
    if table.empty or "rho_i_used" not in table:
        return BASELINE["rho_i"]
    value = float(table.loc[0, "rho_i_used"])
    return value if math.isfinite(value) else BASELINE["rho_i"]


def course_benchmark_shock_std() -> dict[str, float]:
    """Read shock sizes fitted to the didactic FEVD benchmark.

    The fitted values depend on the current rho_i calibration. The CSV is produced
    by calibrate_shocks.py; fixed fallbacks keep model generation runnable before
    that calibration script has been executed.
    """

    path = TABLES / "shock_sigmas_comparison.csv"
    if not path.exists():
        return dict(COURSE_BENCHMARK_SHOCK_STD_FALLBACK)
    table = pd.read_csv(path)
    selected = table[table["calibration"] == "didactic_benchmark"]
    if selected.empty:
        return dict(COURSE_BENCHMARK_SHOCK_STD_FALLBACK)
    values = {
        str(row["shock"]): float(row["sigma_calibrated"])
        for _, row in selected.iterrows()
    }
    return {
        shock: values.get(shock, fallback)
        for shock, fallback in COURSE_BENCHMARK_SHOCK_STD_FALLBACK.items()
    }


def find_octave() -> Path | None:
    configured = os.environ.get("OCTAVE_BIN")
    candidates = [
        configured,
        shutil.which("octave-cli"),
        shutil.which("octave"),
        r"C:\Program Files\GNU Octave\Octave-11.1.0\mingw64\bin\octave-cli.exe",
        r"C:\Program Files\GNU Octave\Octave-10.3.0\mingw64\bin\octave-cli.exe",
    ]
    for candidate in candidates:
        if candidate and Path(candidate).exists():
            return Path(candidate)
    octave_root = Path(r"C:\Program Files\GNU Octave")
    if octave_root.exists():
        matches = sorted(octave_root.glob(r"Octave-*\mingw64\bin\octave-cli.exe"))
        if matches:
            return matches[-1]
    return None


def find_dynare_matlab() -> Path | None:
    configured = os.environ.get("DYNARE_MATLAB")
    candidates = [
        configured,
        r"C:\dynare\7.1\matlab",
        r"C:\dynare\7.0\matlab",
        r"C:\Program Files\Dynare\7.1\matlab",
    ]
    for candidate in candidates:
        if candidate and (Path(candidate) / "dynare.m").exists():
            return Path(candidate)
    dynare_root = Path(r"C:\dynare")
    if dynare_root.exists():
        matches = sorted(dynare_root.glob(r"*\matlab\dynare.m"))
        if matches:
            return matches[-1].parent
    return None


def _state_coefficients(params: dict) -> tuple[float, float, float]:
    """Exact minimum-state-variable solution of the linear NK model.

    The only predetermined state is the lagged policy rate, so in the absence of
    shocks the dynamics collapse to a single stable eigenvalue a_i:
        x_t = a_x*i_{t-1},  pi_t = a_pi*i_{t-1},  i_t = a_i*i_{t-1}.
    a_i is the stable root (|a_i| < 1) of the system's characteristic cubic, which is
    unique under Blanchard-Kahn determinacy (one predetermined variable, two jumpers).
    """

    beta = params["beta"]
    sigma = params["sigma"]
    kappa = params["kappa"]
    rho_i = params["rho_i"]
    phi_pi = params["phi_pi"]
    phi_x = params["phi_x"]

    # Characteristic cubic P(a)=0 from a_x = a_x*a_i + ..., a_pi, and the Taylor rule:
    #   P(a) = (a-rho)*(beta a^2 - (1+beta+kappa/sigma) a + 1)
    #          + ((1-rho)/sigma)*((phi_pi*kappa + phi_x) a - phi_x*beta a^2)
    q2, q1, q0 = beta, -(1.0 + beta + kappa / sigma), 1.0
    c3 = q2
    c2 = q1 - rho_i * q2
    c1 = q0 - rho_i * q1
    c0 = -rho_i * q0
    g = (1.0 - rho_i) / sigma
    c2 += -g * phi_x * beta
    c1 += g * (phi_pi * kappa + phi_x)
    roots = np.roots([c3, c2, c1, c0])

    stable = [r.real for r in roots if abs(r.imag) < 1e-8 and abs(r) < 1.0 - 1e-10]
    if not stable:
        stable = [min(roots, key=lambda z: abs(z)).real]
    a_i = float(min(stable, key=lambda v: abs(v - rho_i)))

    determinant = (1.0 - a_i) * (1.0 - beta * a_i) - a_i * kappa / sigma
    a_x = (-a_i / sigma) * (1.0 - beta * a_i) / determinant
    a_pi = (-a_i * kappa / sigma) / determinant
    return a_x, a_pi, a_i


def characteristic_roots(params: dict) -> np.ndarray:
    """Eigenvalues (roots of the characteristic cubic) of the linear NK system.

    With one predetermined variable (the lagged rate) and two jumpers, Blanchard-Kahn
    determinacy requires exactly one root inside the unit circle.
    """
    beta = params["beta"]
    sigma = params["sigma"]
    kappa = params["kappa"]
    rho_i = params["rho_i"]
    phi_pi = params["phi_pi"]
    phi_x = params["phi_x"]
    q2, q1, q0 = beta, -(1.0 + beta + kappa / sigma), 1.0
    c3 = q2
    c2 = q1 - rho_i * q2
    c1 = q0 - rho_i * q1
    c0 = -rho_i * q0
    g = (1.0 - rho_i) / sigma
    c2 += -g * phi_x * beta
    c1 += g * (phi_pi * kappa + phi_x)
    return np.roots([c3, c2, c1, c0])


def fallback_irfs(
    params: dict, scenario: str, horizon: int = 20
) -> pd.DataFrame:
    """Generate model-based synthetic IRFs when Dynare output is unavailable."""

    a_x, a_pi, a_i = _state_coefficients(params)
    sigma = params["sigma"]
    rho_i = params["rho_i"]
    phi_pi = params["phi_pi"]
    phi_x = params["phi_x"]
    beta = params["beta"]

    c_i = a_x - (1.0 - a_pi) / sigma
    impact_matrix = np.array(
        [
            [1.0, 0.0, -c_i],
            [-params["kappa"], 1.0, -beta * a_pi],
            [
                -(1.0 - rho_i) * phi_x,
                -(1.0 - rho_i) * phi_pi,
                1.0,
            ],
        ]
    )

    records = []
    variables = ("x", "pi", "i")
    shocks = ("e_x", "e_pi", "e_i")
    for shock_index, shock in enumerate(shocks):
        innovation = np.zeros(3)
        innovation[shock_index] = SHOCK_STD[shock]
        impact = np.linalg.solve(impact_matrix, innovation)
        responses = [impact]
        lagged_rate = impact[2]
        for _ in range(1, horizon):
            current = np.array([a_x, a_pi, a_i]) * lagged_rate
            responses.append(current)
            lagged_rate = current[2]
        for period, response in enumerate(responses):
            for variable, value in zip(variables, response):
                records.append(
                    {
                        "scenario": scenario,
                        "source": "synthetic_model_fallback",
                        "variable": variable,
                        "shock": shock,
                        "horizon": period,
                        "response": float(value),
                    }
                )
    return pd.DataFrame.from_records(records)


def fallback_moments(params: dict, observations: int = 50000) -> pd.DataFrame:
    """Simulate deviation moments from the fallback solution."""

    a_x, a_pi, _ = _state_coefficients(params)
    sigma = params["sigma"]
    rho_i = params["rho_i"]
    beta = params["beta"]
    c_i = a_x - (1.0 - a_pi) / sigma
    impact_matrix = np.array(
        [
            [1.0, 0.0, -c_i],
            [-params["kappa"], 1.0, -beta * a_pi],
            [
                -(1.0 - rho_i) * params["phi_x"],
                -(1.0 - rho_i) * params["phi_pi"],
                1.0,
            ],
        ]
    )

    rng = np.random.default_rng(20260614)
    values = np.zeros((observations, 3))
    lagged_rate = 0.0
    std = np.array([SHOCK_STD["e_x"], SHOCK_STD["e_pi"], SHOCK_STD["e_i"]])
    for t in range(observations):
        rhs = rng.normal(0.0, std)
        rhs[2] += rho_i * lagged_rate
        values[t] = np.linalg.solve(impact_matrix, rhs)
        lagged_rate = values[t, 2]

    rows = []
    for index, variable in enumerate(("x", "pi", "i")):
        series = values[:, index]
        mean_value = float(series.mean())
        if variable == "i":
            mean_value += params["rstar"]
        rows.append(
            {
                "variable": variable,
                "mean": mean_value,
                "std_dev": float(series.std(ddof=1)),
                "variance": float(series.var(ddof=1)),
                "autocorrelation_1": float(np.corrcoef(series[1:], series[:-1])[0, 1]),
                "source": "synthetic_model_fallback",
            }
        )
    return pd.DataFrame(rows)
\end{lstlisting}
\clearpage
\part*{Dados (BCCh / FRED / filtro de Kalman)}\addcontentsline{toc}{part}{Dados (BCCh / FRED / filtro de Kalman)}

\section{python/build\_chile\_dataset.py}
\begin{lstlisting}[language=Python]
"""Build the real Chilean quarterly dataset from the Banco Central de Chile (BCCh).

Source: BCCh "Base de Datos Estadisticos" REST web service (SieteRestWS).
Series actually downloaded (confirmed available):
  - TPM  (monetary policy rate, daily, % p.a.):      F022.TPM.TIN.D001.NO.Z.D
  - IPC  (CPI monthly variation, %):                 F074.IPC.VAR.Z.Z.C.M
  - PIB  (real GDP, chained vol., seasonally adj.,
          spliced, reference 2018, quarterly):       F032.PIB.FLU.R.CLP.EP18.Z.Z.1.T

Credentials are read from the environment variables BCCH_USER / BCCH_PASS, or from a
local, git-ignored file data/raw/bcch_credentials.json with {"user": ..., "pass": ...}.
Credentials are never written to data/clean or printed.

If credentials are missing or the download fails, the script falls back to a clearly
labelled SYNTHETIC AR(1) series (for pipeline testing only) and flags is_synthetic=True.

Outputs (all in data/clean):
  - chile_policy_rate.csv      quarterly policy rate (% p.a.) for the AR(1) of rho_i
  - chile_macro_quarterly.csv  full quarterly panel (rate, inflation, gdp, gap, real rate)
  - chile_observables.csv      model-unit deviations (pi, i, x) for Dynare estimation
  - dataset_metadata.json      provenance (institution, series ids, units, dates, transforms)
"""

from __future__ import annotations

import json
import datetime as dt
from pathlib import Path

import numpy as np
import pandas as pd

from common import DATA_CLEAN, DATA_RAW, ensure_directories, kalman_gap

BCCH_BASE = "https://si3.bcentral.cl/SieteRestWS/SieteRestWS.ashx"
SERIES = {
    "tpm": "F022.TPM.TIN.D001.NO.Z.D",
    "ipc_var": "F074.IPC.VAR.Z.Z.C.M",
    "pib": "F032.PIB.FLU.R.CLP.EP18.Z.Z.1.T",
}
START_DATE = "2001-01-01"
INFLATION_TARGET_ANNUAL = 0.03  # Meta de inflacion del BCCh: 3% anual.
HP_LAMBDA = 1600  # Convencion trimestral.


# --------------------------------------------------------------------------- #
# Credentials and HTTP access
# --------------------------------------------------------------------------- #
def get_credentials() -> tuple[str, str] | None:
    import os

    user = os.environ.get("BCCH_USER")
    password = os.environ.get("BCCH_PASS")
    if user and password:
        return user, password
    creds_file = DATA_RAW / "bcch_credentials.json"
    if creds_file.exists():
        try:
            payload = json.loads(creds_file.read_text(encoding="utf-8"))
            user = payload.get("user")
            password = payload.get("pass") or payload.get("password")
            if user and password:
                return str(user), str(password)
        except Exception as exc:  # noqa: BLE001
            print(f"Warning: could not read bcch_credentials.json: {exc}")
    return None


def fetch_series(series_id: str, user: str, password: str,
                 first: str = START_DATE, last: str | None = None) -> pd.Series:
    """Download one BCCh series and return a float Series indexed by date."""
    import urllib.parse
    import urllib.request

    if last is None:
        last = dt.date.today().isoformat()
    query = urllib.parse.urlencode(
        {
            "user": user,
            "pass": password,
            "function": "GetSeries",
            "timeseries": series_id,
            "firstdate": first,
            "lastdate": last,
        }
    )
    url = f"{BCCH_BASE}?{query}"
    with urllib.request.urlopen(url, timeout=60) as response:
        payload = json.loads(response.read().decode("utf-8", "replace"))
    if payload.get("Codigo") != 0:
        raise RuntimeError(
            f"BCCh error for {series_id}: {payload.get('Descripcion')}"
        )
    series_block = payload.get("Series") or {}
    observations = series_block.get("Obs") or []
    dates, values = [], []
    for obs in observations:
        if str(obs.get("statusCode", "")).upper() not in ("OK", ""):
            continue
        raw_value = str(obs.get("value", "")).replace(",", ".").strip()
        if raw_value in ("", "NaN", "ND"):
            continue
        try:
            value = float(raw_value)
        except ValueError:
            continue
        dates.append(pd.to_datetime(obs["indexDateString"], format="%d-%m-%Y"))
        values.append(value)
    if not dates:
        raise RuntimeError(f"BCCh returned no usable observations for {series_id}.")
    return pd.Series(values, index=pd.DatetimeIndex(dates), name=series_id).sort_index()


# --------------------------------------------------------------------------- #
# Real dataset assembly
# --------------------------------------------------------------------------- #
def build_real_panel(user: str, password: str) -> tuple[pd.DataFrame, dict]:
    tpm_daily = fetch_series(SERIES["tpm"], user, password)
    ipc_var_m = fetch_series(SERIES["ipc_var"], user, password)
    pib_q = fetch_series(SERIES["pib"], user, password)

    # Policy rate: quarterly mean of the daily TPM (annual %).
    tpm_q = tpm_daily.groupby(tpm_daily.index.to_period("Q")).mean()

    # Inflation: compound three monthly CPI variations into a quarterly fraction.
    ipc_frac = ipc_var_m / 100.0
    grouped = ipc_frac.groupby(ipc_frac.index.to_period("Q"))
    infl_q = grouped.apply(lambda s: float(np.prod(1.0 + s.to_numpy()) - 1.0))
    months_per_quarter = grouped.size()
    infl_q = infl_q[months_per_quarter == 3]  # keep only complete quarters

    # GDP: real, seasonally adjusted, quarterly -> log -> HP cycle (output gap).
    pib_q.index = pib_q.index.to_period("Q")
    pib_q = pib_q[~pib_q.index.duplicated(keep="last")].sort_index()

    panel = pd.DataFrame(
        {
            "tpm_annual_pct": tpm_q,
            "infl_q": infl_q,
            "pib_sa": pib_q,
        }
    ).dropna()
    if len(panel) < 20:
        raise RuntimeError(
            f"Only {len(panel)} overlapping quarters; expected a long sample."
        )

    panel["log_pib"] = np.log(panel["pib_sa"])
    cycle, trend = kalman_gap(panel["log_pib"].to_numpy())  # Kalman UC, not HP
    panel["pib_trend_kalman"] = trend
    panel["output_gap"] = cycle

    # Quarterly nominal rate as a fraction (model unit), Fisher-consistent.
    panel["i_q"] = (1.0 + panel["tpm_annual_pct"] / 100.0) ** 0.25 - 1.0
    panel["infl_annual_pct"] = ((1.0 + panel["infl_q"]) ** 4 - 1.0) * 100.0
    panel["real_rate_q"] = panel["i_q"] - panel["infl_q"]

    panel.index = panel.index.astype(str)
    panel.index.name = "date"

    metadata = {
        "institution": "Banco Central de Chile (BCCh)",
        "service": "Base de Datos Estadisticos - SieteRestWS REST API",
        "url_pattern": f"{BCCH_BASE}?function=GetSeries&timeseries=<ID>&firstdate=...&lastdate=...",
        "official_bde_url": "https://si3.bcentral.cl/Siete",
        "policy_framework_url": "https://www.bcentral.cl/web/banco-central/areas/politica-monetaria",
        "terms_of_use_url": "https://www.bcentral.cl/web/banco-central/condiciones-de-uso",
        "series_ids": SERIES,
        "frequency": "quarterly (built from daily TPM, monthly CPI, quarterly GDP)",
        "period": f"{panel.index[0]} to {panel.index[-1]}",
        "observations": int(len(panel)),
        "access_date": dt.date.today().isoformat(),
        "units": {
            "tpm_annual_pct": "TPM, percent per year (quarterly average of daily series)",
            "i_q": "nominal policy rate, quarterly fraction = (1+TPM/100)^(1/4)-1",
            "infl_q": "quarterly CPI inflation, fraction (compounded monthly variations)",
            "infl_annual_pct": "annualised quarterly inflation, percent",
            "pib_sa": "real GDP, chained volume, seasonally adjusted, ref 2018 (CLP bn)",
            "output_gap": "HP-filter cycle of log real GDP (lambda=1600), fraction",
            "real_rate_q": "ex-post real quarterly rate = i_q - infl_q",
        },
        "transformations": (
            "TPM: daily->quarterly mean. CPI: monthly % -> quarterly compounded fraction. "
            f"GDP: log + HP filter (lambda={HP_LAMBDA}) for the output gap."
        ),
        "inflation_target_annual": INFLATION_TARGET_ANNUAL,
        "is_synthetic": False,
        "note": "Real official data from BCCh; credentials are not stored in this file.",
    }
    return panel, metadata


def synthetic_panel() -> tuple[pd.DataFrame, dict]:
    """Clearly-labelled synthetic fallback for pipeline testing only."""
    rng = np.random.default_rng(20260614)
    periods = pd.period_range("2001Q3", "2026Q1", freq="Q")
    long_run = 4.5
    rho = 0.92
    tpm = np.empty(len(periods))
    tpm[0] = long_run
    for i in range(1, len(periods)):
        tpm[i] = max(0.25, long_run + rho * (tpm[i - 1] - long_run) + rng.normal(0, 0.5))
    infl_q = 0.0074 + 0.5 * (tpm / 100.0 / 4 - 0.0074) + rng.normal(0, 0.004, len(periods))
    gap = np.zeros(len(periods))
    for i in range(1, len(periods)):
        gap[i] = 0.7 * gap[i - 1] + rng.normal(0, 0.01)
    panel = pd.DataFrame(
        {
            "tpm_annual_pct": tpm,
            "infl_q": infl_q,
            "pib_sa": 100.0 * np.exp(np.cumsum(0.005 + gap * 0.0)),
            "output_gap": gap,
        },
        index=periods.astype(str),
    )
    panel.index.name = "date"
    panel["log_pib"] = np.log(panel["pib_sa"])
    panel["pib_trend_hp"] = panel["log_pib"]
    panel["i_q"] = (1.0 + panel["tpm_annual_pct"] / 100.0) ** 0.25 - 1.0
    panel["infl_annual_pct"] = ((1.0 + panel["infl_q"]) ** 4 - 1.0) * 100.0
    panel["real_rate_q"] = panel["i_q"] - panel["infl_q"]
    metadata = {
        "institution": "SYNTHETIC (no real data)",
        "is_synthetic": True,
        "source": "synthetic_ar1_pipeline_test_not_chilean_data",
        "access_date": dt.date.today().isoformat(),
        "note": "No BCCh credentials available; generated a labelled synthetic panel.",
    }
    return panel, metadata


# --------------------------------------------------------------------------- #
def main() -> None:
    ensure_directories()
    credentials = get_credentials()
    if credentials is not None:
        try:
            panel, metadata = build_real_panel(*credentials)
            print(f"Downloaded real BCCh data: {metadata['period']} "
                  f"({metadata['observations']} quarters).")
        except Exception as exc:  # noqa: BLE001
            print(f"WARNING: BCCh download failed ({exc}). Using synthetic fallback.")
            panel, metadata = synthetic_panel()
    else:
        print("WARNING: no BCCh credentials (set BCCH_USER/BCCH_PASS or "
              "data/raw/bcch_credentials.json). Using synthetic fallback.")
        panel, metadata = synthetic_panel()

    is_synth = bool(metadata.get("is_synthetic", False))
    source_label = (
        "synthetic_ar1_pipeline_test_not_chilean_data"
        if is_synth
        else f"BCCh:{SERIES['tpm']}"
    )

    # 1) Policy-rate file consumed by the AR(1) estimation.
    policy = pd.DataFrame(
        {
            "date": panel.index,
            "policy_rate": panel["tpm_annual_pct"].to_numpy(),
            "is_synthetic": is_synth,
            "data_source": source_label,
        }
    )
    policy.to_csv(DATA_CLEAN / "chile_policy_rate.csv", index=False)

    # 2) Full quarterly macro panel.
    macro_cols = [
        "tpm_annual_pct", "i_q", "infl_q", "infl_annual_pct",
        "pib_sa", "log_pib", "pib_trend_hp", "output_gap", "real_rate_q",
    ]
    panel.reset_index()[["date"] + macro_cols].to_csv(
        DATA_CLEAN / "chile_macro_quarterly.csv", index=False
    )

    # 3) Model-unit observables (deviations from steady state) for estimation.
    observables = pd.DataFrame(
        {
            "date": panel.index,
            "pi": (panel["infl_q"] - ((1.0 + INFLATION_TARGET_ANNUAL) ** 0.25 - 1.0)).to_numpy(),
            "i": (panel["i_q"] - panel["i_q"].mean()).to_numpy(),
            "x": (panel["output_gap"] - panel["output_gap"].mean()).to_numpy(),
        }
    )
    observables.to_csv(DATA_CLEAN / "chile_observables.csv", index=False)

    metadata["files"] = {
        "policy_rate": "chile_policy_rate.csv",
        "macro_panel": "chile_macro_quarterly.csv",
        "observables": "chile_observables.csv (deviations: pi, i, x)",
    }
    (DATA_CLEAN / "dataset_metadata.json").write_text(
        json.dumps(metadata, indent=2, ensure_ascii=False), encoding="utf-8"
    )
    print(f"Wrote chile_policy_rate.csv, chile_macro_quarterly.csv, "
          f"chile_observables.csv ({len(panel)} quarters). synthetic={is_synth}")


if __name__ == "__main__":
    main()
\end{lstlisting}
\clearpage
\section{python/reprocess\_dataset.py}
\begin{lstlisting}[language=Python]
"""Re-derive the output gap (Kalman, not HP) and the model observables from the
cached BCCh panel, WITHOUT re-hitting the BCCh API.

Two methodological changes requested:
  * the output gap is an Unobserved-Components (Kalman) cycle, replacing the HP filter;
  * the inflation observable is centred on the 3% TARGET (not the sample mean), so the
    model's steady state is the target and forecasts converge to 3% by construction.

The rate stays demeaned by its sample mean (= realised neutral; Fisher-implied r*).
build_chile_dataset.py carries the same logic for a fresh pull, but it needs BCCh
credentials; this script works from the committed cache (data/clean/*.csv).
"""

from __future__ import annotations

import json

import numpy as np
import pandas as pd

from common import DATA_CLEAN, ensure_directories, kalman_gap

INFLATION_TARGET_ANNUAL = 0.03


def main() -> None:
    ensure_directories()
    panel = pd.read_csv(DATA_CLEAN / "chile_macro_quarterly.csv")

    # --- Output gap by Kalman (Unobserved Components) instead of HP ---
    cycle, trend = kalman_gap(panel["log_pib"].to_numpy())
    panel["output_gap_hp"] = panel["output_gap"]  # keep HP for reference
    panel["output_gap"] = cycle
    panel["pib_trend_kalman"] = trend
    panel.to_csv(DATA_CLEAN / "chile_macro_quarterly.csv", index=False)

    # --- Observables: inflation centred on the 3% target ---
    target_q = (1.0 + INFLATION_TARGET_ANNUAL) ** 0.25 - 1.0
    observables = pd.DataFrame(
        {
            "date": panel["date"],
            "pi": (panel["infl_q"] - target_q).to_numpy(),
            "i": (panel["i_q"] - panel["i_q"].mean()).to_numpy(),
            "x": (panel["output_gap"] - panel["output_gap"].mean()).to_numpy(),
        }
    )
    observables.to_csv(DATA_CLEAN / "chile_observables.csv", index=False)

    # --- Provenance ---
    meta_path = DATA_CLEAN / "dataset_metadata.json"
    meta = json.loads(meta_path.read_text(encoding="utf-8"))
    meta["output_gap_method"] = (
        "Unobserved-Components Kalman smoother (smooth I(2) trend + AR(2) cycle), "
        "replacing the HP filter; COVID stays in the cycle"
    )
    meta["inflation_centering"] = (
        "observable pi = quarterly inflation minus the 3% annual target "
        "(steady state = target; model forecasts converge to 3% by construction)"
    )
    meta_path.write_text(json.dumps(meta, indent=2, ensure_ascii=False), encoding="utf-8")

    mean_infl_annual = ((1.0 + panel["infl_q"].mean()) ** 4 - 1.0) * 100.0
    covid = float(cycle[panel["date"].to_numpy() == "2020Q2"][0]) * 100.0
    print(f"Kalman output gap: std {cycle.std()*100:.2f}% | COVID 2020Q2 "
          f"{covid:.1f}% | last {cycle[-1]*100:.2f}%")
    print(f"Inflation centred on 3% target (sample mean was {mean_infl_annual:.2f}%); "
          f"pi observable mean now {observables['pi'].mean()*400:.2f} annual p.p. above target")
    print("Wrote chile_macro_quarterly.csv, chile_observables.csv, dataset_metadata.json")


if __name__ == "__main__":
    main()
\end{lstlisting}
\clearpage
\section{python/build\_open\_economy\_dataset.py}
\begin{lstlisting}[language=Python]
"""Build the open-economy extension dataset for Chile.

Public monthly series are downloaded as CSV from FRED:
  - CCUSMA02CLM618N: CLP per USD, source OECD Main Economic Indicators.
  - PCOPPUSDM: global copper price, source IMF Primary Commodity Prices.

The raw/clean CSV is kept local by .gitignore. Public outputs contain only
derived summary statistics, calibration targets, metadata and a figure.
"""

from __future__ import annotations

import datetime as dt
import json
import urllib.request
from io import BytesIO

import matplotlib

matplotlib.use("Agg")
import matplotlib.pyplot as plt
import numpy as np
import pandas as pd
import statsmodels.api as sm
from statsmodels.tsa.filters.hp_filter import hpfilter

from common import DATA_CLEAN, FIGURES, TABLES, ensure_directories


SERIES = {
    "exchange_rate_clp_per_usd": "CCUSMA02CLM618N",
    "copper_usd_per_metric_ton": "PCOPPUSDM",
}
FRED_CSV = "https://fred.stlouisfed.org/graph/fredgraph.csv?id={series_id}"


def download_fred(series_id: str) -> pd.Series:
    request = urllib.request.Request(
        FRED_CSV.format(series_id=series_id),
        headers={"User-Agent": "nk-monetary-policy-chile/1.0"},
    )
    with urllib.request.urlopen(request, timeout=90) as response:
        payload = response.read()
    frame = pd.read_csv(BytesIO(payload))
    frame.columns = ["date", "value"]
    frame["date"] = pd.to_datetime(frame["date"])
    frame["value"] = pd.to_numeric(frame["value"], errors="coerce")
    return frame.dropna().set_index("date")["value"].sort_index()


def ar1(series: pd.Series) -> tuple[float, float]:
    y = series.iloc[1:].to_numpy()
    x = sm.add_constant(series.iloc[:-1].to_numpy())
    result = sm.OLS(y, x).fit()
    residual_std = float(np.std(result.resid, ddof=2))
    return float(result.params[1]), residual_std


def main() -> None:
    ensure_directories()
    fx_monthly = download_fred(SERIES["exchange_rate_clp_per_usd"])
    copper_monthly = download_fred(SERIES["copper_usd_per_metric_ton"])

    fx_q = fx_monthly.groupby(fx_monthly.index.to_period("Q")).mean()
    copper_q = copper_monthly.groupby(copper_monthly.index.to_period("Q")).mean()
    open_panel = pd.DataFrame({"fx_clp_usd": fx_q, "copper_usd_ton": copper_q})
    macro = pd.read_csv(DATA_CLEAN / "chile_macro_quarterly.csv")
    macro.index = pd.PeriodIndex(macro["date"], freq="Q")
    open_panel = open_panel.join(
        macro[["output_gap", "infl_q", "i_q", "tpm_annual_pct"]],
        how="inner",
    ).dropna()

    open_panel["log_fx"] = np.log(open_panel["fx_clp_usd"])
    open_panel["log_copper"] = np.log(open_panel["copper_usd_ton"])
    from common import kalman_gap  # Kalman UC cycle, not HP
    fx_cycle, fx_trend = kalman_gap(open_panel["log_fx"].to_numpy())
    copper_cycle, copper_trend = kalman_gap(open_panel["log_copper"].to_numpy())
    open_panel["fx_gap"] = fx_cycle
    open_panel["fx_trend"] = fx_trend
    open_panel["copper_gap"] = copper_cycle
    open_panel["copper_trend"] = copper_trend
    open_panel["fx_depreciation_q"] = open_panel["log_fx"].diff()

    rho_copper, std_copper = ar1(open_panel["copper_gap"])
    rho_fx, std_fx_gap = ar1(open_panel["fx_gap"])
    std_fx_depreciation = float(open_panel["fx_depreciation_q"].dropna().std(ddof=1))

    calibration = pd.DataFrame(
        [
            {
                "rho_copper": float(np.clip(rho_copper, 0.20, 0.95)),
                "rho_fx_gap_descriptive": rho_fx,
                "std_e_copper": float(np.clip(std_copper, 0.01, 0.12)),
                "std_e_fx": float(np.clip(std_fx_gap, 0.01, 0.10)),
                "std_fx_depreciation_q": std_fx_depreciation,
                "alpha_q_is": 0.05,
                "alpha_copper_is": 0.03,
                "gamma_q_phillips": 0.03,
                "eta_i_uip": 1.0,
                "phi_q_taylor": 0.05,
                "calibration_scope": "illustrative open-economy extension",
            }
        ]
    )
    calibration.to_csv(TABLES / "open_economy_calibration.csv", index=False)

    summary_rows = []
    for column in [
        "fx_clp_usd",
        "copper_usd_ton",
        "fx_gap",
        "copper_gap",
        "fx_depreciation_q",
    ]:
        values = open_panel[column].dropna()
        summary_rows.append(
            {
                "series": column,
                "observations": len(values),
                "mean": float(values.mean()),
                "std_dev": float(values.std(ddof=1)),
                "minimum": float(values.min()),
                "maximum": float(values.max()),
                "sample_start": str(values.index[0]),
                "sample_end": str(values.index[-1]),
            }
        )
    pd.DataFrame(summary_rows).to_csv(
        TABLES / "open_economy_data_summary.csv", index=False
    )

    clean = open_panel.reset_index(names="date")
    clean["date"] = clean["date"].astype(str)
    clean.to_csv(DATA_CLEAN / "chile_open_economy.csv", index=False)
    metadata = {
        "access_date": dt.date.today().isoformat(),
        "is_synthetic": False,
        "frequency": "quarterly averages built from monthly series",
        "series": {
            "exchange_rate": {
                "id": SERIES["exchange_rate_clp_per_usd"],
                "title": "CLP per USD, average daily rate",
                "source": "OECD Main Economic Indicators via FRED",
                "url": "https://fred.stlouisfed.org/series/CCUSMA02CLM618N",
            },
            "copper": {
                "id": SERIES["copper_usd_per_metric_ton"],
                "title": "Global price of copper, USD per metric ton",
                "source": "IMF Primary Commodity Prices via FRED",
                "url": "https://fred.stlouisfed.org/series/PCOPPUSDM",
            },
        },
        "transformations": (
            "Monthly levels to quarterly averages; natural logs; HP-filter cycles "
            "with lambda=1600. Calibration uses AR(1) persistence and innovation std."
        ),
        "warning": (
            "Open-economy parameters are illustrative and are not jointly estimated "
            "structural coefficients."
        ),
    }
    (DATA_CLEAN / "open_economy_metadata.json").write_text(
        json.dumps(metadata, ensure_ascii=False, indent=2), encoding="utf-8"
    )

    x = np.arange(len(open_panel))
    ticks = np.arange(0, len(open_panel), 8)
    labels = [str(open_panel.index[index]) for index in ticks]
    fig, axes = plt.subplots(2, 2, figsize=(13, 8), sharex=True)
    axes[0, 0].plot(x, open_panel["fx_clp_usd"], color="#1E4E79")
    axes[0, 0].set_title("Cambio nominal: pesos chilenos por US$ 1")
    axes[0, 0].set_ylabel("CLP/USD")
    axes[0, 1].plot(x, open_panel["copper_usd_ton"], color="#B45309")
    axes[0, 1].set_title("Preco global do cobre")
    axes[0, 1].set_ylabel("US$/tonelada")
    axes[1, 0].plot(x, 100 * open_panel["fx_gap"], color="#1E4E79")
    axes[1, 0].axhline(0, color="black", linewidth=0.7)
    axes[1, 0].set_title("Gap cambial HP (alta = peso depreciado)")
    axes[1, 0].set_ylabel("% log")
    axes[1, 1].plot(x, 100 * open_panel["copper_gap"], color="#B45309")
    axes[1, 1].axhline(0, color="black", linewidth=0.7)
    axes[1, 1].set_title("Gap do preco do cobre HP")
    axes[1, 1].set_ylabel("% log")
    for ax in axes.ravel():
        ax.grid(alpha=0.25)
        ax.set_xticks(ticks)
        ax.set_xticklabels(labels, rotation=70, fontsize=8)
    fig.suptitle("Chile como economia pequena e aberta: cambio e cobre")
    fig.text(
        0.01,
        0.01,
        "Fontes: OECD (cambio) e FMI (cobre), distribuidas pelo FRED. "
        "Gaps HP sao proxies ciclicas, nao valores de equilibrio observados.",
        fontsize=8,
    )
    fig.tight_layout(rect=(0, 0.03, 1, 0.95))
    fig.savefig(FIGURES / "open_economy_data.png", dpi=180)
    plt.close(fig)
    print(calibration.to_string(index=False))
    print(f"Wrote open-economy data for {len(open_panel)} quarters.")


if __name__ == "__main__":
    main()
\end{lstlisting}
\clearpage
\section{python/discover\_ipc\_series.py}
\begin{lstlisting}[language=Python]
"""Discover BCCh IPC core/sub-index series IDs via SieteRestWS SearchSeries.

Run once credentials are available (env BCCH_USER/BCCH_PASS or
data/raw/bcch_credentials.json). Prints every MONTHLY series whose Spanish title
mentions the CPI cores we need (sin volatiles, servicios, bienes, energia,
alimentos, volatil), so we can lock the exact ``F074.IPC...`` ids for the puller.

Credentials are never printed.
"""

from __future__ import annotations

import json
import urllib.parse
import urllib.request

from build_chile_dataset import BCCH_BASE, get_credentials

KEYWORDS = (
    "ipc", "volatil", "volatil", "sin volatiles", "sin volatiles",
    "servicio", "bien", "energ", "aliment", "subyacente", "nucleo", "nucleo",
)


def search(frequency: str, user: str, password: str) -> list[dict]:
    query = urllib.parse.urlencode(
        {"user": user, "pass": password, "function": "SearchSeries",
         "frequency": frequency}
    )
    url = f"{BCCH_BASE}?{query}"
    with urllib.request.urlopen(url, timeout=90) as response:
        payload = json.loads(response.read().decode("utf-8", "replace"))
    return payload.get("SeriesInfos") or []


def main() -> None:
    creds = get_credentials()
    if not creds:
        raise SystemExit("No BCCh credentials (set BCCH_USER/BCCH_PASS or "
                         "data/raw/bcch_credentials.json).")
    user, password = creds
    infos = search("MONTHLY", user, password)
    print(f"MONTHLY series returned: {len(infos)}")
    hits = []
    for info in infos:
        sid = info.get("seriesId", "")
        title = (info.get("spanishTitle") or info.get("frequencyCode") or "")
        low = f"{sid} {title}".lower()
        if sid.startswith("F074.IPC") or ("ipc" in low and
                                          any(k in low for k in KEYWORDS)):
            hits.append((sid, title))
    print(f"IPC-related monthly series: {len(hits)}")
    for sid, title in sorted(hits):
        print(f"  {sid:32s} | {title}")


if __name__ == "__main__":
    main()
\end{lstlisting}
\clearpage
\part*{Geracao dos modelos Dynare}\addcontentsline{toc}{part}{Geracao dos modelos Dynare}

\section{python/generate\_dynare\_models.py}
\begin{lstlisting}[language=Python]
"""Generate Dynare sensitivity models from the baseline specification."""

from __future__ import annotations

from pathlib import Path

import pandas as pd

from common import (
    BASELINE,
    ROOT,
    SHOCK_STD,
    TABLES,
    calibration_rho,
    complete_parameters,
    course_benchmark_shock_std,
    ensure_directories,
)

TEMPLATE = """// Auto-generated scenario: {scenario}
var x pi i;
varexo e_x e_pi e_i;
parameters beta sigma kappa rho_i phi_pi phi_x rstar;

beta = {beta:.12g};
sigma = {sigma:.12g};
kappa = {kappa:.12g};
rho_i = {rho_i:.12g};
phi_pi = {phi_pi:.12g};
phi_x = {phi_x:.12g};
rstar = {rstar:.12g};

model(linear);
x  = x(+1) - (1/sigma)*(i - pi(+1) - rstar) + e_x;
pi = beta*pi(+1) + kappa*x + e_pi;
i  = rho_i*i(-1) + (1-rho_i)*(rstar + phi_pi*pi + phi_x*x) + e_i;
end;

initval;
x = 0;
pi = 0;
i = rstar;
end;

shocks;
var e_x;  stderr {std_e_x:.12g};
var e_pi; stderr {std_e_pi:.12g};
var e_i;  stderr {std_e_i:.12g};
end;

steady;
check;
verbatim;
export_stability('{scenario}');
end;
stoch_simul(order=1, irf=20);
verbatim;
export_results('{scenario}');
end;
"""


def scenario_name(prefix: str, value: float, digits: int = 2) -> str:
    text = f"{value:.{digits}f}".replace(".", "p")
    return f"{prefix}_{text}"


def main() -> None:
    ensure_directories()
    generated = ROOT / "dynare" / "generated"
    generated.mkdir(parents=True, exist_ok=True)

    rho_i = calibration_rho()
    baseline_shocks = {
        "std_e_x": SHOCK_STD["e_x"],
        "std_e_pi": SHOCK_STD["e_pi"],
        "std_e_i": SHOCK_STD["e_i"],
    }

    # Regenerate the baseline .mod so it stays consistent with the estimated rho_i
    # and the current shock calibration (single source of truth for the baseline).
    baseline_params = complete_parameters({"rho_i": rho_i})
    (ROOT / "dynare" / "nk_chile_base.mod").write_text(
        TEMPLATE.format(
            scenario="baseline",
            **baseline_shocks,
            **baseline_params,
        ),
        encoding="ascii",
    )

    scenarios: list[dict] = []
    for annual in (0.02, 0.03, 0.04):
        params = complete_parameters({"rstar_annual": annual, "rho_i": rho_i})
        scenarios.append(
            {
                "scenario": scenario_name("rstar", annual * 100, 0),
                "scenario_type": "rstar",
                **baseline_shocks,
                **params,
            }
        )
    for kappa in (0.07, 0.10, 0.13):
        params = complete_parameters({"kappa": kappa, "rho_i": rho_i})
        scenarios.append(
            {
                "scenario": scenario_name("kappa", kappa, 2),
                "scenario_type": "kappa",
                **baseline_shocks,
                **params,
            }
        )
    for phi_pi in [round(value, 1) for value in
                   [1.3 + 0.1 * index for index in range(10)]]:
        params = complete_parameters({"phi_pi": phi_pi, "rho_i": rho_i})
        scenarios.append(
            {
                "scenario": scenario_name("phi_pi", phi_pi, 1),
                "scenario_type": "phi_pi",
                **baseline_shocks,
                **params,
            }
        )

    # "Calibrate or estimate rho_i": retain the empirical baseline and add the
    # lecture-style calibrated route as a clean one-parameter robustness case.
    scenarios.append(
        {
            "scenario": "rho_calibrated_0p80",
            "scenario_type": "rho_i",
            **baseline_shocks,
            **complete_parameters({"rho_i": BASELINE["rho_i"]}),
        }
    )

    # Alternative shock calibration fitted to the didactic FEVD benchmark,
    # holding all structural parameters at the Chilean empirical baseline.
    benchmark_shocks = course_benchmark_shock_std()
    scenarios.append(
        {
            "scenario": "fevd_didactic_benchmark",
            "scenario_type": "fevd",
            "std_e_x": benchmark_shocks["e_x"],
            "std_e_pi": benchmark_shocks["e_pi"],
            "std_e_i": benchmark_shocks["e_i"],
            **baseline_params,
        }
    )

    for row in scenarios:
        content = TEMPLATE.format(**row)
        (generated / f"{row['scenario']}.mod").write_text(
            content, encoding="ascii"
        )

    expected_models = {f"{row['scenario']}.mod" for row in scenarios}
    for old_model in generated.glob("*.mod"):
        if old_model.name not in expected_models:
            try:
                old_model.unlink()
            except PermissionError:
                print(
                    f"Warning: OneDrive locked obsolete model {old_model.name}; "
                    "remove it manually before publishing."
                )

    manifest_columns = [
        "scenario",
        "scenario_type",
        "rstar_annual",
        "rstar",
        "beta",
        "sigma",
        "kappa",
        "rho_i",
        "phi_pi",
        "phi_x",
        "std_e_x",
        "std_e_pi",
        "std_e_i",
    ]
    pd.DataFrame(scenarios)[manifest_columns].to_csv(
        TABLES / "scenario_manifest.csv", index=False
    )
    print(f"Generated {len(scenarios)} models in {generated}")


if __name__ == "__main__":
    main()
\end{lstlisting}
\clearpage
\section{python/generate\_macro\_extension\_models.py}
\begin{lstlisting}[language=Python]
"""Generate Dynare files for the open-economy and hybrid-NKPC extensions."""

from __future__ import annotations

import pandas as pd

from common import ROOT, SHOCK_STD, TABLES, calibration_rho, complete_parameters


OPEN_TEMPLATE = """// Illustrative small-open-economy extension for Chile.
var x pi i q copper;
varexo e_x e_pi e_i e_q e_copper;
parameters beta sigma kappa rho_i phi_pi phi_x alpha_q alpha_copper
           gamma_q eta_i phi_q rho_copper;

beta = {beta:.12g};
sigma = {sigma:.12g};
kappa = {kappa:.12g};
rho_i = {rho_i:.12g};
phi_pi = {phi_pi:.12g};
phi_x = {phi_x:.12g};
alpha_q = {alpha_q:.12g};
alpha_copper = {alpha_copper:.12g};
gamma_q = {gamma_q:.12g};
eta_i = {eta_i:.12g};
phi_q = {phi_q:.12g};
rho_copper = {rho_copper:.12g};

model(linear);
  // Domestic demand: real rate, competitiveness and copper-income channel.
  x = x(+1) - (1/sigma)*(i-pi(+1)) + alpha_q*q + alpha_copper*copper + e_x;
  // Domestic inflation with direct exchange-rate pass-through.
  pi = beta*pi(+1) + kappa*x + gamma_q*q + e_pi;
  // Taylor rule may lean mildly against depreciation.
  i = rho_i*i(-1) + (1-rho_i)*(phi_pi*pi + phi_x*x + phi_q*q) + e_i;
  // UIP-style exchange-rate equation; q>0 means a weaker peso.
  q = q(+1) - eta_i*i + e_q;
  // Exogenous persistent copper-price gap.
  copper = rho_copper*copper(-1) + e_copper;
end;

steady;
check;
verbatim;
  export_stability('open_economy');
end;

shocks;
  var e_x; stderr {std_e_x:.12g};
  var e_pi; stderr {std_e_pi:.12g};
  var e_i; stderr {std_e_i:.12g};
  var e_q; stderr {std_e_q:.12g};
  var e_copper; stderr {std_e_copper:.12g};
end;

stoch_simul(order=1, irf=20, nograph);
verbatim;
  export_extension('open_economy');
end;
"""


HYBRID_TEMPLATE = """// Hybrid Phillips-curve extension with intrinsic inflation persistence.
var x pi i;
varexo e_x e_pi e_i;
parameters beta sigma kappa rho_i phi_pi phi_x gamma_pi;

beta = {beta:.12g};
sigma = {sigma:.12g};
kappa = {kappa:.12g};
rho_i = {rho_i:.12g};
phi_pi = {phi_pi:.12g};
phi_x = {phi_x:.12g};
gamma_pi = {gamma_pi:.12g};

model(linear);
  x = x(+1) - (1/sigma)*(i-pi(+1)) + e_x;
  pi = gamma_pi*pi(-1) + (1-gamma_pi)*(beta*pi(+1)+kappa*x) + e_pi;
  i = rho_i*i(-1) + (1-rho_i)*(phi_pi*pi+phi_x*x) + e_i;
end;

steady;
check;
verbatim;
  export_stability('hybrid_nkpc');
end;

shocks;
  var e_x; stderr {std_e_x:.12g};
  var e_pi; stderr {std_e_pi:.12g};
  var e_i; stderr {std_e_i:.12g};
end;

stoch_simul(order=1, irf=20, nograph);
verbatim;
  export_extension('hybrid_nkpc');
end;
"""


def main() -> None:
    params = complete_parameters({"rho_i": calibration_rho()})
    calibration_path = TABLES / "open_economy_calibration.csv"
    if calibration_path.exists():
        open_cal = pd.read_csv(calibration_path).iloc[0].to_dict()
    else:
        open_cal = {
            "rho_copper": 0.75,
            "std_e_copper": 0.08,
            "std_e_fx": 0.04,
            "alpha_q_is": 0.05,
            "alpha_copper_is": 0.03,
            "gamma_q_phillips": 0.03,
            "eta_i_uip": 1.0,
            "phi_q_taylor": 0.05,
        }
    open_values = {
        **params,
        "alpha_q": float(open_cal["alpha_q_is"]),
        "alpha_copper": float(open_cal["alpha_copper_is"]),
        "gamma_q": float(open_cal["gamma_q_phillips"]),
        "eta_i": float(open_cal["eta_i_uip"]),
        "phi_q": float(open_cal["phi_q_taylor"]),
        "rho_copper": float(open_cal["rho_copper"]),
        "std_e_x": SHOCK_STD["e_x"],
        "std_e_pi": SHOCK_STD["e_pi"],
        "std_e_i": SHOCK_STD["e_i"],
        "std_e_q": float(open_cal["std_e_fx"]),
        "std_e_copper": float(open_cal["std_e_copper"]),
    }
    (ROOT / "dynare" / "nk_chile_open.mod").write_text(
        OPEN_TEMPLATE.format(**open_values), encoding="ascii"
    )

    # Backward OLS estimate is around 0.34; use 0.35 as an illustrative
    # intrinsic-persistence share, not as a structural estimate of indexation.
    hybrid_values = {
        **params,
        "gamma_pi": 0.35,
        "std_e_x": SHOCK_STD["e_x"],
        "std_e_pi": SHOCK_STD["e_pi"],
        "std_e_i": SHOCK_STD["e_i"],
    }
    (ROOT / "dynare" / "nk_chile_hybrid.mod").write_text(
        HYBRID_TEMPLATE.format(**hybrid_values), encoding="ascii"
    )
    pd.DataFrame(
        [
            {"scenario": "open_economy", **open_values},
            {"scenario": "hybrid_nkpc", **hybrid_values},
        ]
    ).to_csv(TABLES / "macro_extension_calibration.csv", index=False)
    print("Generated nk_chile_open.mod and nk_chile_hybrid.mod")


if __name__ == "__main__":
    main()
\end{lstlisting}
\clearpage
\part*{Modelos Dynare (.mod)}\addcontentsline{toc}{part}{Modelos Dynare (.mod)}

\section{dynare/nk\_chile\_base.mod}
\begin{lstlisting}[language={}]
// Auto-generated scenario: baseline
var x pi i;
varexo e_x e_pi e_i;
parameters beta sigma kappa rho_i phi_pi phi_x rstar;

beta = 0.992637536145;
sigma = 1;
kappa = 0.1;
rho_i = 0.934095195259;
phi_pi = 1.75;
phi_x = 0.5;
rstar = 0.00741707177773;

model(linear);
x  = x(+1) - (1/sigma)*(i - pi(+1) - rstar) + e_x;
pi = beta*pi(+1) + kappa*x + e_pi;
i  = rho_i*i(-1) + (1-rho_i)*(rstar + phi_pi*pi + phi_x*x) + e_i;
end;

initval;
x = 0;
pi = 0;
i = rstar;
end;

shocks;
var e_x;  stderr 0.005;
var e_pi; stderr 0.00277;
var e_i;  stderr 0.00068;
end;

steady;
check;
verbatim;
export_stability('baseline');
end;
stoch_simul(order=1, irf=20);
verbatim;
export_results('baseline');
end;
\end{lstlisting}
\clearpage
\section{dynare/nk\_chile\_estim.mod}
\begin{lstlisting}[language={}]
// Posterior-mode estimation of the minimal New Keynesian model on Chilean data.
// Written in pure deviation form (r* drops out), so the steady state is x=pi=i=0 and
// matches the demeaned observables (pi, i, x are quarterly fractions, mean-zero).
// Three observables and three structural shocks => no stochastic singularity.

var x pi i;
varexo e_x e_pi e_i;
parameters beta sigma kappa rho_i phi_pi phi_x;

beta   = 0.992637536145;   // 1/(1+r*), r* from the 3% annual scenario
sigma  = 1.0;
kappa  = 0.10;
rho_i  = 0.90;
phi_pi = 1.75;
phi_x  = 0.50;

model(linear);
x  = x(+1) - (1/sigma)*(i - pi(+1)) + e_x;
pi = beta*pi(+1) + kappa*x + e_pi;
i  = rho_i*i(-1) + (1-rho_i)*(phi_pi*pi + phi_x*x) + e_i;
end;

steady;

varobs pi i x;

estimated_params;
  sigma,       gamma_pdf,     1.00,  0.40;
  kappa,       gamma_pdf,     0.10,  0.05;
  rho_i,       beta_pdf,      0.80,  0.10;
  phi_pi,      gamma_pdf,     1.50,  0.25;
  phi_x,       gamma_pdf,     0.50,  0.25;
  stderr e_x,  inv_gamma_pdf, 0.005, inf;
  stderr e_pi, inv_gamma_pdf, 0.0025, inf;
  stderr e_i,  inv_gamma_pdf, 0.002, inf;
end;

// Posterior-mode (MAP) estimation: fast and robust in Octave. Standard errors come from
// the mode Hessian (recovered in Python). mh_replic can be raised to add a full MCMC.
estimation(datafile='chile_observables_dynare.csv', mode_compute=4,
           mh_replic=0, plot_priors=0, nograph);

verbatim;
export_bayesian('estimation');
end;
\end{lstlisting}
\clearpage
\section{dynare/nk\_chile\_hybrid.mod}
\begin{lstlisting}[language={}]
// Hybrid Phillips-curve extension with intrinsic inflation persistence.
var x pi i;
varexo e_x e_pi e_i;
parameters beta sigma kappa rho_i phi_pi phi_x gamma_pi;

beta = 0.992637536145;
sigma = 1;
kappa = 0.1;
rho_i = 0.934095195259;
phi_pi = 1.75;
phi_x = 0.5;
gamma_pi = 0.35;

model(linear);
  x = x(+1) - (1/sigma)*(i-pi(+1)) + e_x;
  pi = gamma_pi*pi(-1) + (1-gamma_pi)*(beta*pi(+1)+kappa*x) + e_pi;
  i = rho_i*i(-1) + (1-rho_i)*(phi_pi*pi+phi_x*x) + e_i;
end;

steady;
check;
verbatim;
  export_stability('hybrid_nkpc');
end;

shocks;
  var e_x; stderr 0.005;
  var e_pi; stderr 0.00277;
  var e_i; stderr 0.00068;
end;

stoch_simul(order=1, irf=20, nograph);
verbatim;
  export_extension('hybrid_nkpc');
end;
\end{lstlisting}
\clearpage
\section{dynare/nk\_chile\_hybrid\_estim.mod}
\begin{lstlisting}[language={}]
// Posterior-mode estimation of the hybrid-Phillips Chilean NK model.
// This model uses the same observables and common priors as nk_chile_estim.mod,
// adding one estimated indexation parameter gamma_pi. The matched data sample
// makes its marginal likelihood comparable with the forward-looking baseline.

var x pi i;
varexo e_x e_pi e_i;
parameters beta sigma kappa rho_i phi_pi phi_x gamma_pi;

beta     = 0.992637536145;
sigma    = 1.0;
kappa    = 0.10;
rho_i    = 0.90;
phi_pi   = 1.75;
phi_x    = 0.50;
gamma_pi = 0.35;

model(linear);
  x  = x(+1) - (1/sigma)*(i-pi(+1)) + e_x;
  pi = gamma_pi*pi(-1)
       + (1-gamma_pi)*(beta*pi(+1)+kappa*x) + e_pi;
  i  = rho_i*i(-1) + (1-rho_i)*(phi_pi*pi+phi_x*x) + e_i;
end;

steady;
check;
varobs pi i x;

estimated_params;
  sigma,       gamma_pdf,     1.00,  0.40;
  kappa,       gamma_pdf,     0.10,  0.05;
  rho_i,       beta_pdf,      0.80,  0.10;
  phi_pi,      gamma_pdf,     1.50,  0.25;
  phi_x,       gamma_pdf,     0.50,  0.25;
  gamma_pi,    beta_pdf,      0.35,  0.15;
  stderr e_x,  inv_gamma_pdf, 0.005, inf;
  stderr e_pi, inv_gamma_pdf, 0.0025, inf;
  stderr e_i,  inv_gamma_pdf, 0.002, inf;
end;

estimation(datafile='chile_observables_dynare.csv', mode_compute=4,
           mh_replic=0, plot_priors=0, nograph, nodisplay);
\end{lstlisting}
\clearpage
\section{dynare/nk\_chile\_forecast.mod}
\begin{lstlisting}[language={}]
// Native Dynare forecast for the calibrated Chilean New Keynesian model.
// Implements the "projections from observed data" workflow: read quarterly
// observables, keep the calibration fixed (no Bayesian estimation), and project
// pi, i and x eight quarters ahead, both unconditionally and under controlled
// interest-rate and inflation paths.
//
// The model is written in pure deviation form (r* drops out), so the steady
// state is x=pi=i=0 and matches the demeaned observables. Three observables and
// three structural shocks => no stochastic singularity.

var x pi i;
varexo e_x e_pi e_i;
parameters beta sigma kappa rho_i phi_pi phi_x;

// Baseline calibration (same values used by the stoch_simul baseline):
beta   = 0.992637536145;   // 1/(1+r*), r* from the 3% annual scenario
sigma  = 1.0;              // log-utility benchmark
kappa  = 0.10;             // Phillips-curve slope
rho_i  = 0.934095195259;   // interest-rate smoothing, AR(1) of the quarterly TPM
phi_pi = 1.75;             // inflation response (Taylor principle satisfied)
phi_x  = 0.50;             // output-gap response

model(linear);
  // IS (expectational), deviation form
  x  = x(+1) - (1/sigma)*(i - pi(+1)) + e_x;
  // New Keynesian Phillips curve
  pi = beta*pi(+1) + kappa*x + e_pi;
  // Taylor rule with interest-rate smoothing
  i  = rho_i*i(-1) + (1-rho_i)*(phi_pi*pi + phi_x*x) + e_i;
end;

steady;
check;

// Calibrated shock sizes (same as the baseline stoch_simul):
shocks;
  var e_x;  stderr 0.005;
  var e_pi; stderr 0.00277;
  var e_i;  stderr 0.00068;
end;

varobs pi i x;

// The three shock standard deviations are declared so the estimation machinery
// has something to "estimate", but mode_compute=0 together with
// estimated_params_init(use_calibration) means nothing is optimised: every
// parameter stays at its calibrated value. The priors below are never used.
estimated_params;
  stderr e_x,  inv_gamma_pdf, 0.005,   inf;
  stderr e_pi, inv_gamma_pdf, 0.00277, inf;
  stderr e_i,  inv_gamma_pdf, 0.00068, inf;
end;

estimated_params_init(use_calibration);
  stderr e_x,  0.005;
  stderr e_pi, 0.00277;
  stderr e_i,  0.00068;
end;

// Kalman smoother over the history + 8-quarter unconditional forecast, at the
// fixed calibration. forecast=8 fills oo_.forecast.Mean / HPDinf / HPDsup.
estimation(datafile='chile_observables_dynare.csv', mode_compute=0, mh_replic=0,
           smoother, forecast=8, nograph, nodisplay) x pi i;

verbatim;
  export_forecast('forecast');
end;

// ----------------------------------------------------------------------------
// Conditional projections. The controlled exogenous variable is the monetary
// policy shock e_i, which is solved for so that the constrained path is hit.
// Values are in model (deviation) units; the steady state is zero.
// ----------------------------------------------------------------------------

// Scenario 1: hold the policy-rate gap at its last observed level for 2 quarters.
conditional_forecast_paths;
  var i;
  periods 1 2;
  values  0.000988 0.000988;
end;
conditional_forecast(parameter_set=calibration, controlled_varexo=(e_i),
                     periods=8, replic=2000);
verbatim;
  export_conditional('cond_hold_2q');
end;

// Scenario 2: tighten and hold the policy-rate gap ~1 p.p. (annual) above its
// historical mean for 4 quarters.
conditional_forecast_paths;
  var i;
  periods 1 2 3 4;
  values  0.0025 0.0025 0.0025 0.0025;
end;
conditional_forecast(parameter_set=calibration, controlled_varexo=(e_i),
                     periods=8, replic=2000);
verbatim;
  export_conditional('cond_hike_4q');
end;

// Scenario 3: force inflation to the 3% official target for 4 quarters.
// The steady state corresponds to the sample-mean inflation (~3.79% annual), so
// the 3% target maps to a NEGATIVE inflation gap of (1.03)^(1/4)-1 minus the
// mean quarterly inflation = -0.0019233. Imposing this gap is a genuine
// disinflation experiment (not vacuous), enforced through the policy shock e_i.
conditional_forecast_paths;
  var pi;
  periods 1 2 3 4;
  values  -0.0019233 -0.0019233 -0.0019233 -0.0019233;
end;
conditional_forecast(parameter_set=calibration, controlled_varexo=(e_i),
                     periods=8, replic=2000);
verbatim;
  export_conditional('cond_pi_to_target_4q');
end;
\end{lstlisting}
\clearpage
\section{dynare/nk\_chile\_history.mod}
\begin{lstlisting}[language={}]
// Historical analysis of the calibrated Chilean New Keynesian model.
// Runs the Kalman smoother over 2001Q1-2026Q1 at the fixed calibration and
// asks the macro-narrative question: which structural shocks (demand,
// cost-push, monetary) actually drove inflation, the output gap and the policy
// rate over the sample? Uses Dynare's shock_decomposition, plus the smoothed
// structural shock series. Same deviation-form model as the forecast file.

var x pi i;
varexo e_x e_pi e_i;
parameters beta sigma kappa rho_i phi_pi phi_x;

beta   = 0.992637536145;
sigma  = 1.0;
kappa  = 0.10;
rho_i  = 0.934095195259;
phi_pi = 1.75;
phi_x  = 0.50;

model(linear);
  x  = x(+1) - (1/sigma)*(i - pi(+1)) + e_x;
  pi = beta*pi(+1) + kappa*x + e_pi;
  i  = rho_i*i(-1) + (1-rho_i)*(phi_pi*pi + phi_x*x) + e_i;
end;

steady;
check;

shocks;
  var e_x;  stderr 0.005;
  var e_pi; stderr 0.00277;
  var e_i;  stderr 0.00068;
end;

varobs pi i x;

estimated_params;
  stderr e_x,  inv_gamma_pdf, 0.005,   inf;
  stderr e_pi, inv_gamma_pdf, 0.00277, inf;
  stderr e_i,  inv_gamma_pdf, 0.00068, inf;
end;

estimated_params_init(use_calibration);
  stderr e_x,  0.005;
  stderr e_pi, 0.00277;
  stderr e_i,  0.00068;
end;

// Kalman smoother at the fixed calibration (no estimation).
estimation(datafile='chile_observables_dynare.csv', mode_compute=0, mh_replic=0,
           smoother, nograph, nodisplay) x pi i;

// Historical decomposition: contribution of each shock to each variable.
shock_decomposition(parameter_set=calibration, nograph) pi i x;

verbatim;
  export_history('history');
end;
\end{lstlisting}
\clearpage
\section{dynare/nk\_chile\_history\_hybrid.mod}
\begin{lstlisting}[language={}]
// Historical shock decomposition for the HYBRID-NKPC model (intrinsic inflation
// persistence). Same smoother + shock_decomposition exercise as nk_chile_history.mod,
// but with the indexation term gamma_pi*pi(-1). Used to ask whether the story of
// Chilean IPC inflation changes once inflation has memory.

var x pi i;
varexo e_x e_pi e_i;
parameters beta sigma kappa rho_i phi_pi phi_x gamma_pi;

beta     = 0.992637536145;
sigma    = 1.0;
kappa    = 0.10;
rho_i    = 0.934095195259;
phi_pi   = 1.75;
phi_x    = 0.50;
gamma_pi = 0.35;

model(linear);
  x  = x(+1) - (1/sigma)*(i - pi(+1)) + e_x;
  pi = gamma_pi*pi(-1) + (1-gamma_pi)*(beta*pi(+1) + kappa*x) + e_pi;
  i  = rho_i*i(-1) + (1-rho_i)*(phi_pi*pi + phi_x*x) + e_i;
end;

steady;
check;

shocks;
  var e_x;  stderr 0.005;
  var e_pi; stderr 0.00277;
  var e_i;  stderr 0.00068;
end;

varobs pi i x;

estimated_params;
  stderr e_x,  inv_gamma_pdf, 0.005,   inf;
  stderr e_pi, inv_gamma_pdf, 0.00277, inf;
  stderr e_i,  inv_gamma_pdf, 0.00068, inf;
end;

estimated_params_init(use_calibration);
  stderr e_x,  0.005;
  stderr e_pi, 0.00277;
  stderr e_i,  0.00068;
end;

estimation(datafile='chile_observables_dynare.csv', mode_compute=0, mh_replic=0,
           smoother, nograph, nodisplay) x pi i;

shock_decomposition(parameter_set=calibration, nograph) pi i x;

verbatim;
  export_history('history_hybrid');
end;
\end{lstlisting}
\clearpage
\section{dynare/nk\_chile\_mcmc.mod}
\begin{lstlisting}[language={}]
// Full Bayesian MCMC extension of the minimal Chilean NK model.
var x pi i;
varexo e_x e_pi e_i;
parameters beta sigma kappa rho_i phi_pi phi_x;

beta   = 0.992637536145;
sigma  = 1.0;
kappa  = 0.10;
rho_i  = 0.90;
phi_pi = 1.75;
phi_x  = 0.50;

model(linear);
  x  = x(+1) - (1/sigma)*(i-pi(+1)) + e_x;
  pi = beta*pi(+1) + kappa*x + e_pi;
  i  = rho_i*i(-1) + (1-rho_i)*(phi_pi*pi+phi_x*x) + e_i;
end;

steady;
check;
varobs pi i x;

estimated_params;
  sigma,       gamma_pdf,     1.00,  0.40;
  kappa,       gamma_pdf,     0.10,  0.05;
  rho_i,       beta_pdf,      0.80,  0.10;
  phi_pi,      gamma_pdf,     1.50,  0.25;
  phi_x,       gamma_pdf,     0.50,  0.25;
  stderr e_x,  inv_gamma_pdf, 0.005, inf;
  stderr e_pi, inv_gamma_pdf, 0.0025, inf;
  stderr e_i,  inv_gamma_pdf, 0.002, inf;
end;

// Two independent chains. Ten thousand retained proposals per chain are enough
// for a serious course extension, but convergence is assessed rather than assumed.
estimation(datafile='chile_observables_dynare.csv', mode_compute=4,
           mh_replic=10000, mh_nblocks=2, mh_drop=0.30, mh_jscale=0.30,
           mh_conf_sig=0.90, plot_priors=0, nograph, nodisplay);

verbatim;
  export_mcmc('mcmc');
end;
\end{lstlisting}
\clearpage
\section{dynare/nk\_chile\_open.mod}
\begin{lstlisting}[language={}]
// Illustrative small-open-economy extension for Chile.
var x pi i q copper;
varexo e_x e_pi e_i e_q e_copper;
parameters beta sigma kappa rho_i phi_pi phi_x alpha_q alpha_copper
           gamma_q eta_i phi_q rho_copper;

beta = 0.992637536145;
sigma = 1;
kappa = 0.1;
rho_i = 0.934095195259;
phi_pi = 1.75;
phi_x = 0.5;
alpha_q = 0.05;
alpha_copper = 0.03;
gamma_q = 0.03;
eta_i = 1;
phi_q = 0.05;
rho_copper = 0.941182126643;

model(linear);
  // Domestic demand: real rate, competitiveness and copper-income channel.
  x = x(+1) - (1/sigma)*(i-pi(+1)) + alpha_q*q + alpha_copper*copper + e_x;
  // Domestic inflation with direct exchange-rate pass-through.
  pi = beta*pi(+1) + kappa*x + gamma_q*q + e_pi;
  // Taylor rule may lean mildly against depreciation.
  i = rho_i*i(-1) + (1-rho_i)*(phi_pi*pi + phi_x*x + phi_q*q) + e_i;
  // UIP-style exchange-rate equation; q>0 means a weaker peso.
  q = q(+1) - eta_i*i + e_q;
  // Exogenous persistent copper-price gap.
  copper = rho_copper*copper(-1) + e_copper;
end;

steady;
check;
verbatim;
  export_stability('open_economy');
end;

shocks;
  var e_x; stderr 0.005;
  var e_pi; stderr 0.00277;
  var e_i; stderr 0.00068;
  var e_q; stderr 0.0460980810142;
  var e_copper; stderr 0.12;
end;

stoch_simul(order=1, irf=20, nograph);
verbatim;
  export_extension('open_economy');
end;
\end{lstlisting}
\clearpage
\part*{Solucao e calibracao}\addcontentsline{toc}{part}{Solucao e calibracao}

\section{python/hybrid\_solution.py}
\begin{lstlisting}[language=Python]
"""Exact first-order reduced form of the hybrid-NKPC Chilean model.

The hybrid model adds intrinsic inflation persistence:

    x  = E x(+1) - (1/sigma)(i - E pi(+1)) + e_x
    pi = gamma*pi(-1) + (1-gamma)*(beta*E pi(+1) + kappa*x) + e_pi
    i  = rho*i(-1) + (1-rho)(phi_pi*pi + phi_x*x) + e_i

Now there are TWO predetermined states, s_t = [i_{t-1}, pi_{t-1}], so the minimal
state-variable solution is  y_t = T s_{t-1} + R eps_t  with T a 3x2 matrix. We
solve T by undetermined coefficients (a 6-equation fixed point) and validate the
solution against Dynare's hybrid IRFs (it reproduces them to ~1e-6).
"""

from __future__ import annotations

import numpy as np

from common import calibration_rho, complete_parameters

E1 = np.array([1.0, 0.0])  # picks i_{t-1}
E2 = np.array([0.0, 1.0])  # picks pi_{t-1}
I2 = np.eye(2)


def hybrid_params(gamma_pi: float = 0.35) -> dict:
    params = complete_parameters({"rho_i": calibration_rho()})
    params["gamma_pi"] = float(gamma_pi)
    return params


def solve_hybrid(params: dict, max_iter: int = 5000, tol: float = 1e-14) -> dict:
    """Unique stable MSV solution by contraction on the state transition A.

    Given a guess for the state transition A (s_t = A s_{t-1}, s=[i,pi]), the
    structural equations are LINEAR in T=[Tx;Tpi;Ti] (because the expectations
    become E_t y_{t+1}=T A s_{t-1}). We solve that linear system for T, update
    A=[Ti;Tpi], and iterate. For a determinate model this converges to the unique
    forward-stable solution (avoiding the spurious roots a blind root-finder hits).
    """
    beta, sigma, kappa = params["beta"], params["sigma"], params["kappa"]
    rho, phi_pi, phi_x = params["rho_i"], params["phi_pi"], params["phi_x"]
    g = params["gamma_pi"]

    def residual(v: np.ndarray, A: np.ndarray) -> np.ndarray:
        Tx, Tpi, Ti = v[0:2], v[2:4], v[4:6]
        r_is = Tx @ (I2 - A) + (1.0 / sigma) * Ti - (1.0 / sigma) * (Tpi @ A)
        r_pc = Tpi - (1.0 - g) * beta * (Tpi @ A) - (1.0 - g) * kappa * Tx - g * E2
        r_tr = Ti - (1.0 - rho) * phi_pi * Tpi - (1.0 - rho) * phi_x * Tx - rho * E1
        return np.concatenate([r_is, r_pc, r_tr])

    def solve_T(A: np.ndarray) -> np.ndarray:
        # residual is affine in v: L v + c = 0  ->  v = solve(L, -c)
        c = residual(np.zeros(6), A)
        L = np.zeros((6, 6))
        for k in range(6):
            ek = np.zeros(6)
            ek[k] = 1.0
            L[:, k] = residual(ek, A) - c
        return np.linalg.solve(L, -c)

    A = np.zeros((2, 2))
    v = np.zeros(6)
    for _ in range(max_iter):
        v = solve_T(A)
        A_new = np.vstack([v[4:6], v[2:4]])  # [Ti; Tpi]
        if np.max(np.abs(A_new - A)) < tol:
            A = A_new
            break
        A = A_new

    Tx, Tpi, Ti = v[0:2], v[2:4], v[4:6]
    S = np.vstack([Ti, Tpi])
    T = np.vstack([Tx, Tpi, Ti])  # rows x,pi,i; cols [i_{-1}, pi_{-1}]

    M = np.array([
        [1.0, -Tx[1] - Tpi[1] / sigma, -Tx[0] + 1.0 / sigma - Tpi[0] / sigma],
        [-(1.0 - g) * kappa, 1.0 - (1.0 - g) * beta * Tpi[1], -(1.0 - g) * beta * Tpi[0]],
        [-(1.0 - rho) * phi_x, -(1.0 - rho) * phi_pi, 1.0],
    ])
    R = np.linalg.inv(M)
    return {"T": T, "S": S, "R": R, "eig": np.linalg.eigvals(S),
            "stable": bool(np.max(np.abs(np.linalg.eigvals(S))) < 1.0)}


def hybrid_irf(sol: dict, shock_std: dict, horizon: int = 20) -> dict:
    T, S, R = sol["T"], sol["S"], sol["R"]
    out = {}
    for j, shock in enumerate(("e_x", "e_pi", "e_i")):
        e = np.zeros(3)
        e[j] = shock_std[shock]
        y0 = R @ e
        responses = [y0]
        s = np.array([y0[2], y0[1]])  # [i_0, pi_0]
        for _ in range(1, horizon):
            yk = T @ s
            responses.append(yk)
            s = np.array([yk[2], yk[1]])
        out[shock] = np.array(responses)
    return out


def hybrid_forecast(sol: dict, last_i_dev: float, last_pi_dev: float,
                    horizon: int) -> np.ndarray:
    """h-step forecast of [x, pi, i] deviations from the origin state."""
    s = np.array([last_i_dev, last_pi_dev])
    Sp = np.linalg.matrix_power(sol["S"], horizon - 1)
    return sol["T"] @ Sp @ s


def _validate() -> None:
    import pandas as pd

    from common import DYNARE_OUTPUTS

    params = hybrid_params(0.35)
    sol = solve_hybrid(params)
    print(f"Stable: {sol['stable']}; |eig(S)| = {np.abs(sol['eig']).round(4)}")

    shock_std = {"e_x": 0.005, "e_pi": 0.00277, "e_i": 0.00068}
    mine = hybrid_irf(sol, shock_std, horizon=20)

    dyn = pd.read_csv(DYNARE_OUTPUTS / "hybrid_nkpc" / "irfs.csv")
    var_index = {"x": 0, "pi": 1, "i": 2}
    max_diff = 0.0
    for shock in ("e_x", "e_pi", "e_i"):
        for var, vi in var_index.items():
            col = f"{var}_{shock}"
            if col not in dyn:
                continue
            diff = np.max(np.abs(mine[shock][:, vi] - dyn[col].to_numpy()[:20]))
            max_diff = max(max_diff, diff)
    print(f"Max |IRF difference| vs Dynare hybrid: {max_diff:.3e}")
    print("VALIDATION OK" if max_diff < 1e-5 else "VALIDATION FAILED")


if __name__ == "__main__":
    _validate()
\end{lstlisting}
\clearpage
\section{python/calibrate\_shocks.py}
\begin{lstlisting}[language=Python]
"""Calibrate shock standard deviations to two transparent FEVD targets.

The assignment asks for shock sizes that generate plausible impulse responses and
moments. We make that choice transparent by specifying an illustrative variance-share
target: inflation is mostly driven by cost-push shocks, the policy rate mostly by its
own innovation, and the output gap is split between demand and monetary shocks.

FEVD shares depend only on the products sigma_k^2 * V_k(v), where V_k(v) is the asymptotic
variance of variable v generated by a unit-variance shock k. We compute V_k(v) from the
exact linear-RE solution and then search the shock sizes (anchoring sigma_ex = 0.005) that
best match the target shares at the estimated Chilean calibration.

The first target is the didactic benchmark stated in the assignment material. The
second is the researcher's alternative target originally used in this repository.
Neither is a historical decomposition of Chile or an official BCCh estimate.
"""

from __future__ import annotations

import numpy as np
import pandas as pd
import matplotlib

matplotlib.use("Agg")
import matplotlib.pyplot as plt
from scipy.optimize import least_squares

from common import (
    FIGURES,
    INITIAL_SHOCK_STD,
    TABLES,
    _state_coefficients,
    calibration_rho,
    complete_parameters,
    ensure_directories,
)

VARS = ("x", "pi", "i")
SHOCKS = ("e_x", "e_pi", "e_i")

TARGETS = {
    "didactic_benchmark": {
        "x": {"e_x": 0.375, "e_pi": 0.100, "e_i": 0.524},
        "pi": {"e_x": 0.014, "e_pi": 0.823, "e_i": 0.162},
        "i": {"e_x": 0.087, "e_pi": 0.147, "e_i": 0.766},
    },
    "researcher_alternative": {
        "x": {"e_x": 0.45, "e_pi": 0.10, "e_i": 0.45},
        "pi": {"e_x": 0.02, "e_pi": 0.80, "e_i": 0.18},
        "i": {"e_x": 0.10, "e_pi": 0.15, "e_i": 0.75},
    },
}


def per_shock_variance(params: dict, horizon: int = 4000) -> np.ndarray:
    """V[v, k] = asymptotic variance of variable v from a unit-variance shock k."""
    a_x, a_pi, a_i = _state_coefficients(params)
    sigma, beta, kappa = params["sigma"], params["beta"], params["kappa"]
    rho_i, phi_pi, phi_x = params["rho_i"], params["phi_pi"], params["phi_x"]
    c_i = a_x - (1.0 - a_pi) / sigma
    impact_matrix = np.array(
        [
            [1.0, 0.0, -c_i],
            [-kappa, 1.0, -beta * a_pi],
            [-(1.0 - rho_i) * phi_x, -(1.0 - rho_i) * phi_pi, 1.0],
        ]
    )
    propagation = np.array([a_x, a_pi, a_i])
    variance = np.zeros((3, 3))
    for k in range(3):
        innovation = np.zeros(3)
        innovation[k] = 1.0
        response = np.linalg.solve(impact_matrix, innovation)
        accum = response ** 2
        lag = response[2]
        for _ in range(1, horizon):
            response = propagation * lag
            accum = accum + response ** 2
            lag = response[2]
        variance[:, k] = accum
    return variance


def fevd_from_sigmas(variance: np.ndarray, sigmas: np.ndarray) -> np.ndarray:
    contributions = variance * (sigmas ** 2)  # column k scaled by sigma_k^2
    return contributions / contributions.sum(axis=1, keepdims=True)


def calibrate_to_target(
    variance: np.ndarray,
    initial_sigmas: np.ndarray,
    target: np.ndarray,
) -> tuple[np.ndarray, np.ndarray, float]:
    """Fit e_pi and e_i standard deviations while anchoring sigma_e_x."""

    target = target / target.sum(axis=1, keepdims=True)

    def residuals(free_log: np.ndarray) -> np.ndarray:
        sigmas = np.array(
            [initial_sigmas[0], np.exp(free_log[0]), np.exp(free_log[1])]
        )
        return (fevd_from_sigmas(variance, sigmas) - target).ravel()

    start = np.log([initial_sigmas[1], initial_sigmas[2]])
    solution = least_squares(residuals, start, method="lm")
    calibrated = np.array(
        [initial_sigmas[0], np.exp(solution.x[0]), np.exp(solution.x[1])]
    )
    achieved = fevd_from_sigmas(variance, calibrated)
    rmse_pp = 100.0 * float(np.sqrt(np.mean((achieved - target) ** 2)))
    return calibrated, achieved, rmse_pp


def plot_comparison(table: pd.DataFrame) -> None:
    fig, axes = plt.subplots(2, 3, figsize=(13, 7), sharey=True)
    labels = {"e_x": "Demand", "e_pi": "Cost-push", "e_i": "Monetary"}
    for row_index, calibration in enumerate(TARGETS):
        subset = table[table["calibration"] == calibration]
        for column_index, variable in enumerate(VARS):
            ax = axes[row_index, column_index]
            frame = subset[subset["variable"] == variable].set_index("shock")
            shocks = list(SHOCKS)
            positions = np.arange(len(shocks))
            ax.bar(
                positions - 0.18,
                frame.loc[shocks, "target_pct"],
                width=0.36,
                label="Target",
                color="#9CA3AF",
            )
            ax.bar(
                positions + 0.18,
                frame.loc[shocks, "achieved_pct"],
                width=0.36,
                label="Achieved",
                color="#1E4E79",
            )
            ax.set_xticks(positions, [labels[shock] for shock in shocks], rotation=20)
            ax.set_title(f"{calibration.replace('_', ' ').title()} - {variable}")
            ax.grid(axis="y", alpha=0.25)
            if column_index == 0:
                ax.set_ylabel("Variance share (%)")
            if row_index == 0 and column_index == 2:
                ax.legend(fontsize=8)
    fig.suptitle("FEVD calibration: stated targets versus model-achieved shares")
    fig.tight_layout(rect=(0, 0, 1, 0.96))
    fig.savefig(FIGURES / "fevd_calibration_comparison.png", dpi=180)
    plt.close(fig)


def main() -> None:
    ensure_directories()
    params = complete_parameters({"rho_i": calibration_rho()})
    variance = per_shock_variance(params)

    initial_sigmas = np.array(
        [INITIAL_SHOCK_STD["e_x"], INITIAL_SHOCK_STD["e_pi"], INITIAL_SHOCK_STD["e_i"]]
    )
    fevd_initial = fevd_from_sigmas(variance, initial_sigmas)

    comparison_rows = []
    sigma_rows = []
    fitted: dict[str, np.ndarray] = {}
    for calibration, target_dict in TARGETS.items():
        raw_target = np.array(
            [[target_dict[variable][shock] for shock in SHOCKS] for variable in VARS]
        )
        target = raw_target / raw_target.sum(axis=1, keepdims=True)
        calibrated, achieved, rmse_pp = calibrate_to_target(
            variance, initial_sigmas, target
        )
        fitted[calibration] = calibrated
        for variable_index, variable in enumerate(VARS):
            for shock_index, shock in enumerate(SHOCKS):
                comparison_rows.append(
                    {
                        "calibration": calibration,
                        "variable": variable,
                        "shock": shock,
                        "initial_pct": 100.0 * fevd_initial[variable_index, shock_index],
                        "target_pct": 100.0 * target[variable_index, shock_index],
                        "achieved_pct": 100.0 * achieved[variable_index, shock_index],
                        "difference_pp": 100.0
                        * (
                            achieved[variable_index, shock_index]
                            - target[variable_index, shock_index]
                        ),
                        "rmse_all_shares_pp": rmse_pp,
                        "rho_i": params["rho_i"],
                    }
                )
        for shock, sigma_initial, sigma_calibrated in zip(
            SHOCKS, initial_sigmas, calibrated
        ):
            sigma_rows.append(
                {
                    "calibration": calibration,
                    "shock": shock,
                    "sigma_initial": sigma_initial,
                    "sigma_calibrated": sigma_calibrated,
                    "rho_i": params["rho_i"],
                }
            )

    comparison = pd.DataFrame(comparison_rows)
    sigmas_comparison = pd.DataFrame(sigma_rows)
    comparison.to_csv(TABLES / "fevd_calibration_comparison.csv", index=False)
    sigmas_comparison.to_csv(TABLES / "shock_sigmas_comparison.csv", index=False)
    plot_comparison(comparison)

    # Preserve the original headline outputs for backwards compatibility. The
    # baseline remains the researcher alternative; the benchmark is a robustness case.
    baseline_name = "researcher_alternative"
    baseline = comparison[comparison["calibration"] == baseline_name].copy()
    baseline = baseline.rename(
        columns={
            "initial_pct": "fevd_initial_sigmas_pct",
            "target_pct": "illustrative_target_pct",
            "achieved_pct": "fevd_calibrated_sigmas_pct",
        }
    )
    baseline["calibration_type"] = baseline_name
    baseline[
        [
            "variable",
            "shock",
            "fevd_initial_sigmas_pct",
            "illustrative_target_pct",
            "fevd_calibrated_sigmas_pct",
            "calibration_type",
        ]
    ].to_csv(TABLES / "shock_calibration.csv", index=False)
    pd.DataFrame(
        {
            "shock": SHOCKS,
            "sigma_initial": initial_sigmas,
            "sigma_calibrated": fitted[baseline_name],
            "calibration_type": baseline_name,
        }
    ).to_csv(TABLES / "shock_sigmas.csv", index=False)

    print(f"rho_i (Chile baseline) = {params['rho_i']:.4f}")
    print(comparison.round(3).to_string(index=False))
    print("\nShock-standard-deviation comparison:")
    print(sigmas_comparison.round(6).to_string(index=False))
    print("Wrote FEVD comparison tables and figure.")


if __name__ == "__main__":
    main()
\end{lstlisting}
\clearpage
\section{python/analyze\_determinacy.py}
\begin{lstlisting}[language=Python]
"""Map Blanchard-Kahn determinacy and convergence speed against phi_pi.

The assignment's phi_pi grid (1.3-2.2) is entirely determinate, so a verbal "phi_pi>1"
rule is not enough. Here we sweep phi_pi (including values below 1) and, for each, count
the eigenvalues inside the unit circle from the model's characteristic cubic. With one
predetermined variable the model is determinate iff exactly one root is stable. We also
record the largest stable-root modulus, which governs how fast the economy converges
after a shock. This complements (and is cross-checked by) Dynare's per-scenario check.

Outputs: outputs/tables/determinacy_map.csv and outputs/figures/determinacy_map.png.
"""

from __future__ import annotations

import matplotlib

matplotlib.use("Agg")
import matplotlib.pyplot as plt
import numpy as np
import pandas as pd

from common import (
    FIGURES,
    TABLES,
    calibration_rho,
    characteristic_roots,
    complete_parameters,
    ensure_directories,
)


def main() -> None:
    ensure_directories()
    rho_i = calibration_rho()
    grid = np.round(np.arange(0.5, 3.0001, 0.05), 2)
    rows = []
    for phi_pi in grid:
        params = complete_parameters({"rho_i": rho_i, "phi_pi": float(phi_pi)})
        roots = characteristic_roots(params)
        moduli = np.abs(roots)
        n_stable = int(np.sum(moduli < 1.0 - 1e-9))
        stable_moduli = moduli[moduli < 1.0 - 1e-9]
        dominant = float(stable_moduli.max()) if stable_moduli.size else np.nan
        rows.append(
            {
                "phi_pi": float(phi_pi),
                "n_stable_roots": n_stable,
                "n_unstable_roots": int(np.sum(moduli > 1.0 + 1e-9)),
                "determinate": n_stable == 1,
                "dominant_stable_modulus": dominant,
            }
        )
    table = pd.DataFrame(rows)
    table.to_csv(TABLES / "determinacy_map.csv", index=False)

    threshold = table.loc[table["determinate"], "phi_pi"].min()

    fig, ax1 = plt.subplots(figsize=(8.8, 5.0))
    colors = ["#16a34a" if d else "#b91c1c" for d in table["determinate"]]
    ax1.bar(table["phi_pi"], table["n_stable_roots"], width=0.045, color=colors, alpha=0.65)
    ax1.axhline(1.0, color="black", lw=0.8, ls=":")
    ax1.set_xlabel("phi_pi (Taylor-rule inflation response, phi_x = 0.5)")
    ax1.set_ylabel("Number of stable eigenvalues")
    ax1.set_ylim(0, 3.2)
    ax2 = ax1.twinx()
    ax2.plot(table["phi_pi"], table["dominant_stable_modulus"], color="#1E4E79",
             lw=2, marker="o", ms=3, label="Dominant stable |eigenvalue|")
    ax2.set_ylabel("Dominant stable |eigenvalue| (persistence)")
    ax1.axvline(1.0, color="gray", lw=1, ls="--")
    ax1.set_title(
        f"Determinacy and convergence vs phi_pi (rho_i={rho_i:.2f}); "
        f"determinate for phi_pi >= {threshold:.2f}"
    )
    ax2.legend(loc="upper right", fontsize=8)
    fig.text(0.01, 0.01, "Green = determinate (1 stable root); red = indeterminate/explosive. "
             "Source: characteristic roots of the calibrated model (cross-checked with Dynare).",
             fontsize=7)
    fig.tight_layout(rect=(0, 0.03, 1, 1))
    fig.savefig(FIGURES / "determinacy_map.png", dpi=180)
    plt.close(fig)

    print(table.to_string(index=False))
    print(f"Determinate for phi_pi >= {threshold:.2f}")
    print(f"Wrote {TABLES / 'determinacy_map.csv'} and {FIGURES / 'determinacy_map.png'}")


if __name__ == "__main__":
    main()
\end{lstlisting}
\clearpage
\part*{Estimacao dos parametros}\addcontentsline{toc}{part}{Estimacao dos parametros}

\section{python/estimate\_nkpc.py}
\begin{lstlisting}[language=Python]
"""Estimate the New Keynesian Phillips curve slope kappa (Chile, real data).

Inflation pi_t and the output gap x_t are quarterly fractions from the BCCh panel.
Three estimators of the slope kappa are reported:

  * Backward / accelerationist OLS (HAC):  pi_t = c + gamma_b*pi_{t-1} + kappa*x_t.
  * Forward-looking IV (2SLS):              pi_t = c + gamma_f*pi_{t+1} + kappa*x_t,
        with E_t pi_{t+1} instrumented by lagged inflation, gap and rate.
  * Restricted (gamma_f = beta = 0.9926):   (pi_t - beta*pi_{t+1}) = c + kappa*x_t.

Output: outputs/tables/nkpc_estimates.csv. The estimates are compared with the
calibration grid kappa in {0.07, 0.10, 0.13}.
"""

from __future__ import annotations

import json
import math

import numpy as np
import pandas as pd

from common import DATA_CLEAN, TABLES, complete_parameters, ensure_directories

KAPPA_GRID = (0.07, 0.10, 0.13)


def _tsls_hac(y, X, Z, maxlags: int = 4):
    ZtZ_inv = np.linalg.inv(Z.T @ Z)
    Pz = Z @ ZtZ_inv @ Z.T
    XtPzX_inv = np.linalg.inv(X.T @ Pz @ X)
    beta = XtPzX_inv @ X.T @ Pz @ y
    resid = y - X @ beta
    fitted_x = Pz @ X
    scores = fitted_x * resid[:, None]
    meat = scores.T @ scores
    for lag in range(1, min(maxlags, len(scores) - 1) + 1):
        weight = 1.0 - lag / (maxlags + 1.0)
        gamma = scores[lag:].T @ scores[:-lag]
        meat += weight * (gamma + gamma.T)
    cov = XtPzX_inv @ meat @ XtPzX_inv
    r2 = 1.0 - float(resid @ resid) / float(((y - y.mean()) ** 2).sum())
    return beta, np.sqrt(np.diag(cov)), r2


def _first_stage_f(endogenous, instruments) -> float:
    import statsmodels.api as sm

    result = sm.OLS(endogenous, instruments).fit()
    restrictions = np.zeros((instruments.shape[1] - 1, instruments.shape[1]))
    restrictions[:, 1:] = np.eye(instruments.shape[1] - 1)
    return float(result.f_test(restrictions).fvalue)


def main() -> None:
    ensure_directories()
    panel = pd.read_csv(DATA_CLEAN / "chile_macro_quarterly.csv")
    metadata = json.loads(
        (DATA_CLEAN / "dataset_metadata.json").read_text(encoding="utf-8")
    )
    is_synthetic = bool(metadata.get("is_synthetic", True))
    data_source = str(metadata.get("institution", metadata.get("source", "unknown")))
    base = pd.DataFrame({"pi": panel["infl_q"], "x": panel["output_gap"]})
    base["pi_lead"] = base["pi"].shift(-1)
    base["pi_l1"] = base["pi"].shift(1)
    base["pi_l2"] = base["pi"].shift(2)
    base["x_l1"] = base["x"].shift(1)
    base["x_l2"] = base["x"].shift(2)

    beta_disc = complete_parameters()["beta"]
    rows = []

    import statsmodels.api as sm

    # Backward-looking OLS (HAC).
    bw = base.dropna(subset=["pi", "pi_l1", "x"]).reset_index(drop=True)
    Xb = sm.add_constant(np.column_stack([bw["pi_l1"], bw["x"]]))
    ols = sm.OLS(bw["pi"].to_numpy(), Xb).fit(cov_type="HAC", cov_kwds={"maxlags": 4})
    rows.append(
        {
            "method": "backward_OLS_HAC",
            "gamma": float(ols.params[1]),
            "kappa": float(ols.params[2]),
            "se_kappa": float(ols.bse[2]),
            "r_squared": float(ols.rsquared),
            "observations": int(len(bw)),
            "first_stage_F_pi_lead": np.nan,
            "first_stage_F_x": np.nan,
            "is_synthetic": is_synthetic,
            "data_source": data_source,
        }
    )

    # Forward-looking IV (instrument pi_{t+1} with lags).
    fw = base.dropna().reset_index(drop=True)
    n = len(fw)
    y = fw["pi"].to_numpy()
    X = np.column_stack([np.ones(n), fw["pi_lead"], fw["x"]])
    Z = np.column_stack(
        [np.ones(n), fw["pi_l1"], fw["pi_l2"], fw["x_l1"], fw["x_l2"]]
    )
    beta_iv, se_iv, r2_iv = _tsls_hac(y, X, Z)
    rows.append(
        {
            "method": "forward_IV_2SLS_HAC",
            "gamma": float(beta_iv[1]),
            "kappa": float(beta_iv[2]),
            "se_kappa": float(se_iv[2]),
            "r_squared": r2_iv,
            "observations": n,
            "first_stage_F_pi_lead": _first_stage_f(fw["pi_lead"].to_numpy(), Z),
            "first_stage_F_x": _first_stage_f(fw["x"].to_numpy(), Z),
            "is_synthetic": is_synthetic,
            "data_source": data_source,
        }
    )

    # Restricted: impose gamma_f = beta.
    rs = base.dropna(subset=["pi", "pi_lead", "x"]).reset_index(drop=True)
    y_r = (rs["pi"] - beta_disc * rs["pi_lead"]).to_numpy()
    Xr = sm.add_constant(rs["x"].to_numpy())
    olsr = sm.OLS(y_r, Xr).fit(cov_type="HAC", cov_kwds={"maxlags": 4})
    rows.append(
        {
            "method": f"restricted_gamma_eq_beta({beta_disc:.4f})",
            "gamma": beta_disc,
            "kappa": float(olsr.params[1]),
            "se_kappa": float(olsr.bse[1]),
            "r_squared": float(olsr.rsquared),
            "observations": int(len(rs)),
            "first_stage_F_pi_lead": np.nan,
            "first_stage_F_x": np.nan,
            "is_synthetic": is_synthetic,
            "data_source": data_source,
        }
    )

    table = pd.DataFrame(rows)
    table["kappa_grid_nearest"] = table["kappa"].apply(
        lambda k: min(KAPPA_GRID, key=lambda g: abs(g - k))
    )
    table.to_csv(TABLES / "nkpc_estimates.csv", index=False)
    print(table.round(4).to_string(index=False))
    print(f"Wrote {TABLES / 'nkpc_estimates.csv'}")


if __name__ == "__main__":
    main()
\end{lstlisting}
\clearpage
\section{python/estimate\_taylor\_rule.py}
\begin{lstlisting}[language=Python]
"""Estimate a reduced-form Taylor rule with interest-rate smoothing (Chile).

Specification (all variables in quarterly fractions, from the real BCCh panel):

    i_t = c + a * i_{t-1} + b * pi_t + d * x_t + eps_t,

with the structural mapping rho_i = a, phi_pi = b / (1 - a), phi_x = d / (1 - a).

Two estimators are reported:
  * OLS with Newey-West (HAC) standard errors.
  * 2SLS/IV instrumenting current inflation and the gap with their own lags (and the
    lagged rate), which addresses the simultaneity between the policy rate and the
    contemporaneous inflation/gap it reacts to.

Structural-parameter standard errors use the delta method. Output:
outputs/tables/taylor_rule_estimates.csv.
"""

from __future__ import annotations

import json
import math

import numpy as np
import pandas as pd

from common import DATA_CLEAN, TABLES, ensure_directories


def _delta_ratio(params: np.ndarray, cov: np.ndarray, i_a: int, i_num: int) -> tuple[float, float]:
    """phi = params[i_num] / (1 - params[i_a]); delta-method standard error."""
    a = params[i_a]
    num = params[i_num]
    phi = num / (1.0 - a)
    grad = np.zeros(len(params))
    grad[i_a] = num / (1.0 - a) ** 2
    grad[i_num] = 1.0 / (1.0 - a)
    var = float(grad @ cov @ grad)
    return phi, math.sqrt(var) if var > 0 else math.nan


def _tsls_hac(y: np.ndarray, X: np.ndarray, Z: np.ndarray, maxlags: int = 4):
    ZtZ_inv = np.linalg.inv(Z.T @ Z)
    Pz = Z @ ZtZ_inv @ Z.T
    XtPzX_inv = np.linalg.inv(X.T @ Pz @ X)
    beta = XtPzX_inv @ X.T @ Pz @ y
    resid = y - X @ beta
    fitted_x = Pz @ X
    scores = fitted_x * resid[:, None]
    meat = scores.T @ scores
    for lag in range(1, min(maxlags, len(scores) - 1) + 1):
        weight = 1.0 - lag / (maxlags + 1.0)
        gamma = scores[lag:].T @ scores[:-lag]
        meat += weight * (gamma + gamma.T)
    cov = XtPzX_inv @ meat @ XtPzX_inv
    r2 = 1.0 - float(resid @ resid) / float(((y - y.mean()) ** 2).sum())
    return beta, cov, r2


def _first_stage_f(endogenous: np.ndarray, instruments: np.ndarray) -> float:
    import statsmodels.api as sm

    result = sm.OLS(endogenous, instruments).fit()
    restrictions = np.zeros((instruments.shape[1] - 2, instruments.shape[1]))
    restrictions[:, 2:] = np.eye(instruments.shape[1] - 2)
    return float(result.f_test(restrictions).fvalue)


def main() -> None:
    ensure_directories()
    panel = pd.read_csv(DATA_CLEAN / "chile_macro_quarterly.csv")
    metadata = json.loads(
        (DATA_CLEAN / "dataset_metadata.json").read_text(encoding="utf-8")
    )
    is_synthetic = bool(metadata.get("is_synthetic", True))
    data_source = str(metadata.get("institution", metadata.get("source", "unknown")))
    base = pd.DataFrame(
        {
            "i": panel["i_q"],
            "pi": panel["infl_q"],
            "x": panel["output_gap"],
        }
    )
    for col in ("i", "pi", "x"):
        base[f"{col}_l1"] = base[col].shift(1)
        base[f"{col}_l2"] = base[col].shift(2)
    reg = base.dropna().reset_index(drop=True)
    n = len(reg)

    y = reg["i"].to_numpy()
    X = np.column_stack([np.ones(n), reg["i_l1"], reg["pi"], reg["x"]])  # c, a, b, d
    Z = np.column_stack(
        [np.ones(n), reg["i_l1"], reg["pi_l1"], reg["pi_l2"], reg["x_l1"], reg["x_l2"]]
    )

    rows = []

    import statsmodels.api as sm

    ols = sm.OLS(y, X).fit(cov_type="HAC", cov_kwds={"maxlags": 4})
    cov_ols = np.asarray(ols.cov_params())
    phi_pi, se_pp = _delta_ratio(ols.params, cov_ols, 1, 2)
    phi_x, se_px = _delta_ratio(ols.params, cov_ols, 1, 3)
    rows.append(
        {
            "method": "OLS_HAC",
            "c": float(ols.params[0]),
            "rho_i": float(ols.params[1]),
            "se_rho_i": float(ols.bse[1]),
            "phi_pi": phi_pi,
            "se_phi_pi": se_pp,
            "phi_x": phi_x,
            "se_phi_x": se_px,
            "r_squared": float(ols.rsquared),
            "observations": n,
            "first_stage_F_pi": np.nan,
            "first_stage_F_x": np.nan,
            "is_synthetic": is_synthetic,
            "data_source": data_source,
        }
    )

    beta_iv, cov_iv, r2_iv = _tsls_hac(y, X, Z)
    se_iv = np.sqrt(np.diag(cov_iv))
    phi_pi_iv, se_pp_iv = _delta_ratio(beta_iv, cov_iv, 1, 2)
    phi_x_iv, se_px_iv = _delta_ratio(beta_iv, cov_iv, 1, 3)
    rows.append(
        {
            "method": "IV_2SLS_HAC",
            "c": float(beta_iv[0]),
            "rho_i": float(beta_iv[1]),
            "se_rho_i": float(se_iv[1]),
            "phi_pi": phi_pi_iv,
            "se_phi_pi": se_pp_iv,
            "phi_x": phi_x_iv,
            "se_phi_x": se_px_iv,
            "r_squared": r2_iv,
            "observations": n,
            "first_stage_F_pi": _first_stage_f(reg["pi"].to_numpy(), Z),
            "first_stage_F_x": _first_stage_f(reg["x"].to_numpy(), Z),
            "is_synthetic": is_synthetic,
            "data_source": data_source,
        }
    )

    table = pd.DataFrame(rows)
    table.to_csv(TABLES / "taylor_rule_estimates.csv", index=False)
    print(table.round(4).to_string(index=False))
    print(f"Wrote {TABLES / 'taylor_rule_estimates.csv'}")


if __name__ == "__main__":
    main()
\end{lstlisting}
\clearpage
\section{python/estimate\_rhoi\_chile.py}
\begin{lstlisting}[language=Python]
"""Estimate the AR(1) persistence of the quarterly policy rate and pick rho_i.

Implements exactly the assignment's step: regress the quarterly policy rate on its
own lag, i_t = alpha + rho_i * i_{t-1} + u_t, by OLS with Newey-West (HAC) standard
errors. With real BCCh data the estimate is used directly in the calibration; with a
synthetic series (or an implausible estimate) it falls back to the baseline 0.80.
"""

from __future__ import annotations

import math

import numpy as np
import pandas as pd

from common import BASELINE, DATA_CLEAN, TABLES, ensure_directories


def main() -> None:
    ensure_directories()
    input_path = DATA_CLEAN / "chile_policy_rate.csv"
    if not input_path.exists():
        raise FileNotFoundError(
            "Missing data/clean/chile_policy_rate.csv. "
            "Run python/build_chile_dataset.py first."
        )

    frame = pd.read_csv(input_path)
    required = {"policy_rate", "is_synthetic"}
    missing = required.difference(frame.columns)
    if missing:
        raise ValueError(f"Clean dataset is missing columns: {sorted(missing)}")

    rate = pd.to_numeric(frame["policy_rate"], errors="coerce")
    valid = rate.dropna()
    if len(valid) < 11:
        raise ValueError("At least 11 valid policy-rate observations are required.")

    dates = frame["date"].astype(str) if "date" in frame.columns else None
    sample_start = dates.iloc[1] if dates is not None else ""
    sample_end = dates.iloc[-1] if dates is not None else ""

    current = valid.to_numpy()[1:]
    lagged = valid.to_numpy()[:-1]
    n_obs = len(current)

    se_hac = math.nan
    method = "statsmodels_ols_hac"
    try:
        import statsmodels.api as sm

        design = sm.add_constant(lagged)
        ols = sm.OLS(current, design).fit()
        hac = sm.OLS(current, design).fit(cov_type="HAC", cov_kwds={"maxlags": 4})
        alpha = float(ols.params[0])
        rho_estimate = float(ols.params[1])
        se_ols = float(ols.bse[1])
        se_hac = float(hac.bse[1])
        r_squared = float(ols.rsquared)
    except ImportError:
        method = "numpy_lstsq_fallback"
        design = np.column_stack([np.ones(n_obs), lagged])
        coefficients, _, _, _ = np.linalg.lstsq(design, current, rcond=None)
        alpha, rho_estimate = map(float, coefficients)
        residual = current - design @ coefficients
        sigma2 = float(residual @ residual / (n_obs - 2))
        covariance = sigma2 * np.linalg.inv(design.T @ design)
        se_ols = float(np.sqrt(covariance[1, 1]))
        r_squared = 1.0 - float(residual @ residual) / float(
            ((current - current.mean()) ** 2).sum()
        )

    se_for_ci = se_hac if math.isfinite(se_hac) else se_ols
    ci_low = rho_estimate - 1.96 * se_for_ci
    ci_high = rho_estimate + 1.96 * se_for_ci
    half_life = (
        math.log(0.5) / math.log(rho_estimate)
        if 0.0 < rho_estimate < 1.0
        else math.nan
    )
    long_run_mean = alpha / (1.0 - rho_estimate) if rho_estimate != 1.0 else math.nan

    is_synthetic = bool(
        frame["is_synthetic"].astype(str).str.lower().isin(["true", "1"]).all()
    )
    data_source = (
        str(frame["data_source"].iloc[0])
        if "data_source" in frame.columns
        else "unknown"
    )
    if is_synthetic or not 0.0 <= rho_estimate < 0.999:
        rho_used = BASELINE["rho_i"]
        basis = "fallback_0.80_synthetic_or_implausible_estimate"
    else:
        rho_used = min(max(rho_estimate, 0.0), 0.98)
        basis = "estimated_ar1_real_bcch_tpm"

    output = pd.DataFrame(
        [
            {
                "alpha": alpha,
                "rho_i_estimate": rho_estimate,
                "rho_i_se_ols": se_ols,
                "rho_i_se_hac": se_hac,
                "rho_i_ci95_low": ci_low,
                "rho_i_ci95_high": ci_high,
                "half_life_quarters": half_life,
                "long_run_mean_pct": long_run_mean,
                "rho_i_used": rho_used,
                "r_squared": r_squared,
                "observations": n_obs,
                "sample_start": sample_start,
                "sample_end": sample_end,
                "is_synthetic": is_synthetic,
                "data_source": data_source,
                "calibration_basis": basis,
                "estimation_method": method,
            }
        ]
    )
    destination = TABLES / "rhoi_estimate.csv"
    output.to_csv(destination, index=False)
    print(output.T.to_string())
    print(f"Wrote {destination}")


if __name__ == "__main__":
    main()
\end{lstlisting}
\clearpage
\section{python/estimate\_rstar.py}
\begin{lstlisting}[language=Python]
"""Light empirical anchor for the neutral real rate r* (Chile, real data).

The minimal NK model treats r* as a calibrated constant; the assignment asks for the
{2,3,4}% annual scenarios. Here we provide a data-based reference point so the scenario
grid can be judged against the realised ex-post real rate:

  * ex-post real quarterly rate r_t = i_q - infl_q (full IT sample and last 5 years);
  * HP-filter trend of r_t (end-of-sample value as a slow-moving r* proxy).

These are descriptive anchors (not a Laubach-Williams estimate); the report makes that
explicit. Output: outputs/tables/rstar_estimates.csv.
"""

from __future__ import annotations

import json
import numpy as np
import pandas as pd

from common import DATA_CLEAN, TABLES, ensure_directories


def annualise(q: float) -> float:
    return ((1.0 + q) ** 4 - 1.0) * 100.0


def main() -> None:
    ensure_directories()
    panel = pd.read_csv(DATA_CLEAN / "chile_macro_quarterly.csv")
    metadata = json.loads(
        (DATA_CLEAN / "dataset_metadata.json").read_text(encoding="utf-8")
    )
    is_synthetic = bool(metadata.get("is_synthetic", True))
    data_source = str(metadata.get("institution", metadata.get("source", "unknown")))
    real_q = panel["real_rate_q"].dropna().to_numpy()

    from statsmodels.tsa.filters.hp_filter import hpfilter

    cycle, trend = hpfilter(panel["real_rate_q"].dropna(), lamb=1600)

    rows = []
    full_mean = float(np.mean(real_q))
    rows.append(
        {
            "method": "ex_post_real_rate_mean_full_sample",
            "rstar_quarterly": full_mean,
            "rstar_annual_pct": annualise(full_mean),
            "beta_implied": 1.0 / (1.0 + full_mean),
            "observations": int(len(real_q)),
        }
    )
    last20 = real_q[-20:]
    last_mean = float(np.mean(last20))
    rows.append(
        {
            "method": "ex_post_real_rate_mean_last_5y",
            "rstar_quarterly": last_mean,
            "rstar_annual_pct": annualise(last_mean),
            "beta_implied": 1.0 / (1.0 + last_mean),
            "observations": int(len(last20)),
        }
    )
    trend_end = float(trend.iloc[-1])
    rows.append(
        {
            "method": "hp_trend_real_rate_end_of_sample",
            "rstar_quarterly": trend_end,
            "rstar_annual_pct": annualise(trend_end),
            "beta_implied": 1.0 / (1.0 + trend_end),
            "observations": int(len(real_q)),
        }
    )

    table = pd.DataFrame(rows)
    table["is_synthetic"] = is_synthetic
    table["data_source"] = data_source
    table.to_csv(TABLES / "rstar_estimates.csv", index=False)
    print(table.round(4).to_string(index=False))
    print(f"Wrote {TABLES / 'rstar_estimates.csv'}")


if __name__ == "__main__":
    main()
\end{lstlisting}
\clearpage
\section{python/estimate\_time\_varying\_rstar.py}
\begin{lstlisting}[language=Python]
"""Estimate a transparent time-varying statistical proxy for Chilean r-star.

Expected inflation is generated recursively with an AR model, so the ex-ante
real policy rate does not use future inflation. A local-level state-space model
then separates a slow-moving latent component from high-frequency noise, with
the lagged output gap as a cyclical control.

This is deliberately labelled a proxy. It is not a full Laubach-Williams model:
potential growth, the IS relation and inflation expectations are not jointly
estimated in a structural system.
"""

from __future__ import annotations

import matplotlib

matplotlib.use("Agg")
import matplotlib.pyplot as plt
import numpy as np
import pandas as pd
import statsmodels.api as sm
from statsmodels.tsa.statespace.structural import UnobservedComponents

from common import DATA_CLEAN, FIGURES, TABLES, ensure_directories


def recursive_inflation_expectation(pi: np.ndarray, max_lags: int = 4) -> np.ndarray:
    expected = np.full(len(pi), np.nan)
    for t in range(len(pi)):
        if t < 8:
            expected[t] = float(np.mean(pi[: t + 1]))
            continue
        lags = min(max_lags, t - 2)
        y = pi[lags:t]
        x = np.column_stack(
            [pi[lags - lag - 1:t - lag - 1] for lag in range(lags)]
        )
        model = sm.OLS(y, sm.add_constant(x)).fit()
        predictors = np.r_[1.0, [pi[t - lag - 1] for lag in range(lags)]]
        expected[t] = float(model.predict(predictors)[0])
    return expected


def annualized_real_rate(nominal_q: np.ndarray, expected_inflation_q: np.ndarray) -> np.ndarray:
    return (((1.0 + nominal_q) / (1.0 + expected_inflation_q)) ** 4 - 1.0) * 100.0


def main() -> None:
    ensure_directories()
    macro = pd.read_csv(DATA_CLEAN / "chile_macro_quarterly.csv")
    pi_q = macro["infl_q"].to_numpy(dtype=float)
    expected_pi_q = recursive_inflation_expectation(pi_q)
    real_ex_ante = annualized_real_rate(
        macro["i_q"].to_numpy(dtype=float), expected_pi_q
    )
    gap_lag = (100.0 * macro["output_gap"]).shift(1).fillna(0.0)

    model = UnobservedComponents(
        real_ex_ante,
        level="local level",
        exog=pd.DataFrame({"output_gap_lag_pp": gap_lag}),
    )
    result = model.fit(disp=False, maxiter=2000)
    smoothed = np.asarray(result.level.smoothed)
    filtered = np.asarray(result.level.filtered)
    variance = np.maximum(np.asarray(result.level.smoothed_cov), 0.0)
    se = np.sqrt(variance)

    output = pd.DataFrame(
        {
            "date": macro["date"],
            "policy_rate_annual_pct": macro["tpm_annual_pct"],
            "expected_inflation_annual_pct": ((1.0 + expected_pi_q) ** 4 - 1.0) * 100.0,
            "real_rate_ex_ante_annual_pct": real_ex_ante,
            "rstar_filtered_pct": filtered,
            "rstar_smoothed_pct": smoothed,
            "rstar_lower_95_pct": smoothed - 1.96 * se,
            "rstar_upper_95_pct": smoothed + 1.96 * se,
            "output_gap_pct": 100.0 * macro["output_gap"],
        }
    )
    output.to_csv(TABLES / "rstar_time_varying.csv", index=False)

    summary = pd.DataFrame(
        [
            {
                "method": "local_level_state_space_proxy",
                "sample_start": macro["date"].iloc[0],
                "sample_end": macro["date"].iloc[-1],
                "observations": len(macro),
                "converged": bool(result.mle_retvals.get("converged", False)),
                "latest_rstar_smoothed_pct": float(smoothed[-1]),
                "latest_rstar_filtered_pct": float(filtered[-1]),
                "sample_mean_rstar_pct": float(np.mean(smoothed)),
                "sample_min_rstar_pct": float(np.min(smoothed)),
                "sample_max_rstar_pct": float(np.max(smoothed)),
                "variance_irregular": float(result.params["sigma2.irregular"]),
                "variance_level": float(result.params["sigma2.level"]),
                "gap_coefficient": float(result.params["beta.output_gap_lag_pp"]),
                "fixed_grid_low_pct": 2.0,
                "fixed_grid_baseline_pct": 3.0,
                "fixed_grid_high_pct": 4.0,
                "interpretation": "statistical proxy; not structural Laubach-Williams r-star",
            }
        ]
    )
    summary.to_csv(TABLES / "rstar_time_varying_summary.csv", index=False)

    x = np.arange(len(macro))
    ticks = np.arange(0, len(macro), 8)
    labels = macro.loc[ticks, "date"]
    fig, axes = plt.subplots(2, 1, figsize=(12, 8), sharex=True)
    axes[0].plot(x, real_ex_ante, color="#9CA3AF", linewidth=1,
                 label="Juro real ex ante (proxy)")
    axes[0].fill_between(
        x,
        output["rstar_lower_95_pct"],
        output["rstar_upper_95_pct"],
        color="#93C5FD",
        alpha=0.35,
        label="IC 95% do nivel latente",
    )
    axes[0].plot(x, smoothed, color="#1E4E79", linewidth=2.2,
                 label="r* suavizado (proxy)")
    axes[0].plot(x, filtered, color="#16A34A", linewidth=1.4, linestyle="--",
                 label="r* filtrado")
    axes[0].axhspan(2.0, 4.0, color="#F59E0B", alpha=0.10, label="Grade fixa 2%-4%")
    axes[0].axhline(3.0, color="#B45309", linestyle=":", linewidth=1.3)
    axes[0].set_ylabel("% a.a.")
    axes[0].set_title("Taxa real ex ante e componente neutro variavel")
    axes[0].legend(ncol=3, fontsize=8)
    axes[0].grid(alpha=0.25)

    axes[1].plot(x, macro["infl_annual_pct"], color="#B45309", linewidth=1.5,
                 label="Inflacao realizada anualizada")
    axes[1].plot(x, output["expected_inflation_annual_pct"], color="#1E4E79",
                 linewidth=1.5, label="Inflacao esperada AR recursiva")
    axes[1].axhline(3.0, color="black", linestyle=":", linewidth=1,
                    label="Meta de 3%")
    axes[1].set_ylabel("% a.a.")
    axes[1].set_title("Proxy de expectativas usada no juro real ex ante")
    axes[1].legend(ncol=3, fontsize=8)
    axes[1].grid(alpha=0.25)
    axes[1].set_xticks(ticks)
    axes[1].set_xticklabels(labels, rotation=70, fontsize=8)
    fig.suptitle("Chile: proxy estatistica de r* variavel no tempo")
    fig.text(
        0.01,
        0.01,
        "Modelo local-level com hiato defasado. Nao e uma estimativa estrutural "
        "Laubach-Williams; expectativas sao previsoes AR recursivas.",
        fontsize=8,
    )
    fig.tight_layout(rect=(0, 0.03, 1, 0.95))
    fig.savefig(FIGURES / "rstar_time_varying.png", dpi=180)
    plt.close(fig)
    print(summary.to_string(index=False))


if __name__ == "__main__":
    main()
\end{lstlisting}
\clearpage
\part*{Execucao (Octave/Dynare + Bayesiano)}\addcontentsline{toc}{part}{Execucao (Octave/Dynare + Bayesiano)}

\section{python/run\_dynare\_batch.py}
\begin{lstlisting}[language=Python]
"""Locate Octave and Dynare and run the Dynare models robustly on Windows.

To avoid the OneDrive file-locking and accented-path problems that occur when Dynare
preprocesses inside the synced project folder, this wrapper:

  1. stages the requested .mod files, the export helpers and dynare/run_models.m into a
     disposable ASCII temp directory (copy is done in Python, which reads the accented
     source paths reliably);
  2. runs Octave entirely inside that ASCII directory with NK_REPO_ROOT pointing there;
  3. copies the resulting CSVs back into outputs/dynare/ in the repository.

By default only the baseline runs; pass --all to also run every generated scenario.
A hard timeout terminates the Octave process tree.
"""

from __future__ import annotations

import argparse
import os
import shutil
import subprocess
import sys
import tempfile
from pathlib import Path

from common import (
    DYNARE_OUTPUTS,
    LOGS,
    ROOT,
    ensure_directories,
    find_dynare_matlab,
    find_octave,
)


def octave_literal(path) -> str:
    return str(path).replace("\\", "/").replace("'", "''")


def stage_work_dir(run_all: bool) -> Path:
    work = Path(tempfile.gettempdir()) / "nk_dynare_work"
    if work.exists():
        shutil.rmtree(work, ignore_errors=True)
    (work / "outputs" / "dynare").mkdir(parents=True, exist_ok=True)
    (work / "outputs" / "logs").mkdir(parents=True, exist_ok=True)
    dynare_dir = ROOT / "dynare"
    for helper in ("export_results.m", "export_stability.m", "run_models.m"):
        shutil.copy(dynare_dir / helper, work / helper)
    shutil.copy(dynare_dir / "nk_chile_base.mod", work / "nk_chile_base.mod")
    if run_all:
        for mod in sorted((dynare_dir / "generated").glob("*.mod")):
            shutil.copy(mod, work / mod.name)
    return work


def copy_results_back(work: Path) -> list[str]:
    staged = work / "outputs" / "dynare"
    copied = []
    if not staged.exists():
        return copied
    for scenario_dir in sorted(staged.iterdir()):
        if not scenario_dir.is_dir():
            continue
        dest = DYNARE_OUTPUTS / scenario_dir.name
        dest.mkdir(parents=True, exist_ok=True)
        for csv in scenario_dir.glob("*.csv"):
            shutil.copy(csv, dest / csv.name)
        copied.append(scenario_dir.name)
    return copied


def kill_tree(proc: subprocess.Popen) -> None:
    try:
        subprocess.run(
            ["taskkill", "/F", "/T", "/PID", str(proc.pid)],
            capture_output=True,
            check=False,
        )
    except Exception:  # noqa: BLE001
        proc.kill()


def main() -> int:
    parser = argparse.ArgumentParser()
    parser.add_argument("--all", action="store_true",
                        help="Run the baseline and every generated scenario.")
    parser.add_argument("--timeout", type=int, default=None,
                        help="Maximum runtime in seconds (default 150 baseline, 600 all).")
    parser.add_argument("--strict", action="store_true",
                        help="Return nonzero when Octave or Dynare is unavailable.")
    args = parser.parse_args()
    ensure_directories()

    octave = find_octave()
    dynare_matlab = find_dynare_matlab()
    if octave is None or dynare_matlab is None:
        message = (
            f"Dynare batch skipped. Octave found: {octave is not None}; "
            f"Dynare found: {dynare_matlab is not None}. "
            "Set OCTAVE_BIN and DYNARE_MATLAB if installed elsewhere."
        )
        print(f"WARNING: {message}")
        (LOGS / "dynare_batch.log").write_text(message + "\n", encoding="utf-8")
        return 1 if args.strict else 0

    work = stage_work_dir(args.all)
    expression = (
        f"addpath('{octave_literal(dynare_matlab)}'); run('run_models.m');"
    )
    command = [str(octave), "--quiet", "--eval", expression]
    environment = os.environ.copy()
    environment["NK_REPO_ROOT"] = str(work)
    timeout = args.timeout or (600 if args.all else 150)
    print(f"Octave: {octave}")
    print(f"Work dir (ASCII): {work}")
    print(f"Mode: {'all scenarios' if args.all else 'baseline only'}; timeout {timeout}s")

    proc = subprocess.Popen(
        command, cwd=str(work), env=environment, text=True,
        stdout=subprocess.PIPE, stderr=subprocess.PIPE,
    )
    try:
        stdout, stderr = proc.communicate(timeout=timeout)
        timed_out = False
    except subprocess.TimeoutExpired:
        kill_tree(proc)
        stdout, stderr = proc.communicate()
        timed_out = True

    copied = copy_results_back(work)
    log_text = (
        f"COMMAND: {' '.join(command)}\nWORK: {work}\nTIMEOUT: {timed_out}\n"
        f"COPIED: {copied}\n\nSTDOUT\n{stdout}\n\nSTDERR\n{stderr}\n"
    )
    (LOGS / "dynare_wrapper.log").write_text(log_text, encoding="utf-8")
    print(stdout[-2000:] if stdout else "")
    if stderr:
        print(stderr[-1500:], file=sys.stderr)
    print(f"Scenarios copied back: {copied}")
    if timed_out:
        print(f"ERROR: Octave exceeded {timeout}s and was terminated.", file=sys.stderr)
        return 124
    return proc.returncode if proc.returncode is not None else 0


if __name__ == "__main__":
    raise SystemExit(main())
\end{lstlisting}
\clearpage
\section{python/run\_bayesian.py}
\begin{lstlisting}[language=Python]
"""Run posterior-mode estimation robustly via ASCII staging.

Stages nk_chile_estim.mod, export_bayesian.m and a numeric-only observables CSV into a
disposable ASCII temp directory (to avoid OneDrive/accent issues), runs the Dynare
estimation there, and copies the posterior CSV back to outputs/tables/. A hard timeout
terminates the Octave process tree.
"""

from __future__ import annotations

import argparse
import json
import os
import shutil
import subprocess
import sys
import tempfile
from pathlib import Path

import pandas as pd

from common import (
    DATA_CLEAN,
    LOGS,
    ROOT,
    TABLES,
    ensure_directories,
    find_dynare_matlab,
    find_octave,
)


def octave_literal(path) -> str:
    return str(path).replace("\\", "/").replace("'", "''")


def kill_tree(proc: subprocess.Popen) -> None:
    try:
        subprocess.run(["taskkill", "/F", "/T", "/PID", str(proc.pid)],
                       capture_output=True, check=False)
    except Exception:  # noqa: BLE001
        proc.kill()


def augment_with_std(work: Path, produced: Path) -> None:
    """Add Laplace standard errors and 95% intervals from the saved mode Hessian.

    Dynare saves the posterior mode and Hessian (hh) in <model>_mode.mat. The rows of
    xparam1/hh follow the same order as the CSV written by export_bayesian.m, so we map
    by row index. Done in Python (scipy) because it is far more robust than reading the
    Hessian out of Octave structures.
    """
    if not produced.exists():
        return
    mode_files = list(work.rglob("*_mode.mat"))
    if not mode_files:
        print("No mode .mat found; reporting posterior mode without standard errors.")
        return
    try:
        import numpy as np
        from scipy.io import loadmat

        frame = pd.read_csv(produced)
        payload = loadmat(str(mode_files[0]))
        hessian = payload.get("hh")
        if hessian is None or hessian.shape[0] != len(frame):
            print(f"Mode Hessian shape {None if hessian is None else hessian.shape} "
                  f"!= {len(frame)} rows; skipping std.")
            return
        std = np.sqrt(np.diag(np.linalg.inv(hessian)))
        xparam = payload.get("xparam1")
        if xparam is not None:
            xparam = np.asarray(xparam).ravel()
            if len(xparam) == len(frame):
                frame["posterior_mode"] = xparam  # authoritative mode (all params, in order)
        frame["posterior_std"] = std
        frame["ci95_low"] = frame["posterior_mode"] - 1.96 * std
        frame["ci95_high"] = frame["posterior_mode"] + 1.96 * std
        frame.to_csv(produced, index=False)
        print(f"Added mode/std from {mode_files[0].name}.")
    except Exception as exc:  # noqa: BLE001
        print(f"Standard-error augmentation failed ({exc}); reporting mode only.")


def main() -> int:
    parser = argparse.ArgumentParser()
    parser.add_argument("--timeout", type=int, default=900)
    args = parser.parse_args()
    ensure_directories()

    octave = find_octave()
    dynare_matlab = find_dynare_matlab()
    if octave is None or dynare_matlab is None:
        print("WARNING: Octave/Dynare not found; Bayesian estimation skipped.")
        return 0

    observables = DATA_CLEAN / "chile_observables.csv"
    if not observables.exists():
        print("ERROR: chile_observables.csv missing; run build_chile_dataset.py first.",
              file=sys.stderr)
        return 1

    work = Path(tempfile.gettempdir()) / "nk_estim_work"
    if work.exists():
        shutil.rmtree(work, ignore_errors=True)
    (work / "outputs" / "tables").mkdir(parents=True, exist_ok=True)
    (work / "outputs" / "logs").mkdir(parents=True, exist_ok=True)
    shutil.copy(ROOT / "dynare" / "nk_chile_estim.mod", work / "nk_chile_estim.mod")
    shutil.copy(ROOT / "dynare" / "export_bayesian.m", work / "export_bayesian.m")
    # Numeric-only observables for Dynare's CSV reader (no date column).
    pd.read_csv(observables)[["pi", "i", "x"]].to_csv(
        work / "chile_observables_dynare.csv", index=False
    )

    expression = (
        f"addpath('{octave_literal(dynare_matlab)}'); "
        "dynare('nk_chile_estim.mod','noclearall','nolog');"
    )
    command = [str(octave), "--quiet", "--eval", expression]
    environment = os.environ.copy()
    environment["NK_REPO_ROOT"] = str(work)
    print(f"Octave: {octave}\nWork dir: {work}\nTimeout: {args.timeout}s")

    proc = subprocess.Popen(command, cwd=str(work), env=environment, text=True,
                            stdout=subprocess.PIPE, stderr=subprocess.PIPE)
    try:
        stdout, stderr = proc.communicate(timeout=args.timeout)
        timed_out = False
    except subprocess.TimeoutExpired:
        kill_tree(proc)
        stdout, stderr = proc.communicate()
        timed_out = True

    produced = work / "outputs" / "tables" / "bayesian_estimates.csv"
    augment_with_std(work, produced)
    copied = False
    if produced.exists():
        metadata = json.loads(
            (DATA_CLEAN / "dataset_metadata.json").read_text(encoding="utf-8")
        )
        frame = pd.read_csv(produced)
        frame["is_synthetic"] = bool(metadata.get("is_synthetic", True))
        frame["data_source"] = str(
            metadata.get("institution", metadata.get("source", "unknown"))
        )
        frame["estimation_scope"] = "posterior_mode_laplace_not_mcmc"
        frame.to_csv(produced, index=False)
        shutil.copy(produced, TABLES / "bayesian_estimates.csv")
        copied = True

    (LOGS / "bayesian_run.log").write_text(
        f"TIMEOUT: {timed_out}\nCOPIED: {copied}\nRC: {proc.returncode}\n\n"
        f"STDOUT\n{stdout}\n\nSTDERR\n{stderr}\n",
        encoding="utf-8",
    )
    print(stdout[-2500:] if stdout else "")
    if stderr:
        print(stderr[-1200:], file=sys.stderr)
    print(f"bayesian_estimates.csv copied back: {copied}; timed_out: {timed_out}")
    if timed_out and not copied:
        return 124
    return 0


if __name__ == "__main__":
    raise SystemExit(main())
\end{lstlisting}
\clearpage
\section{python/run\_mcmc.py}
\begin{lstlisting}[language=Python]
"""Run the full Dynare MCMC extension and compute chain diagnostics."""

from __future__ import annotations

import argparse
import os
import re
import shutil
import subprocess
import sys
import tempfile
from pathlib import Path

import numpy as np
import pandas as pd
from scipy.io import loadmat

from common import DATA_CLEAN, LOGS, ROOT, TABLES, find_dynare_matlab, find_octave


PARAMETER_ORDER = [
    "stderr e_x",
    "stderr e_pi",
    "stderr e_i",
    "sigma",
    "kappa",
    "rho_i",
    "phi_pi",
    "phi_x",
]


def octave_literal(path: Path) -> str:
    return str(path).replace("\\", "/").replace("'", "''")


def split_rhat(chains: list[np.ndarray]) -> np.ndarray:
    if len(chains) < 2:
        return np.full(chains[0].shape[1], np.nan)
    n = min(len(chain) for chain in chains)
    arrays = np.stack([chain[-n:] for chain in chains])
    chain_means = arrays.mean(axis=1)
    between = n * chain_means.var(axis=0, ddof=1)
    within = arrays.var(axis=1, ddof=1).mean(axis=0)
    variance = ((n - 1) / n) * within + between / n
    return np.sqrt(variance / within)


def effective_sample_size(chains: list[np.ndarray]) -> np.ndarray:
    n = min(len(chain) for chain in chains)
    arrays = np.stack([chain[-n:] for chain in chains])
    m, _, k = arrays.shape
    result = np.empty(k)
    for parameter in range(k):
        autocorrelations = []
        for chain in arrays[:, :, parameter]:
            centered = chain - chain.mean()
            fft_size = 1 << (2 * n - 1).bit_length()
            spectrum = np.fft.rfft(centered, fft_size)
            covariance = np.fft.irfft(spectrum * spectrum.conjugate(), fft_size)[:n]
            covariance /= np.arange(n, 0, -1)
            autocorrelations.append(covariance / covariance[0])
        mean_acf = np.mean(autocorrelations, axis=0)
        autocorr_sum = 0.0
        for rho in mean_acf[1 : min(1001, n)]:
            if rho <= 0:
                break
            autocorr_sum += float(rho)
        result[parameter] = m * n / (1.0 + 2.0 * autocorr_sum)
    return result


def load_chains(work: Path) -> list[np.ndarray]:
    files = sorted(work.rglob("*_mh*_blck*.mat"))
    grouped: dict[int, list[np.ndarray]] = {}
    for path in files:
        match = re.search(r"blck(\d+)", path.name)
        block = int(match.group(1)) if match else len(grouped) + 1
        payload = loadmat(path)
        draws = payload.get("x2")
        if draws is None or draws.ndim != 2:
            continue
        grouped.setdefault(block, []).append(np.asarray(draws, dtype=float))
    return [np.vstack(parts) for _, parts in sorted(grouped.items())]


def save_chain_outputs(chains: list[np.ndarray], burn_fraction: float = 0.30) -> None:
    retained = [chain[int(len(chain) * burn_fraction) :] for chain in chains]
    draws = np.vstack(retained)
    posterior = pd.DataFrame(
        {
            "parameter": PARAMETER_ORDER,
            "type": ["shock_std"] * 3 + ["param"] * 5,
            "posterior_mean": draws.mean(axis=0),
            "posterior_median": np.median(draws, axis=0),
            "posterior_std": draws.std(axis=0, ddof=1),
            "hpd90_low": np.quantile(draws, 0.05, axis=0),
            "hpd90_high": np.quantile(draws, 0.95, axis=0),
        }
    )
    posterior.to_csv(TABLES / "mcmc_posterior.csv", index=False)

    diagnostics = pd.DataFrame(
        {
            "parameter": PARAMETER_ORDER,
            "rhat": split_rhat(retained),
            "effective_sample_size": effective_sample_size(retained),
            "chains": len(retained),
            "draws_per_chain_loaded": min(len(chain) for chain in retained),
        }
    )
    diagnostics["rhat_pass_1p05"] = diagnostics["rhat"] < 1.05
    diagnostics.to_csv(TABLES / "mcmc_diagnostics.csv", index=False)
    print(diagnostics.to_string(index=False))


def main() -> int:
    parser = argparse.ArgumentParser()
    parser.add_argument("--timeout", type=int, default=1800)
    parser.add_argument(
        "--recover-dir",
        type=Path,
        help="Post-process completed Dynare chain files without rerunning MCMC.",
    )
    args = parser.parse_args()
    if args.recover_dir:
        chains = load_chains(args.recover_dir)
        if len(chains) < 2 or any(
            chain.shape[1] != len(PARAMETER_ORDER) for chain in chains
        ):
            print(f"ERROR: invalid MCMC chain shapes: {[c.shape for c in chains]}")
            return 1
        save_chain_outputs(chains)
        print(f"Recovered MCMC outputs from {args.recover_dir}")
        return 0

    octave = find_octave()
    dynare = find_dynare_matlab()
    if octave is None or dynare is None:
        print("ERROR: Octave/Dynare not found.", file=sys.stderr)
        return 1

    work = Path(tempfile.gettempdir()) / "nk_mcmc_work"
    if work.exists():
        shutil.rmtree(work, ignore_errors=True)
    (work / "outputs" / "tables").mkdir(parents=True, exist_ok=True)
    for filename in ["nk_chile_mcmc.mod", "export_mcmc.m"]:
        shutil.copy(ROOT / "dynare" / filename, work / filename)
    pd.read_csv(DATA_CLEAN / "chile_observables.csv")[["pi", "i", "x"]].to_csv(
        work / "chile_observables_dynare.csv", index=False
    )

    expression = (
        f"addpath('{octave_literal(dynare)}'); "
        "dynare('nk_chile_mcmc.mod','noclearall','nolog');"
    )
    command = [str(octave), "--quiet", "--eval", expression]
    env = os.environ.copy()
    env["NK_REPO_ROOT"] = str(work)
    proc = subprocess.Popen(
        command,
        cwd=work,
        env=env,
        text=True,
        stdout=subprocess.PIPE,
        stderr=subprocess.PIPE,
    )
    try:
        stdout, stderr = proc.communicate(timeout=args.timeout)
        timed_out = False
    except subprocess.TimeoutExpired:
        subprocess.run(
            ["taskkill", "/F", "/T", "/PID", str(proc.pid)],
            capture_output=True,
            check=False,
        )
        stdout, stderr = proc.communicate()
        timed_out = True

    posterior = work / "outputs" / "tables" / "mcmc_posterior.csv"
    copied = False
    if posterior.exists():
        shutil.copy(posterior, TABLES / posterior.name)
        copied = True

    chains = load_chains(work)
    if chains and all(chain.shape[1] == len(PARAMETER_ORDER) for chain in chains):
        save_chain_outputs(chains)
    else:
        print(
            f"WARNING: could not load two {len(PARAMETER_ORDER)}-column MCMC chains; "
            f"shapes={[chain.shape for chain in chains]}",
            file=sys.stderr,
        )

    (LOGS / "mcmc_run.log").write_text(
        f"TIMEOUT: {timed_out}\nRC: {proc.returncode}\nCOPIED: {copied}\n"
        f"CHAIN_SHAPES: {[chain.shape for chain in chains]}\n\n"
        f"STDOUT\n{stdout}\n\nSTDERR\n{stderr}\n",
        encoding="utf-8",
    )
    print(stdout[-4000:] if stdout else "")
    if stderr:
        print(stderr[-1600:], file=sys.stderr)
    return 0 if copied and not timed_out else 1


if __name__ == "__main__":
    raise SystemExit(main())
\end{lstlisting}
\clearpage
\section{python/run\_forecast.py}
\begin{lstlisting}[language=Python]
"""Run the native Dynare forecast (estimation forecast=8 + conditional_forecast).

Mirrors run_bayesian.py: stages nk_chile_forecast.mod, the export helpers and a
numeric-only observables CSV into a disposable ASCII temp directory (to avoid the
OneDrive/accent issues that break Dynare's preprocessing inside the synced folder),
runs Dynare there, and copies the produced CSVs back to outputs/dynare/forecast/.
"""

from __future__ import annotations

import argparse
import os
import shutil
import subprocess
import sys
import tempfile
from pathlib import Path

import pandas as pd

from common import (
    DATA_CLEAN,
    DYNARE_OUTPUTS,
    LOGS,
    ROOT,
    ensure_directories,
    find_dynare_matlab,
    find_octave,
)


def octave_literal(path) -> str:
    return str(path).replace("\\", "/").replace("'", "''")


def kill_tree(proc: subprocess.Popen) -> None:
    try:
        subprocess.run(["taskkill", "/F", "/T", "/PID", str(proc.pid)],
                       capture_output=True, check=False)
    except Exception:  # noqa: BLE001
        proc.kill()


def copy_results_back(work: Path) -> list[str]:
    staged = work / "outputs" / "dynare" / "forecast"
    if not staged.exists():
        return []
    dest = DYNARE_OUTPUTS / "forecast"
    dest.mkdir(parents=True, exist_ok=True)
    copied = []
    for item in sorted(staged.glob("*.csv")) + sorted(staged.glob("*.txt")):
        shutil.copy(item, dest / item.name)
        copied.append(item.name)
    return copied


def main() -> int:
    parser = argparse.ArgumentParser()
    parser.add_argument("--timeout", type=int, default=600)
    args = parser.parse_args()
    ensure_directories()

    octave = find_octave()
    dynare_matlab = find_dynare_matlab()
    if octave is None or dynare_matlab is None:
        print("WARNING: Octave/Dynare not found; forecast skipped.", file=sys.stderr)
        return 1

    observables = DATA_CLEAN / "chile_observables.csv"
    if not observables.exists():
        print("ERROR: chile_observables.csv missing; run build_chile_dataset.py first.",
              file=sys.stderr)
        return 1

    work = Path(tempfile.gettempdir()) / "nk_forecast_work"
    if work.exists():
        shutil.rmtree(work, ignore_errors=True)
    (work / "outputs" / "dynare" / "forecast").mkdir(parents=True, exist_ok=True)
    (work / "outputs" / "logs").mkdir(parents=True, exist_ok=True)

    dynare_dir = ROOT / "dynare"
    for helper in ("nk_chile_forecast.mod", "export_forecast.m", "export_conditional.m"):
        shutil.copy(dynare_dir / helper, work / helper)
    # Numeric-only observables for Dynare's CSV reader (no date column).
    pd.read_csv(observables)[["pi", "i", "x"]].to_csv(
        work / "chile_observables_dynare.csv", index=False
    )

    expression = (
        f"addpath('{octave_literal(dynare_matlab)}'); "
        "dynare('nk_chile_forecast.mod','noclearall','nolog');"
    )
    command = [str(octave), "--quiet", "--eval", expression]
    environment = os.environ.copy()
    environment["NK_REPO_ROOT"] = str(work)
    print(f"Octave: {octave}\nWork dir: {work}\nTimeout: {args.timeout}s")

    proc = subprocess.Popen(command, cwd=str(work), env=environment, text=True,
                            stdout=subprocess.PIPE, stderr=subprocess.PIPE)
    try:
        stdout, stderr = proc.communicate(timeout=args.timeout)
        timed_out = False
    except subprocess.TimeoutExpired:
        kill_tree(proc)
        stdout, stderr = proc.communicate()
        timed_out = True

    copied = copy_results_back(work)
    (LOGS / "forecast_run.log").write_text(
        f"TIMEOUT: {timed_out}\nRC: {proc.returncode}\nCOPIED: {copied}\n\n"
        f"STDOUT\n{stdout}\n\nSTDERR\n{stderr}\n",
        encoding="utf-8",
    )
    print(stdout[-3000:] if stdout else "")
    if stderr:
        print(stderr[-1500:], file=sys.stderr)
    print(f"Forecast CSVs copied back: {copied}; timed_out: {timed_out}")
    if timed_out and not copied:
        return 124
    return 0


if __name__ == "__main__":
    raise SystemExit(main())
\end{lstlisting}
\clearpage
\section{python/run\_history.py}
\begin{lstlisting}[language=Python]
"""Run the historical shock decomposition (Kalman smoother + shock_decomposition).

Stages nk_chile_history.mod, export_history.m and a numeric observables CSV into
an ASCII temp directory, runs Dynare, and copies outputs/dynare/history/*.csv back.
Mirrors run_forecast.py.
"""

from __future__ import annotations

import argparse
import os
import shutil
import subprocess
import sys
import tempfile
from pathlib import Path

import pandas as pd

from common import (
    DATA_CLEAN,
    DYNARE_OUTPUTS,
    LOGS,
    ROOT,
    ensure_directories,
    find_dynare_matlab,
    find_octave,
)


def octave_literal(path) -> str:
    return str(path).replace("\\", "/").replace("'", "''")


def kill_tree(proc: subprocess.Popen) -> None:
    try:
        subprocess.run(["taskkill", "/F", "/T", "/PID", str(proc.pid)],
                       capture_output=True, check=False)
    except Exception:  # noqa: BLE001
        proc.kill()


def main() -> int:
    parser = argparse.ArgumentParser()
    parser.add_argument("--timeout", type=int, default=420)
    parser.add_argument("--mod", default="nk_chile_history.mod")
    args = parser.parse_args()
    ensure_directories()

    octave = find_octave()
    dynare_matlab = find_dynare_matlab()
    if octave is None or dynare_matlab is None:
        print("WARNING: Octave/Dynare not found; history skipped.", file=sys.stderr)
        return 1

    observables = DATA_CLEAN / "chile_observables.csv"
    if not observables.exists():
        print("ERROR: chile_observables.csv missing.", file=sys.stderr)
        return 1

    work = Path(tempfile.gettempdir()) / "nk_history_work"
    if work.exists():
        shutil.rmtree(work, ignore_errors=True)
    (work / "outputs" / "dynare").mkdir(parents=True, exist_ok=True)
    (work / "outputs" / "logs").mkdir(parents=True, exist_ok=True)

    dynare_dir = ROOT / "dynare"
    for helper in (args.mod, "export_history.m"):
        shutil.copy(dynare_dir / helper, work / helper)
    pd.read_csv(observables)[["pi", "i", "x"]].to_csv(
        work / "chile_observables_dynare.csv", index=False
    )

    expression = (
        f"addpath('{octave_literal(dynare_matlab)}'); "
        f"dynare('{args.mod}','noclearall','nolog');"
    )
    command = [str(octave), "--quiet", "--eval", expression]
    environment = os.environ.copy()
    environment["NK_REPO_ROOT"] = str(work)
    print(f"Octave: {octave}\nWork dir: {work}\nTimeout: {args.timeout}s")

    proc = subprocess.Popen(command, cwd=str(work), env=environment, text=True,
                            stdout=subprocess.PIPE, stderr=subprocess.PIPE)
    try:
        stdout, stderr = proc.communicate(timeout=args.timeout)
        timed_out = False
    except subprocess.TimeoutExpired:
        kill_tree(proc)
        stdout, stderr = proc.communicate()
        timed_out = True

    staged_root = work / "outputs" / "dynare"
    copied = []
    if staged_root.exists():
        for sub in sorted(staged_root.iterdir()):
            if not sub.is_dir():
                continue
            dest = DYNARE_OUTPUTS / sub.name
            dest.mkdir(parents=True, exist_ok=True)
            for csv in sorted(sub.glob("*.csv")):
                shutil.copy(csv, dest / csv.name)
                copied.append(f"{sub.name}/{csv.name}")

    (LOGS / "history_run.log").write_text(
        f"TIMEOUT: {timed_out}\nRC: {proc.returncode}\nCOPIED: {copied}\n\n"
        f"STDOUT\n{stdout}\n\nSTDERR\n{stderr}\n",
        encoding="utf-8",
    )
    print(stdout[-2500:] if stdout else "")
    if stderr:
        print(stderr[-1200:], file=sys.stderr)
    print(f"History CSVs copied back: {copied}; timed_out: {timed_out}")
    return 0 if copied else (124 if timed_out else 1)


if __name__ == "__main__":
    raise SystemExit(main())
\end{lstlisting}
\clearpage
\section{python/run\_macro\_extensions.py}
\begin{lstlisting}[language=Python]
"""Run the open-economy and hybrid-NKPC Dynare extensions via ASCII staging."""

from __future__ import annotations

import argparse
import os
import shutil
import subprocess
import sys
import tempfile
from pathlib import Path

from common import DYNARE_OUTPUTS, LOGS, ROOT, find_dynare_matlab, find_octave


def octave_literal(path: Path) -> str:
    return str(path).replace("\\", "/").replace("'", "''")


def kill_tree(proc: subprocess.Popen) -> None:
    subprocess.run(
        ["taskkill", "/F", "/T", "/PID", str(proc.pid)],
        capture_output=True,
        check=False,
    )


def main() -> int:
    parser = argparse.ArgumentParser()
    parser.add_argument("--timeout", type=int, default=300)
    args = parser.parse_args()
    octave = find_octave()
    dynare = find_dynare_matlab()
    if octave is None or dynare is None:
        print("ERROR: Octave/Dynare not found.", file=sys.stderr)
        return 1

    work = Path(tempfile.gettempdir()) / "nk_macro_extensions"
    if work.exists():
        shutil.rmtree(work, ignore_errors=True)
    (work / "outputs" / "dynare").mkdir(parents=True, exist_ok=True)
    for name in [
        "nk_chile_open.mod",
        "nk_chile_hybrid.mod",
        "export_extension.m",
        "export_stability.m",
    ]:
        shutil.copy(ROOT / "dynare" / name, work / name)
    runner = (
        "addpath('" + octave_literal(dynare) + "');"
        "dynare('nk_chile_open.mod','noclearall','nolog');"
        "dynare('nk_chile_hybrid.mod','noclearall','nolog');"
    )
    command = [str(octave), "--quiet", "--eval", runner]
    env = os.environ.copy()
    env["NK_REPO_ROOT"] = str(work)
    proc = subprocess.Popen(
        command,
        cwd=work,
        env=env,
        text=True,
        stdout=subprocess.PIPE,
        stderr=subprocess.PIPE,
    )
    try:
        stdout, stderr = proc.communicate(timeout=args.timeout)
        timed_out = False
    except subprocess.TimeoutExpired:
        kill_tree(proc)
        stdout, stderr = proc.communicate()
        timed_out = True

    copied = []
    for scenario in ["open_economy", "hybrid_nkpc"]:
        source = work / "outputs" / "dynare" / scenario
        if not source.exists():
            continue
        destination = DYNARE_OUTPUTS / scenario
        destination.mkdir(parents=True, exist_ok=True)
        for item in source.glob("*"):
            if item.is_file():
                shutil.copy(item, destination / item.name)
        copied.append(scenario)
    (LOGS / "macro_extensions.log").write_text(
        f"TIMEOUT: {timed_out}\nRC: {proc.returncode}\nCOPIED: {copied}\n\n"
        f"STDOUT\n{stdout}\n\nSTDERR\n{stderr}\n",
        encoding="utf-8",
    )
    print(stdout[-3500:] if stdout else "")
    if stderr:
        print(stderr[-1500:], file=sys.stderr)
    print(f"Copied: {copied}; timeout={timed_out}; rc={proc.returncode}")
    return 0 if len(copied) == 2 and not timed_out else 1


if __name__ == "__main__":
    raise SystemExit(main())
\end{lstlisting}
\clearpage
\part*{Previsao e avaliacao}\addcontentsline{toc}{part}{Previsao e avaliacao}

\section{python/forecast\_model.py}
\begin{lstlisting}[language=Python]
"""Visualise and cross-check the native Dynare forecast (slides 48-55).

The forecast itself is produced by Dynare (python/run_forecast.py runs
dynare/nk_chile_forecast.mod):

* the UNCONDITIONAL projection comes from estimation(..., forecast=8); Dynare
  returns oo_.forecast.Mean and the HPDinf/HPDsup credibility bands, launched
  from the smoothed end-of-sample state. We read those CSVs and draw the fan
  chart spliced onto the observed history.
* the CONDITIONAL projections come from Dynare's conditional_forecast command
  (hold the rate, tighten, force inflation to target). Dynare's
  conditional_forecast launches from the steady state, so here we ALSO recompute
  the same three experiments anchored to the current state with the exact
  first-order reduced form (identical to Dynare's solution: the unconditional
  mean matches oo_.forecast to the last digit). This anchored version connects
  continuously to history and recovers the monetary-policy innovation e_i needed
  to enforce each path.

All series are mapped back to interpretable annualised levels using the sample
means that build_chile_dataset.py removed when it created the observables
(pi = infl_q - mean(infl_q); i = i_q - mean(i_q); x = output_gap - mean).
"""

from __future__ import annotations

import json

import matplotlib

matplotlib.use("Agg")
import matplotlib.pyplot as plt
import numpy as np
import pandas as pd

from common import (
    DATA_CLEAN,
    DYNARE_OUTPUTS,
    FIGURES,
    TABLES,
    calibration_rho,
    complete_parameters,
    ensure_directories,
    _state_coefficients,
)

HORIZON = 8
INFLATION_TARGET_ANNUAL = 0.03
FORECAST_DIR = DYNARE_OUTPUTS / "forecast"

# Palette (kept consistent with the rest of the figure set).
COLORS = {
    "unconditional": "#111827",
    "hold": "#FF6B35",
    "hike": "#16A34A",
    "pi_target": "#7C3AED",
}


def annual_pct(quarterly_fraction):
    return ((1.0 + quarterly_fraction) ** 4 - 1.0) * 100.0


def solution_matrices(params: dict) -> tuple[np.ndarray, np.ndarray]:
    """Return the reduced form  y_t = transition * i_{t-1} + impact * eps_t."""
    a_x, a_pi, a_i = _state_coefficients(params)
    sigma, beta, kappa = params["sigma"], params["beta"], params["kappa"]
    rho_i, phi_pi, phi_x = params["rho_i"], params["phi_pi"], params["phi_x"]
    c_i = a_x - (1.0 - a_pi) / sigma
    structural = np.array(
        [
            [1.0, 0.0, -c_i],
            [-kappa, 1.0, -beta * a_pi],
            [-(1.0 - rho_i) * phi_x, -(1.0 - rho_i) * phi_pi, 1.0],
        ]
    )
    transition = np.array([a_x, a_pi, a_i])
    impact = np.linalg.inv(structural)
    return transition, impact


def conditional_path(transition, impact, initial_i_gap, controlled_variable, targets):
    """Hit a constrained endogenous path using e_i, anchored at the current state."""
    index = {"x": 0, "pi": 1, "i": 2}[controlled_variable]
    policy_shock_index = 2
    values = np.zeros((HORIZON, 3))
    implied_policy_shocks = np.zeros(HORIZON)
    previous_i = initial_i_gap
    for horizon in range(HORIZON):
        mean = transition * previous_i
        innovations = np.zeros(3)
        if horizon in targets:
            denominator = impact[index, policy_shock_index]
            innovations[policy_shock_index] = (targets[horizon] - mean[index]) / denominator
        current = mean + impact @ innovations
        values[horizon] = current
        implied_policy_shocks[horizon] = innovations[policy_shock_index]
        previous_i = current[2]
    return values, implied_policy_shocks


def load_unconditional() -> pd.DataFrame:
    return pd.read_csv(FORECAST_DIR / "forecast_unconditional.csv")


def main() -> None:
    ensure_directories()
    params = complete_parameters({"rho_i": calibration_rho()})
    transition, impact = solution_matrices(params)

    macro = pd.read_csv(DATA_CLEAN / "chile_macro_quarterly.csv")
    observables = pd.read_csv(DATA_CLEAN / "chile_observables.csv")
    metadata = json.loads((DATA_CLEAN / "dataset_metadata.json").read_text(encoding="utf-8"))
    data_source = str(metadata.get("institution", metadata.get("source", "unknown")))

    # Level mapping: inflation is now centred on the TARGET, the rate on its mean.
    mean_infl_q = float(macro["infl_q"].mean())
    mean_i_q = float(macro["i_q"].mean())
    mean_tpm_annual = annual_pct(mean_i_q)
    mean_infl_annual = annual_pct(mean_infl_q)
    target_q = (1.0 + INFLATION_TARGET_ANNUAL) ** 0.25 - 1.0
    pi_target_dev = 0.0  # the target IS the steady state now (pi deviation = 0)

    def to_level(variable: str, deviation):
        if variable == "x":
            return 100.0 * np.asarray(deviation)
        if variable == "pi":
            return annual_pct(target_q + np.asarray(deviation))
        return annual_pct(mean_i_q + np.asarray(deviation))

    last_i_gap = float(observables.iloc[-1]["i"])
    last_period = pd.Period(str(observables.iloc[-1]["date"]), freq="Q")
    future = [str(last_period + h) for h in range(1, HORIZON + 1)]

    # ------------------------------------------------------------------ #
    # 1) Unconditional forecast: read Dynare's mean and HPD bands.        #
    # ------------------------------------------------------------------ #
    uncond = load_unconditional()
    uncond = uncond.assign(
        mean_level=lambda d: [to_level(v, m) for v, m in zip(d["variable"], d["mean"])],
        inf_level=lambda d: [to_level(v, m) for v, m in zip(d["variable"], d["hpd_inf"])],
        sup_level=lambda d: [to_level(v, m) for v, m in zip(d["variable"], d["hpd_sup"])],
        date=lambda d: [future[h - 1] for h in d["horizon"]],
    )

    # ------------------------------------------------------------------ #
    # 2) Conditional experiments anchored at the current state.          #
    # ------------------------------------------------------------------ #
    scenarios = {
        "hold": {
            "label": "Hold rate at current level (2q)",
            "controlled": "i",
            "targets": {0: last_i_gap, 1: last_i_gap},
        },
        "hike": {
            "label": "Tighten ~+1pp for 4q",
            "controlled": "i",
            "targets": {h: 0.0025 for h in range(4)},
        },
        "pi_target": {
            "label": "Inflation to 3% target for 4q",
            "controlled": "pi",
            "targets": {h: pi_target_dev for h in range(4)},
        },
    }
    # Unconditional reduced-form mean (matches Dynare oo_.forecast to last digit).
    uncond_dev = np.zeros((HORIZON, 3))
    prev = last_i_gap
    for h in range(HORIZON):
        uncond_dev[h] = transition * prev
        prev = uncond_dev[h, 2]

    results = {"unconditional": (uncond_dev, np.zeros(HORIZON))}
    for key, spec in scenarios.items():
        values, shocks = conditional_path(
            transition, impact, last_i_gap, spec["controlled"], spec["targets"]
        )
        results[key] = (values, shocks)

    # ------------------------------------------------------------------ #
    # 3) Tidy forecast table (deviations + levels), combining both views. #
    # ------------------------------------------------------------------ #
    var_index = {"x": 0, "pi": 1, "i": 2}
    records: list[dict] = []
    for scenario, (values, shocks) in results.items():
        for h in range(HORIZON):
            for variable, idx in var_index.items():
                row = {
                    "horizon": h + 1,
                    "date": future[h],
                    "scenario": scenario,
                    "variable": variable,
                    "value_dev": values[h, idx],
                    "value_level": float(to_level(variable, values[h, idx])),
                    "implied_e_i_q": shocks[h] if variable == "i" else np.nan,
                }
                if scenario == "unconditional":
                    band = uncond[(uncond["horizon"] == h + 1) & (uncond["variable"] == variable)]
                    row["p05_level"] = float(band["inf_level"].iloc[0]) if len(band) else np.nan
                    row["p95_level"] = float(band["sup_level"].iloc[0]) if len(band) else np.nan
                records.append(row)
    table = pd.DataFrame(records)
    table["data_source"] = data_source
    table["mean_tpm_annual_pct"] = mean_tpm_annual
    table["mean_infl_annual_pct"] = mean_infl_annual
    table.to_csv(TABLES / "forecast.csv", index=False)
    table[table["horizon"].isin([1, 4, 8])].to_csv(TABLES / "forecast_summary.csv", index=False)

    # ------------------------------------------------------------------ #
    # Figure 1: history + unconditional Dynare fan chart (annual levels). #
    # ------------------------------------------------------------------ #
    history = macro.tail(16).copy()
    fig, axes = plt.subplots(3, 1, figsize=(11, 9))
    panels = (
        ("i", history["tpm_annual_pct"], "Policy rate TPM (% per year)"),
        ("pi", history["infl_annual_pct"], "Inflation (% per year)"),
        ("x", 100.0 * history["output_gap"], "Output gap (%)"),
    )
    for ax, (variable, hist_series, title) in zip(axes, panels):
        ax.plot(history["date"], hist_series, color="black", marker="o",
                markersize=3, label="Observed history")
        sub = uncond[uncond["variable"] == variable]
        ax.plot(sub["date"], sub["mean_level"], color="#1E4E79", linewidth=2.4,
                label="Dynare forecast mean")
        ax.fill_between(sub["date"], sub["inf_level"], sub["sup_level"],
                        color="#1E4E79", alpha=0.16, label="Dynare HPD band")
        if variable == "pi":
            ax.axhline(3.0, color="#B45309", linestyle="--", linewidth=1, label="3% target")
            ax.axhline(mean_infl_annual, color="#6B7280", linestyle=":", linewidth=1,
                       label=f"Sample mean ({mean_infl_annual:.1f}%)")
        elif variable == "i":
            ax.axhline(mean_tpm_annual, color="#6B7280", linestyle=":", linewidth=1,
                       label=f"Sample mean ({mean_tpm_annual:.1f}%)")
        else:
            ax.axhline(0.0, color="#6B7280", linewidth=0.7)
        ax.set_title(title)
        ax.grid(alpha=0.25)
        ax.tick_params(axis="x", rotation=75, labelsize=7)
        ax.legend(fontsize=7, ncol=2)
    fig.suptitle("Chile: unconditional forecast (Dynare estimation, forecast=8)")
    fig.text(0.01, 0.005,
             f"Source: {data_source}. Native Dynare forecast launched from the smoothed "
             "end-of-sample state. Mechanical model projection, not an official BCCh forecast.",
             fontsize=7)
    fig.tight_layout(rect=(0, 0.03, 1, 0.96))
    fig.savefig(FIGURES / "forecast_fanchart.png", dpi=180)
    plt.close(fig)

    # ------------------------------------------------------------------ #
    # Figure 2: conditional scenarios anchored to the current state.     #
    # ------------------------------------------------------------------ #
    fig, axes = plt.subplots(1, 3, figsize=(14, 4.8))
    keymap = (("i", "Policy rate TPM (% per year)"),
              ("pi", "Inflation (% per year)"),
              ("x", "Output gap (%)"))
    for ax, (variable, title) in zip(axes, keymap):
        idx = var_index[variable]
        htail = macro.tail(4)
        hist_series = {"i": htail["tpm_annual_pct"], "pi": htail["infl_annual_pct"],
                       "x": 100.0 * htail["output_gap"]}[variable]
        ax.plot(range(-3, 1), hist_series, color="black", marker="o", markersize=3,
                linewidth=1.2, label="History")
        for scenario, (values, _) in results.items():
            color = COLORS["unconditional" if scenario == "unconditional" else scenario]
            label = "Unconditional" if scenario == "unconditional" else scenarios[scenario]["label"]
            level = to_level(variable, values[:, idx])
            # splice: connect the last observed point to the forecast path
            xs = list(range(0, HORIZON + 1))
            ys = [float(hist_series.iloc[-1])] + list(level)
            ax.plot(xs, ys, color=color, linewidth=2, label=label,
                    linestyle="-" if scenario == "unconditional" else "-")
        if variable == "pi":
            ax.axhline(3.0, color="#B45309", linestyle="--", linewidth=1)
        elif variable == "i":
            ax.axhline(mean_tpm_annual, color="#6B7280", linestyle=":", linewidth=1)
        else:
            ax.axhline(0.0, color="#6B7280", linewidth=0.8)
        ax.axvline(0, color="#9CA3AF", linewidth=0.8)
        ax.set_title(title)
        ax.set_xlabel("Quarter (0 = last data)")
        ax.grid(alpha=0.25)
    axes[0].legend(fontsize=7)
    fig.suptitle("Conditional scenarios (anchored at the current state; e_i solves each path)")
    fig.tight_layout(rect=(0, 0, 1, 0.94))
    fig.savefig(FIGURES / "conditional_scenarios.png", dpi=180)
    plt.close(fig)

    # ------------------------------------------------------------------ #
    # Figure 3: monetary-policy innovations e_i required by each path.   #
    # ------------------------------------------------------------------ #
    fig, ax = plt.subplots(figsize=(9.5, 4.8))
    for scenario in ("hold", "hike", "pi_target"):
        _, shocks = results[scenario]
        ax.plot(np.arange(1, HORIZON + 1), 100.0 * shocks, marker="o", linewidth=2,
                color=COLORS[scenario], label=scenarios[scenario]["label"])
    ax.axhline(0, color="black", linewidth=0.8)
    ax.set(title="Monetary-policy innovations e_i required by each conditional path",
           xlabel="Forecast quarter", ylabel="e_i (quarterly percentage points)")
    ax.grid(alpha=0.25)
    ax.legend(fontsize=8)
    fig.tight_layout()
    fig.savefig(FIGURES / "conditional_policy_shocks.png", dpi=180)
    plt.close(fig)

    print(f"Sample-mean TPM {mean_tpm_annual:.2f}% | mean inflation {mean_infl_annual:.2f}% "
          f"| 3% target as gap {pi_target_dev:+.5f}")
    print(f"Initial state: last policy-rate gap {100*last_i_gap:+.3f} quarterly pp")
    print(table[table["horizon"].isin([1, 4, 8])
                & table["variable"].isin(["pi", "i", "x"])].to_string(index=False))
    print("Wrote forecast tables and three figures (fan chart, scenarios, policy shocks).")


if __name__ == "__main__":
    main()
\end{lstlisting}
\clearpage
\section{python/evaluate\_forecasts.py}
\begin{lstlisting}[language=Python]
"""Pseudo-out-of-sample forecast evaluation for the Chilean variables.

Expanding-window forecasts compare the fixed calibrated NK solution with
univariate AR(1), random-walk and VAR benchmarks at one- and four-quarter
horizons. The exercise is pseudo-real-time: revised data and the full-sample HP
output gap remain in the dataset, so it is not a genuine vintage-data test.
"""

from __future__ import annotations

import matplotlib

matplotlib.use("Agg")
import matplotlib.pyplot as plt
import numpy as np
import pandas as pd
import statsmodels.api as sm
from statsmodels.tsa.api import VAR

from common import (
    DATA_CLEAN,
    FIGURES,
    TABLES,
    _state_coefficients,
    calibration_rho,
    complete_parameters,
    ensure_directories,
)
from hybrid_solution import hybrid_params, solve_hybrid, hybrid_forecast


VARIABLES = ["x", "pi", "i"]
MODELS = ["NK", "NK hibrido", "AR1", "Random walk", "VAR"]
HORIZONS = [1, 4]
MIN_TRAIN = 40
SCALES = {"x": 100.0, "pi": 400.0, "i": 400.0}
# The hybrid reduced form is calibration-fixed, so solve it once (gamma_pi=0.35,
# validated to reproduce Dynare's hybrid IRFs to machine precision).
HYBRID_SOLUTION = solve_hybrid(hybrid_params(0.35))


def ar1_forecast(series: np.ndarray, horizon: int) -> float:
    y = series[1:]
    x = sm.add_constant(series[:-1])
    result = sm.OLS(y, x).fit()
    value = float(series[-1])
    for _ in range(horizon):
        value = float(result.params[0] + result.params[1] * value)
    return value


def nk_forecast(train: pd.DataFrame, horizon: int) -> np.ndarray:
    params = complete_parameters({"rho_i": calibration_rho()})
    a_x, a_pi, a_i = _state_coefficients(params)
    means = train.mean()
    lagged_rate_gap = float(train["i"].iloc[-1] - means["i"])
    forecast = None
    for _ in range(horizon):
        forecast = np.array([a_x, a_pi, a_i]) * lagged_rate_gap
        lagged_rate_gap = float(forecast[2])
    return forecast + means[VARIABLES].to_numpy()


def nk_hybrid_forecast(train: pd.DataFrame, horizon: int) -> np.ndarray:
    means = train.mean()
    i_dev = float(train["i"].iloc[-1] - means["i"])
    pi_dev = float(train["pi"].iloc[-1] - means["pi"])
    forecast = hybrid_forecast(HYBRID_SOLUTION, i_dev, pi_dev, horizon)
    return forecast + means[VARIABLES].to_numpy()


def var_forecast(train: pd.DataFrame, horizon: int) -> np.ndarray:
    selection = VAR(train[VARIABLES]).select_order(maxlags=4, trend="c")
    lags = max(1, int(selection.selected_orders.get("bic") or 1))
    result = VAR(train[VARIABLES]).fit(lags, trend="c")
    return result.forecast(train[VARIABLES].to_numpy()[-lags:], steps=horizon)[-1]


def main() -> None:
    ensure_directories()
    macro = pd.read_csv(DATA_CLEAN / "chile_macro_quarterly.csv")
    data = pd.DataFrame(
        {
            "date": macro["date"],
            "x": macro["output_gap"],
            "pi": macro["infl_q"],
            "i": macro["i_q"],
        }
    )

    records: list[dict] = []
    for horizon in HORIZONS:
        for origin in range(MIN_TRAIN - 1, len(data) - horizon):
            train = data.iloc[: origin + 1]
            target = data.iloc[origin + horizon]
            predictions = {
                "NK": nk_forecast(train[VARIABLES], horizon),
                "NK hibrido": nk_hybrid_forecast(train[VARIABLES], horizon),
                "AR1": np.array(
                    [ar1_forecast(train[v].to_numpy(), horizon) for v in VARIABLES]
                ),
                "Random walk": train[VARIABLES].iloc[-1].to_numpy(),
            }
            try:
                predictions["VAR"] = var_forecast(train, horizon)
            except Exception:  # noqa: BLE001
                predictions["VAR"] = predictions["AR1"].copy()

            for model, prediction in predictions.items():
                for index, variable in enumerate(VARIABLES):
                    scale = SCALES[variable]
                    actual = float(target[variable] * scale)
                    forecast = float(prediction[index] * scale)
                    records.append(
                        {
                            "origin": train["date"].iloc[-1],
                            "target_date": target["date"],
                            "horizon": horizon,
                            "model": model,
                            "variable": variable,
                            "actual": actual,
                            "forecast": forecast,
                            "error": actual - forecast,
                            "absolute_error": abs(actual - forecast),
                            "squared_error": (actual - forecast) ** 2,
                        }
                    )

    forecasts = pd.DataFrame(records)
    forecasts.to_csv(TABLES / "forecast_oos_predictions.csv", index=False)
    metrics = (
        forecasts.groupby(["horizon", "variable", "model"], as_index=False)
        .agg(
            observations=("error", "size"),
            bias=("error", "mean"),
            mae=("absolute_error", "mean"),
            mse=("squared_error", "mean"),
        )
    )
    metrics["rmse"] = np.sqrt(metrics["mse"])
    random_walk_rmse = metrics[metrics["model"] == "Random walk"][
        ["horizon", "variable", "rmse"]
    ].rename(columns={"rmse": "random_walk_rmse"})
    metrics = metrics.merge(random_walk_rmse, on=["horizon", "variable"])
    metrics["relative_rmse_vs_random_walk"] = (
        metrics["rmse"] / metrics["random_walk_rmse"]
    )
    metrics.to_csv(TABLES / "forecast_oos_metrics.csv", index=False)

    fig, axes = plt.subplots(2, 3, figsize=(14, 8), sharey=False)
    colors = {
        "NK": "#B45309",
        "NK hibrido": "#B91C1C",
        "AR1": "#1E4E79",
        "Random walk": "#6B7280",
        "VAR": "#16A34A",
    }
    for row, horizon in enumerate(HORIZONS):
        for col, variable in enumerate(VARIABLES):
            ax = axes[row, col]
            sub = metrics[
                (metrics["horizon"] == horizon) & (metrics["variable"] == variable)
            ].set_index("model").reindex(MODELS)
            ax.bar(
                np.arange(len(MODELS)),
                sub["rmse"],
                color=[colors[m] for m in MODELS],
            )
            ax.set_xticks(np.arange(len(MODELS)))
            ax.set_xticklabels(MODELS, rotation=25, ha="right", fontsize=8)
            ax.set_title(f"{variable}: horizonte {horizon}")
            ax.set_ylabel("RMSE (p.p.)" if variable != "x" else "RMSE (% PIB)")
            ax.grid(alpha=0.2, axis="y")
    fig.suptitle("Avaliacao pseudo-fora-da-amostra: modelo NK versus benchmarks")
    fig.text(
        0.01,
        0.01,
        "Janela expansiva, treino inicial de 40 trimestres. Dados revisados e hiato HP "
        "de amostra completa: diagnostico pseudo-real-time, nao backtest em vintages.",
        fontsize=8,
    )
    fig.tight_layout(rect=(0, 0.04, 1, 0.95))
    fig.savefig(FIGURES / "forecast_oos_comparison.png", dpi=180)
    plt.close(fig)

    print(
        metrics[
            ["horizon", "variable", "model", "rmse", "mae",
             "relative_rmse_vs_random_walk"]
        ].to_string(index=False)
    )


if __name__ == "__main__":
    main()
\end{lstlisting}
\clearpage
\part*{Analise (SVAR, Bayes, resultados)}\addcontentsline{toc}{part}{Analise (SVAR, Bayes, resultados)}

\section{python/analyze\_svar.py}
\begin{lstlisting}[language=Python]
"""Estimate a recursive VAR and compare its monetary IRF with the NK model.

Identification uses the quarterly Cholesky ordering [x, pi, i]. Output and
inflation may affect the policy rate within the quarter, while a monetary shock
affects them with a lag. This is a transparent reduced-form robustness check,
not an incontrovertible identification of Chilean monetary shocks.
"""

from __future__ import annotations

import matplotlib

matplotlib.use("Agg")
import matplotlib.pyplot as plt
import numpy as np
import pandas as pd
from statsmodels.tsa.api import VAR

from common import DATA_CLEAN, DYNARE_OUTPUTS, FIGURES, TABLES, ensure_directories


VARIABLES = ["x", "pi", "i"]
LABELS = {"x": "Hiato do produto", "pi": "Inflacao", "i": "Taxa de politica"}
SCALES = {"x": 100.0, "pi": 400.0, "i": 400.0}
MONETARY_IMPACT_ANNUAL_PP = 0.25
# Dynare exports horizons 0..19 (20 observations).
HORIZON = 19


def main() -> None:
    ensure_directories()
    observables = pd.read_csv(DATA_CLEAN / "chile_observables.csv")
    data = observables[VARIABLES].astype(float)

    selection = VAR(data).select_order(maxlags=4, trend="c")
    selected_lag = max(1, int(selection.selected_orders.get("bic") or 1))
    lag_candidates = []
    for lag in range(1, 5):
        candidate = VAR(data).fit(lag, trend="c")
        try:
            whiteness = float(candidate.test_whiteness(nlags=8).pvalue)
        except Exception:  # noqa: BLE001
            whiteness = np.nan
        lag_candidates.append((lag, candidate, whiteness))
    acceptable = [item for item in lag_candidates if item[2] >= 0.05]
    if acceptable:
        estimated_lag, result, selected_whiteness = acceptable[0]
        lag_reason = "smallest lag with residual-whiteness p>=0.05"
    else:
        estimated_lag = selected_lag
        result = VAR(data).fit(estimated_lag, trend="c")
        selected_whiteness = float(
            result.test_whiteness(nlags=max(8, estimated_lag + 2)).pvalue
        )
        lag_reason = "BIC; no lag up to 4 passed residual-whiteness test"

    orth = result.orth_ma_rep(HORIZON)
    lower, upper = result.irf_errband_mc(
        orth=True, repl=2000, steps=HORIZON, signif=0.10, seed=20260615
    )
    policy_index = VARIABLES.index("i")
    impact_i = float(orth[0, policy_index, policy_index])
    normalization = MONETARY_IMPACT_ANNUAL_PP / (SCALES["i"] * impact_i)

    baseline = pd.read_csv(DYNARE_OUTPUTS / "baseline" / "irfs.csv")
    dsge_impact = float(baseline.loc[0, "i_e_i"])
    dsge_normalization = MONETARY_IMPACT_ANNUAL_PP / (
        SCALES["i"] * dsge_impact
    )

    records: list[dict] = []
    fig, axes = plt.subplots(1, 3, figsize=(14, 4.4), sharex=True)
    horizons = np.arange(HORIZON + 1)
    for response_index, variable in enumerate(VARIABLES):
        scale = SCALES[variable]
        point = orth[:, response_index, policy_index] * normalization * scale
        lo = lower[:, response_index, policy_index] * normalization * scale
        hi = upper[:, response_index, policy_index] * normalization * scale
        dsge = (
            baseline[f"{variable}_e_i"].to_numpy()[: HORIZON + 1]
            * dsge_normalization
            * scale
        )
        ax = axes[response_index]
        ax.fill_between(horizons, lo, hi, color="#9DB9D1", alpha=0.5,
                        label="SVAR: banda 90%")
        ax.plot(horizons, point, color="#1E4E79", linewidth=2.2,
                label=f"SVAR recursivo, VAR({estimated_lag})")
        ax.plot(horizons, dsge, color="#B45309", linestyle="--", linewidth=2.2,
                label="DSGE/NK")
        ax.axhline(0, color="black", linewidth=0.7)
        ax.set_title(LABELS[variable])
        ax.set_xlabel("Trimestres")
        ax.set_ylabel("p.p." if variable != "x" else "% do PIB")
        ax.grid(alpha=0.25)
        for h in horizons:
            records.append(
                {
                    "horizon": int(h),
                    "variable": variable,
                    "svar_response": float(point[h]),
                    "svar_lower_90": float(lo[h]),
                    "svar_upper_90": float(hi[h]),
                    "dsge_response": float(dsge[h]),
                    "normalization": "policy-rate impact = +0.25 annual p.p.",
                    "ordering": "x,pi,i",
                }
            )

    handles, labels = axes[0].get_legend_handles_labels()
    fig.legend(handles, labels, ncol=3, loc="lower center", frameon=False)
    fig.suptitle(
        "Choque monetario contracionista: SVAR recursivo versus modelo NK"
    )
    fig.text(
        0.01,
        0.01,
        f"SVAR com ordenacao [x, pi, i], VAR({estimated_lag}); bandas Monte Carlo de 90%. "
        "Ambas as IRFs normalizadas para elevar a taxa em 0,25 p.p. anual no impacto.",
        fontsize=8,
    )
    fig.tight_layout(rect=(0, 0.10, 1, 0.94))
    fig.savefig(FIGURES / "svar_vs_dsge_irf.png", dpi=180)
    plt.close(fig)

    try:
        whiteness_p = float(result.test_whiteness(nlags=max(8, estimated_lag + 2)).pvalue)
    except Exception:  # noqa: BLE001
        whiteness_p = np.nan
    try:
        normality_p = float(result.test_normality().pvalue)
    except Exception:  # noqa: BLE001
        normality_p = np.nan

    diagnostics = pd.DataFrame(
        [
            {
                "observations": int(result.nobs),
                "selected_lags_bic": selected_lag,
                "estimated_lags": estimated_lag,
                "lag_choice_reason": lag_reason,
                "aic": float(result.aic),
                "bic": float(result.bic),
                "hqic": float(result.hqic),
                "stable": bool(result.is_stable()),
                "max_inverse_root": float(np.max(np.abs(result.roots) ** -1)),
                "whiteness_test_pvalue": whiteness_p,
                "normality_test_pvalue": normality_p,
                "cholesky_ordering": "x,pi,i",
                "monetary_shock": "third orthogonal innovation",
                "impact_normalization_annual_pp": MONETARY_IMPACT_ANNUAL_PP,
            }
        ]
    )
    pd.DataFrame(records).to_csv(TABLES / "svar_vs_dsge_irfs.csv", index=False)
    diagnostics.to_csv(TABLES / "svar_diagnostics.csv", index=False)

    print(diagnostics.to_string(index=False))
    for variable in VARIABLES:
        frame = pd.DataFrame(records)
        row = frame[(frame["variable"] == variable) & (frame["horizon"] == 4)].iloc[0]
        print(
            f"h=4 {variable}: SVAR={row['svar_response']:.4f}, "
            f"DSGE={row['dsge_response']:.4f}"
        )


if __name__ == "__main__":
    main()
\end{lstlisting}
\clearpage
\section{python/compare\_bayesian\_models.py}
\begin{lstlisting}[language=Python]
"""Compare the forward-looking and hybrid NK models by marginal likelihood.

Both models use the same Chilean observables and common priors. Dynare finds
each posterior mode and Hessian. The script then applies the Laplace
approximation:

    log p(y|M) ~= log p(y, theta_hat|M)
                  + k/2 log(2*pi) - 1/2 log|H(theta_hat)|.

This is a formal marginal-data-density comparison, but an approximation around
one mode. It complements, rather than replaces, the full baseline MCMC.
"""

from __future__ import annotations

import argparse
import os
import shutil
import subprocess
import sys
import tempfile
from pathlib import Path

import matplotlib.pyplot as plt
import numpy as np
import pandas as pd
from scipy.io import loadmat

from common import DATA_CLEAN, FIGURES, LOGS, ROOT, TABLES, find_dynare_matlab, find_octave


MODELS = {
    "forward_nkpc": "nk_chile_estim.mod",
    "hybrid_nkpc": "nk_chile_hybrid_estim.mod",
}


def octave_literal(path: Path) -> str:
    return str(path).replace("\\", "/").replace("'", "''")


def kill_tree(proc: subprocess.Popen) -> None:
    subprocess.run(
        ["taskkill", "/F", "/T", "/PID", str(proc.pid)],
        capture_output=True,
        check=False,
    )


def run_mode(
    model_id: str,
    filename: str,
    work: Path,
    octave: Path,
    dynare: Path,
    timeout: int,
) -> tuple[Path, str]:
    model_work = work / model_id
    model_work.mkdir(parents=True, exist_ok=True)
    shutil.copy(ROOT / "dynare" / filename, model_work / filename)
    pd.read_csv(DATA_CLEAN / "chile_observables.csv")[["pi", "i", "x"]].to_csv(
        model_work / "chile_observables_dynare.csv", index=False
    )
    expression = (
        f"addpath('{octave_literal(dynare)}'); "
        f"dynare('{filename}','noclearall','nolog');"
    )
    proc = subprocess.Popen(
        [str(octave), "--quiet", "--eval", expression],
        cwd=model_work,
        env=os.environ.copy(),
        text=True,
        stdout=subprocess.PIPE,
        stderr=subprocess.PIPE,
    )
    try:
        stdout, stderr = proc.communicate(timeout=timeout)
    except subprocess.TimeoutExpired:
        kill_tree(proc)
        stdout, stderr = proc.communicate()
        raise TimeoutError(f"{model_id} posterior mode exceeded {timeout}s")
    log = f"RC={proc.returncode}\nSTDOUT\n{stdout}\nSTDERR\n{stderr}\n"
    mode_files = list(model_work.rglob("*_mode.mat"))
    if proc.returncode != 0 or not mode_files:
        raise RuntimeError(
            f"{model_id} mode failed (rc={proc.returncode}); "
            f"tail={log[-1800:]}"
        )
    return mode_files[0], log


def laplace_summary(model_id: str, path: Path) -> dict[str, object]:
    payload = loadmat(path)
    mode = np.asarray(payload["xparam1"], dtype=float).ravel()
    hessian = np.asarray(payload["hh"], dtype=float)
    objective = float(np.asarray(payload["fval"]).ravel()[0])
    sign, logdet = np.linalg.slogdet(hessian)
    eigenvalues = np.linalg.eigvalsh(hessian)
    if sign <= 0 or eigenvalues.min() <= 0:
        raise ValueError(f"{model_id}: mode Hessian is not positive definite")
    parameter_names = [
        str(item[0]) if isinstance(item, np.ndarray) else str(item)
        for item in np.asarray(payload["parameter_names"]).ravel()
    ]
    k = len(mode)
    log_joint_mode = -objective
    log_mdd = log_joint_mode + 0.5 * k * np.log(2 * np.pi) - 0.5 * logdet
    gamma = np.nan
    if "gamma_pi" in parameter_names:
        gamma = float(mode[parameter_names.index("gamma_pi")])
    return {
        "model": model_id,
        "parameters": k,
        "log_joint_at_mode": log_joint_mode,
        "hessian_log_determinant": logdet,
        "laplace_log_marginal_likelihood": log_mdd,
        "gamma_pi_mode": gamma,
        "hessian_min_eigenvalue": float(eigenvalues.min()),
        "method": "Laplace approximation at Dynare posterior mode",
        "same_observables": "pi,i,x",
        "observations": len(pd.read_csv(DATA_CLEAN / "chile_observables.csv")),
    }


def make_figure(frame: pd.DataFrame) -> None:
    best = frame["laplace_log_marginal_likelihood"].max()
    relative = frame["laplace_log_marginal_likelihood"] - best
    labels = frame["model"].map(
        {"forward_nkpc": "NKPC prospectiva", "hybrid_nkpc": "NKPC hibrida"}
    )
    colors = ["#7A9E9F" if value < 0 else "#145DA0" for value in relative]
    fig, ax = plt.subplots(figsize=(8.5, 5.2))
    bars = ax.bar(labels, relative, color=colors, width=0.58)
    ax.axhline(0, color="#333333", linewidth=0.9)
    ax.set_ylabel("Log evidencia relativa ao melhor modelo")
    ax.set_title("Comparacao bayesiana de modelos: evidencia marginal (Laplace)")
    ax.grid(axis="y", alpha=0.25)
    for bar, value in zip(bars, relative, strict=True):
        ax.text(
            bar.get_x() + bar.get_width() / 2,
            value + (0.25 if value >= 0 else -0.25),
            f"{value:.2f}",
            ha="center",
            va="bottom" if value >= 0 else "top",
            fontweight="bold",
        )
    fig.text(
        0.5,
        0.01,
        "Mesma amostra e mesmos observaveis. A evidencia penaliza o parametro extra "
        "da Phillips hibrida; aproximacao local, nao estimativa MCMC da evidencia.",
        ha="center",
        fontsize=9,
    )
    fig.tight_layout(rect=(0, 0.06, 1, 1))
    fig.savefig(FIGURES / "bayesian_model_comparison.png", dpi=180, bbox_inches="tight")
    plt.close(fig)


def main() -> int:
    parser = argparse.ArgumentParser()
    parser.add_argument("--timeout", type=int, default=420)
    parser.add_argument(
        "--reuse-baseline",
        action="store_true",
        help="Reuse the existing baseline mode in the system temp directory.",
    )
    args = parser.parse_args()
    octave = find_octave()
    dynare = find_dynare_matlab()
    if octave is None or dynare is None:
        print("ERROR: Octave/Dynare not found.", file=sys.stderr)
        return 1

    work = Path(tempfile.gettempdir()) / "nk_model_comparison_work"
    if work.exists():
        shutil.rmtree(work, ignore_errors=True)
    work.mkdir(parents=True)
    logs: list[str] = []
    summaries: list[dict[str, object]] = []

    for model_id, filename in MODELS.items():
        if model_id == "forward_nkpc" and args.reuse_baseline:
            cached = (
                Path(tempfile.gettempdir())
                / "nk_estim_work"
                / "nk_chile_estim"
                / "Output"
                / "nk_chile_estim_mode.mat"
            )
            if cached.exists():
                mode_file = cached
                log = f"Reused baseline mode: {cached}\n"
            else:
                mode_file, log = run_mode(
                    model_id, filename, work, octave, dynare, args.timeout
                )
        else:
            mode_file, log = run_mode(
                model_id, filename, work, octave, dynare, args.timeout
            )
        logs.append(f"\n===== {model_id} =====\n{log}")
        summaries.append(laplace_summary(model_id, mode_file))

    frame = pd.DataFrame(summaries)
    best = float(frame["laplace_log_marginal_likelihood"].max())
    frame["log_bayes_factor_vs_best"] = (
        frame["laplace_log_marginal_likelihood"] - best
    )
    weights = np.exp(frame["log_bayes_factor_vs_best"])
    frame["posterior_model_probability_equal_prior"] = weights / weights.sum()
    preferred = frame.loc[
        frame["laplace_log_marginal_likelihood"].idxmax(), "model"
    ]
    frame["preferred_model"] = preferred
    frame.to_csv(TABLES / "bayesian_model_comparison.csv", index=False)
    make_figure(frame)
    (LOGS / "bayesian_model_comparison.log").write_text(
        "".join(logs), encoding="utf-8"
    )
    print(frame.to_string(index=False))
    print(f"Preferred model under equal prior odds: {preferred}")
    return 0


if __name__ == "__main__":
    raise SystemExit(main())
\end{lstlisting}
\clearpage
\section{python/analyze\_model\_results.py}
\begin{lstlisting}[language=Python]
"""Build diagnostic tables that connect model outputs to the assignment.

The script compares theoretical and empirical moments, summarizes IRF impact and
convergence metrics, and checks the expected impact signs of the three structural
shocks. It also creates a model-versus-data moments figure.
"""

from __future__ import annotations

import json
import math

import matplotlib

matplotlib.use("Agg")
import matplotlib.pyplot as plt
import numpy as np
import pandas as pd

from common import (
    DATA_CLEAN,
    DYNARE_OUTPUTS,
    FIGURES,
    SHOCK_STD,
    TABLES,
    _state_coefficients,
    calibration_rho,
    complete_parameters,
    ensure_directories,
)

VARIABLES = ("x", "pi", "i")
LABELS = {"x": "Output gap", "pi": "Inflation", "i": "Policy rate"}


def acf1(values: pd.Series) -> float:
    clean = pd.to_numeric(values, errors="coerce").dropna()
    return float(clean.autocorr(lag=1)) if len(clean) > 2 else math.nan


def empirical_moment_rows(panel: pd.DataFrame, sample: str) -> list[dict]:
    series = {
        "x": panel["output_gap"],
        "pi": panel["infl_q"],
        "i": panel["i_q"],
    }
    rows = []
    for variable, values in series.items():
        clean = pd.to_numeric(values, errors="coerce").dropna()
        rows.append(
            {
                "source": "BCCh data",
                "sample": sample,
                "variable": variable,
                "mean_pp_quarterly": 100.0 * float(clean.mean()),
                "std_dev_pp_quarterly": 100.0 * float(clean.std(ddof=1)),
                "variance": float(clean.var(ddof=1)),
                "autocorrelation_1": acf1(clean),
                "observations": int(len(clean)),
            }
        )
    return rows


def model_moment_rows(moments: pd.DataFrame, scenario: str, sample: str) -> list[dict]:
    selected = moments[moments["scenario"] == scenario]
    rows = []
    for _, row in selected.iterrows():
        variable = str(row["variable"])
        model_mean = float(row["mean"])
        if variable == "i":
            # Dynare reports the level r*; comparisons focus on cyclical deviations.
            model_mean = 0.0
        rows.append(
            {
                "source": "Dynare model",
                "sample": sample,
                "variable": variable,
                "mean_pp_quarterly": 100.0 * model_mean,
                "std_dev_pp_quarterly": 100.0 * float(row["std_dev"]),
                "variance": float(row["variance"]),
                "autocorrelation_1": float(row["autocorrelation_1"]),
                "observations": np.nan,
            }
        )
    return rows


def build_moment_comparison() -> pd.DataFrame:
    panel = pd.read_csv(DATA_CLEAN / "chile_macro_quarterly.csv")
    moments = pd.read_csv(TABLES / "moments_all_scenarios.csv")
    dates = panel["date"].astype(str)
    pandemic = dates.between("2020Q2", "2021Q4")

    rows = empirical_moment_rows(panel, "2001Q1-2026Q1")
    rows += empirical_moment_rows(
        panel.loc[~pandemic].reset_index(drop=True),
        "excluding 2020Q2-2021Q4",
    )
    rows += model_moment_rows(moments, "baseline", "rho_i estimated (0.934)")
    rows += model_moment_rows(
        moments,
        "rho_calibrated_0p80",
        "rho_i calibrated (0.80)",
    )
    table = pd.DataFrame(rows)
    table.to_csv(TABLES / "moments_model_vs_data.csv", index=False)
    return table


def plot_moments(table: pd.DataFrame) -> None:
    order = [
        "2001Q1-2026Q1",
        "excluding 2020Q2-2021Q4",
        "rho_i estimated (0.934)",
        "rho_i calibrated (0.80)",
    ]
    colors = ["#111827", "#6B7280", "#1E4E79", "#FF6B35"]
    fig, axes = plt.subplots(1, 2, figsize=(13, 5.2))
    width = 0.19
    positions = np.arange(len(VARIABLES))
    for index, (sample, color) in enumerate(zip(order, colors)):
        selected = table[table["sample"] == sample].set_index("variable")
        offset = (index - 1.5) * width
        axes[0].bar(
            positions + offset,
            selected.loc[list(VARIABLES), "std_dev_pp_quarterly"],
            width=width,
            label=sample,
            color=color,
        )
        axes[1].bar(
            positions + offset,
            selected.loc[list(VARIABLES), "autocorrelation_1"],
            width=width,
            label=sample,
            color=color,
        )
    axes[0].set_title("Quarterly standard deviations")
    axes[0].set_ylabel("Percentage points")
    axes[1].set_title("First-order autocorrelation")
    axes[1].axhline(0, color="black", linewidth=0.7)
    for ax in axes:
        ax.set_xticks(positions, [LABELS[var] for var in VARIABLES])
        ax.grid(axis="y", alpha=0.25)
    axes[0].legend(fontsize=8)
    fig.suptitle("Chile: theoretical moments versus empirical moments")
    fig.tight_layout(rect=(0, 0, 1, 0.95))
    fig.savefig(FIGURES / "moments_model_vs_data.png", dpi=180)
    plt.close(fig)


def first_sign_change(frame: pd.DataFrame) -> float:
    ordered = frame.sort_values("horizon")
    initial = float(ordered.iloc[0]["response"])
    if abs(initial) < 1e-14:
        return math.nan
    initial_sign = np.sign(initial)
    changed = ordered[
        (np.sign(ordered["response"]) != initial_sign)
        & (ordered["response"].abs() > 1e-12)
    ]
    return float(changed.iloc[0]["horizon"]) if not changed.empty else math.nan


def build_irf_metrics() -> pd.DataFrame:
    irfs = pd.read_csv(TABLES / "irfs_long.csv")
    manifest = pd.read_csv(TABLES / "scenario_manifest.csv")
    types = dict(zip(manifest["scenario"], manifest["scenario_type"]))
    types["baseline"] = "baseline"
    rows = []
    for (scenario, variable, shock), frame in irfs.groupby(
        ["scenario", "variable", "shock"]
    ):
        ordered = frame.sort_values("horizon")
        peak_row = ordered.loc[ordered["response"].abs().idxmax()]
        by_horizon = ordered.set_index("horizon")["response"]
        rows.append(
            {
                "scenario": scenario,
                "scenario_type": types.get(scenario, ""),
                "variable": variable,
                "shock": shock,
                "impact_pp": 100.0 * float(ordered.iloc[0]["response"]),
                "peak_response_pp": 100.0 * float(peak_row["response"]),
                "peak_abs_pp": 100.0 * abs(float(peak_row["response"])),
                "peak_horizon": int(peak_row["horizon"]),
                "cumulative_abs_pp": 100.0 * float(ordered["response"].abs().sum()),
                "first_sign_change_horizon": first_sign_change(ordered),
                "response_h4_pp": 100.0 * float(by_horizon.get(4, np.nan)),
                "response_h8_pp": 100.0 * float(by_horizon.get(8, np.nan)),
                "response_h19_pp": 100.0 * float(by_horizon.get(19, np.nan)),
                "source": str(ordered.iloc[0]["source"]),
            }
        )
    table = pd.DataFrame(rows)
    table.to_csv(TABLES / "irf_summary_metrics.csv", index=False)
    return table


def build_sign_checks(metrics: pd.DataFrame) -> pd.DataFrame:
    expected = {
        ("e_x", "x"): 1,
        ("e_x", "pi"): 1,
        ("e_x", "i"): 1,
        ("e_pi", "x"): -1,
        ("e_pi", "pi"): 1,
        ("e_pi", "i"): 1,
        ("e_i", "x"): -1,
        ("e_i", "pi"): -1,
        ("e_i", "i"): 1,
    }
    baseline = metrics[metrics["scenario"] == "baseline"].copy()
    baseline["expected_impact_sign"] = baseline.apply(
        lambda row: expected[(row["shock"], row["variable"])], axis=1
    )
    baseline["observed_impact_sign"] = np.sign(baseline["impact_pp"]).astype(int)
    baseline["impact_sign_matches"] = (
        baseline["expected_impact_sign"] == baseline["observed_impact_sign"]
    )
    output = baseline[
        [
            "shock",
            "variable",
            "expected_impact_sign",
            "observed_impact_sign",
            "impact_sign_matches",
            "impact_pp",
            "first_sign_change_horizon",
        ]
    ]
    output.to_csv(TABLES / "irf_sign_checks.csv", index=False)
    return output


def build_policy_functions() -> pd.DataFrame:
    params = complete_parameters({"rho_i": calibration_rho()})
    a_x, a_pi, a_i = _state_coefficients(params)
    sigma, beta, kappa = params["sigma"], params["beta"], params["kappa"]
    rho_i, phi_pi, phi_x = params["rho_i"], params["phi_pi"], params["phi_x"]
    c_i = a_x - (1.0 - a_pi) / sigma
    structural = np.array(
        [
            [1.0, 0.0, -c_i],
            [-kappa, 1.0, -beta * a_pi],
            [-(1.0 - rho_i) * phi_x, -(1.0 - rho_i) * phi_pi, 1.0],
        ]
    )
    impact = np.linalg.inv(structural)
    transition = np.array([a_x, a_pi, a_i])
    drivers = [
        ("constant", np.array([0.0, 0.0, params["rstar"]])),
        ("i(-1)", transition),
        ("e_x", impact[:, 0]),
        ("e_pi", impact[:, 1]),
        ("e_i", impact[:, 2]),
    ]
    rows = []
    for driver, coefficients in drivers:
        for variable, coefficient in zip(VARIABLES, coefficients):
            rows.append(
                {
                    "driver": driver,
                    "variable": variable,
                    "coefficient_unit_shock": float(coefficient),
                    "coefficient_calibrated_shock": (
                        float(coefficient) * SHOCK_STD[driver]
                        if driver in SHOCK_STD
                        else float(coefficient)
                    ),
                }
            )
    table = pd.DataFrame(rows)
    table.to_csv(TABLES / "policy_transition_coefficients.csv", index=False)
    return table


def build_correlation_comparison() -> pd.DataFrame:
    panel = pd.read_csv(DATA_CLEAN / "chile_macro_quarterly.csv")
    empirical = pd.DataFrame(
        {
            "x": panel["output_gap"],
            "pi": panel["infl_q"],
            "i": panel["i_q"],
        }
    ).corr()
    rows = []
    for variable in VARIABLES:
        for other in VARIABLES:
            rows.append(
                {
                    "source": "BCCh data",
                    "variable": variable,
                    "other_variable": other,
                    "correlation": float(empirical.loc[variable, other]),
                }
            )
    model_path = DYNARE_OUTPUTS / "baseline" / "correlations.csv"
    if model_path.exists():
        model = pd.read_csv(model_path).set_index("variable")
        for variable in VARIABLES:
            for other in VARIABLES:
                rows.append(
                    {
                        "source": "Dynare baseline",
                        "variable": variable,
                        "other_variable": other,
                        "correlation": float(model.loc[variable, other]),
                    }
                )
    table = pd.DataFrame(rows)
    table.to_csv(TABLES / "correlations_model_vs_data.csv", index=False)
    return table


def main() -> None:
    ensure_directories()
    metadata = json.loads(
        (DATA_CLEAN / "dataset_metadata.json").read_text(encoding="utf-8")
    )
    if bool(metadata.get("is_synthetic", True)):
        print("WARNING: empirical moment rows use the labelled synthetic fallback.")
    moment_table = build_moment_comparison()
    plot_moments(moment_table)
    metrics = build_irf_metrics()
    signs = build_sign_checks(metrics)
    policy = build_policy_functions()
    correlations = build_correlation_comparison()
    print(f"Moment comparison rows: {len(moment_table)}")
    print(f"IRF metric rows: {len(metrics)}")
    print(f"Policy-function rows: {len(policy)}")
    print(f"Correlation rows: {len(correlations)}")
    print(
        "Baseline impact sign checks: "
        f"{int(signs['impact_sign_matches'].sum())}/{len(signs)}"
    )


if __name__ == "__main__":
    main()
\end{lstlisting}
\clearpage
\section{python/analyze\_macro\_extensions.py}
\begin{lstlisting}[language=Python]
"""Plot and summarize the open-economy and hybrid-Phillips extensions."""

from __future__ import annotations

import matplotlib

matplotlib.use("Agg")
import matplotlib.pyplot as plt
import numpy as np
import pandas as pd

from common import DYNARE_OUTPUTS, FIGURES, TABLES, ensure_directories


def load_irfs(scenario: str) -> pd.DataFrame:
    return pd.read_csv(DYNARE_OUTPUTS / scenario / "irfs.csv")


def main() -> None:
    ensure_directories()
    open_irf = load_irfs("open_economy")
    hybrid_irf = load_irfs("hybrid_nkpc")
    baseline = load_irfs("baseline")
    h = open_irf["horizon"].to_numpy()

    variables = ["x", "pi", "i", "q", "copper"]
    shocks = ["e_i", "e_q", "e_copper"]
    labels = {
        "x": "Hiato",
        "pi": "Inflacao",
        "i": "Juro",
        "q": "Cambio real/nominal gap",
        "copper": "Cobre",
        "e_i": "Choque monetario",
        "e_q": "Choque cambial/risco",
        "e_copper": "Choque de cobre",
    }
    scales = {"x": 100, "pi": 400, "i": 400, "q": 100, "copper": 100}
    fig, axes = plt.subplots(5, 3, figsize=(14, 14), sharex=True)
    for row, variable in enumerate(variables):
        for col, shock in enumerate(shocks):
            ax = axes[row, col]
            ax.plot(h, open_irf[f"{variable}_{shock}"] * scales[variable],
                    color="#1E4E79", linewidth=2)
            ax.axhline(0, color="black", linewidth=0.6)
            if row == 0:
                ax.set_title(labels[shock])
            if col == 0:
                ax.set_ylabel(labels[variable])
            ax.grid(alpha=0.2)
    fig.suptitle("Modelo NK aberto: transmissao monetaria, cambial e do cobre")
    fig.text(
        0.01,
        0.01,
        "Extensao calibrada ilustrativa. q>0 representa depreciacao; respostas de x/q/cobre "
        "em %, pi/i em p.p. anualizados.",
        fontsize=8,
    )
    fig.tight_layout(rect=(0, 0.03, 1, 0.96))
    fig.savefig(FIGURES / "open_economy_irfs.png", dpi=180)
    plt.close(fig)

    fig, axes = plt.subplots(2, 3, figsize=(14, 7), sharex=True)
    for col, shock in enumerate(["e_x", "e_pi", "e_i"]):
        for row, variable in enumerate(["pi", "x"]):
            ax = axes[row, col]
            scale = 400 if variable == "pi" else 100
            ax.plot(
                baseline["horizon"],
                baseline[f"{variable}_{shock}"] * scale,
                color="#6B7280",
                linewidth=2,
                label="NK prospectivo",
            )
            ax.plot(
                hybrid_irf["horizon"],
                hybrid_irf[f"{variable}_{shock}"] * scale,
                color="#B45309",
                linestyle="--",
                linewidth=2,
                label="NKPC hibrida",
            )
            ax.axhline(0, color="black", linewidth=0.6)
            ax.set_title(f"{variable} a {shock}")
            ax.grid(alpha=0.2)
    axes[0, 0].legend(fontsize=8)
    fig.suptitle("Persistencia intrinseca: Phillips prospectiva versus hibrida")
    fig.tight_layout(rect=(0, 0, 1, 0.95))
    fig.savefig(FIGURES / "hybrid_nkpc_irfs.png", dpi=180)
    plt.close(fig)

    fig, axes = plt.subplots(1, 3, figsize=(13, 4.2), sharex=True)
    for ax, variable in zip(axes, ["x", "pi", "i"]):
        scale = 100 if variable == "x" else 400
        ax.plot(
            baseline["horizon"],
            baseline[f"{variable}_e_i"] * scale,
            color="#6B7280",
            linewidth=2,
            label="Fechado",
        )
        ax.plot(
            open_irf["horizon"],
            open_irf[f"{variable}_e_i"] * scale,
            color="#1E4E79",
            linewidth=2,
            linestyle="--",
            label="Aberto",
        )
        ax.axhline(0, color="black", linewidth=0.6)
        ax.set_title(labels[variable])
        ax.grid(alpha=0.2)
    axes[0].legend()
    fig.suptitle("Choque monetario: modelo fechado versus economia aberta")
    fig.tight_layout(rect=(0, 0, 1, 0.93))
    fig.savefig(FIGURES / "open_vs_closed_monetary_irf.png", dpi=180)
    plt.close(fig)

    rows = []
    for scenario, frame in [
        ("baseline", baseline),
        ("hybrid_nkpc", hybrid_irf),
        ("open_economy", open_irf),
    ]:
        available_shocks = ["e_x", "e_pi", "e_i"]
        if scenario == "open_economy":
            available_shocks += ["e_q", "e_copper"]
        for variable in ["x", "pi", "i"]:
            scale = 100 if variable == "x" else 400
            for shock in available_shocks:
                values = frame[f"{variable}_{shock}"].to_numpy() * scale
                rows.append(
                    {
                        "scenario": scenario,
                        "variable": variable,
                        "shock": shock,
                        "impact": float(values[0]),
                        "peak_abs": float(np.max(np.abs(values))),
                        "cumulative_abs_20q": float(np.sum(np.abs(values))),
                        "response_h4": float(values[4]),
                    }
                )
    pd.DataFrame(rows).to_csv(TABLES / "macro_extension_irf_metrics.csv", index=False)
    print("Wrote macro-extension figures and metrics.")


if __name__ == "__main__":
    main()
\end{lstlisting}
\clearpage
\part*{Nucleos da inflacao e mercado de trabalho}\addcontentsline{toc}{part}{Nucleos da inflacao e mercado de trabalho}

\section{python/build\_inflation\_cores.py}
\begin{lstlisting}[language=Python]
"""Decompose Chilean CPI (IPC) inflation by core/sub-index, from public FRED data.

FRED OECD series (no credentials). The component indices are only published to
~2023 (energy/services/core) so the analysis covers the build-up and the 2021-23
surge but not the most recent disinflation -- stated honestly on the figure.

Produces:
  * data/clean/chile_inflation_cores.csv -- YoY % inflation by component;
  * ipc_cores_panel.png   -- one panel per core vs the 3% target and 2-4% band;
  * ipc_cores_drivers.png -- what drove the surge: components through 2018-2023 +
    the peak-by-component ranking (what was most impactful).
"""

from __future__ import annotations

import io
import urllib.request

import matplotlib

matplotlib.use("Agg")
import matplotlib.pyplot as plt
import numpy as np
import pandas as pd

from common import DATA_CLEAN, FIGURES, ensure_directories

FRED_CSV = "https://fred.stlouisfed.org/graph/fredgraph.csv?id={series_id}"
SERIES = {
    "Cheia (todos os itens)": "CHLCPIALLMINMEI",
    "Nucleo (s/ alim. e energia)": "CHLCPICORMINMEI",
    "Energia": "CHLCPIENGMINMEI",
    "Servicos": "CHLCPGRSE01IXOBM",
}
COLORS = {
    "Cheia (todos os itens)": "#111827",
    "Nucleo (s/ alim. e energia)": "#1E4E79",
    "Energia": "#B45309",
    "Servicos": "#16A34A",
}
TARGET, BAND = 3.0, (2.0, 4.0)


def download_fred(series_id: str) -> pd.Series:
    url = FRED_CSV.format(series_id=series_id)
    with urllib.request.urlopen(url, timeout=60) as response:
        raw = response.read().decode("utf-8", "replace")
    frame = pd.read_csv(io.StringIO(raw))
    frame.columns = ["date", "value"]
    frame["date"] = pd.to_datetime(frame["date"], errors="coerce")
    frame["value"] = pd.to_numeric(frame["value"], errors="coerce")
    frame = frame.dropna()
    return pd.Series(frame["value"].to_numpy(), index=pd.DatetimeIndex(frame["date"]))


def main() -> None:
    ensure_directories()
    yoy = {}
    for name, sid in SERIES.items():
        index = download_fred(sid).sort_index()
        # 12-month change of the monthly index = annualised (YoY) inflation rate.
        yoy[name] = (index / index.shift(12) - 1.0) * 100.0
    data = pd.DataFrame(yoy).dropna(how="all")
    data.index.name = "date"
    data.to_csv(DATA_CLEAN / "chile_inflation_cores.csv")
    last_full = data["Cheia (todos os itens)"].dropna().index.max()

    # ------------------------------------------------------------------ #
    # Figure 1: one panel per core vs the target band.                    #
    # ------------------------------------------------------------------ #
    fig, axes = plt.subplots(2, 2, figsize=(12.5, 7.2), sharex=True)
    for ax, name in zip(axes.ravel(), SERIES):
        series = data[name].dropna().loc["2000-01-01":]  # inflation-targeting era only
        ax.axhspan(BAND[0], BAND[1], color="#FDE68A", alpha=0.35, label="Banda 2--4%")
        ax.axhline(TARGET, color="#B45309", linestyle="--", linewidth=1, label="Meta 3%")
        ax.plot(series.index, series.values, color=COLORS[name], linewidth=1.8)
        latest = series.iloc[-1]
        inside = BAND[0] <= latest <= BAND[1]
        ax.set_title(f"{name}  --  ultimo {latest:.1f}% "
                     f"({'dentro' if inside else 'fora'} da banda)", fontsize=10)
        ax.set_ylabel("% a.a. (YoY)")
        ax.grid(alpha=0.25)
    axes[0, 0].legend(fontsize=8, loc="upper left")
    fig.suptitle("Inflacao do IPC chileno por nucleo (variacao anual) vs meta de 3%")
    fig.text(0.01, 0.01, f"Fonte: OECD via FRED. Componentes ate {last_full:%Y-%m} "
             "(OECD descontinuou; perde a desinflacao recente). Variacao interanual = anualizada.",
             fontsize=7.5)
    fig.tight_layout(rect=(0, 0.03, 1, 0.96))
    fig.savefig(FIGURES / "ipc_cores_panel.png", dpi=170)
    plt.close(fig)

    # ------------------------------------------------------------------ #
    # Figure 2: what drove the surge (2018-2023) + peak ranking.          #
    # ------------------------------------------------------------------ #
    window = data.loc["2018-01-01":]
    fig, (axL, axR) = plt.subplots(1, 2, figsize=(13.5, 5.0),
                                   gridspec_kw={"width_ratios": [2.0, 1.0]})
    for name in SERIES:
        s = window[name].dropna()
        axL.plot(s.index, s.values, color=COLORS[name], linewidth=2.0, label=name)
    axL.axhspan(BAND[0], BAND[1], color="#FDE68A", alpha=0.3)
    axL.axhline(TARGET, color="#B45309", linestyle="--", linewidth=1)
    axL.set_title("O surto de 2021--2023 por nucleo")
    axL.set_ylabel("Inflacao anual (%)")
    axL.grid(alpha=0.25)
    axL.legend(fontsize=8)

    # Peak YoY during 2021-2023 by component = "what was most impactful".
    surge = data.loc["2021-01-01":"2023-12-31"]
    peaks = {name: float(surge[name].max()) for name in SERIES}
    order = sorted(peaks, key=peaks.get)
    axR.barh(range(len(order)), [peaks[n] for n in order],
             color=[COLORS[n] for n in order])
    axR.set_yticks(range(len(order)))
    axR.set_yticklabels([n.split(" (")[0] for n in order], fontsize=8)
    axR.axvline(TARGET, color="#B45309", linestyle="--", linewidth=1)
    axR.set_title("Pico de inflacao 2021--23")
    axR.set_xlabel("% a.a.")
    axR.grid(alpha=0.25, axis="x")
    for i, n in enumerate(order):
        axR.text(peaks[n], i, f" {peaks[n]:.0f}%", va="center", fontsize=8)

    fig.suptitle("Quem mais empurrou a inflacao chilena: energia lidera o pico, nucleo/servicos sao persistentes")
    fig.tight_layout(rect=(0, 0, 1, 0.94))
    fig.savefig(FIGURES / "ipc_cores_drivers.png", dpi=170)
    plt.close(fig)

    print("Peak YoY 2021-23 by component:")
    for n in reversed(order):
        print(f"  {n}: {peaks[n]:.1f}%")
    print(f"Wrote chile_inflation_cores.csv, ipc_cores_panel.png, ipc_cores_drivers.png "
          f"(through {last_full:%Y-%m})")


if __name__ == "__main__":
    main()
\end{lstlisting}
\clearpage
\section{python/build\_inflation\_cores\_bcch.py}
\begin{lstlisting}[language=Python]
"""Chilean CPI cores from the BCCh (real, recent: base 2023=100, up to 2026).

Pulls the BCCh analytical CPI breakdown (general, sin volatiles, goods/services
sin volatiles, volatile, volatile food/energy) as monthly variations (MoM, %)
from the SieteRestWS, then builds the central-bank "inflation momentum" view per
core:

  * MoM  (monthly change, annualised)        -> bars
  * MM3M (3-month change, annualised)         -> line
  * MM6M (6-month change, annualised)         -> line
  * YoY  (12-month change)                    -> line

All four are expressed as annual % so they are directly comparable and sit
against the 2-4% tolerance band (in/out of target). This replaces the FRED cores
(frozen at 2023-12) with real data through the latest BCCh release.

Credentials: env BCCH_USER/BCCH_PASS or data/raw/bcch_credentials.json.
Outputs: data/clean/chile_inflation_cores_bcch.csv, ipc_cores_momentum.png.
"""

from __future__ import annotations

import matplotlib

matplotlib.use("Agg")
import matplotlib.pyplot as plt
import numpy as np
import pandas as pd

from common import DATA_CLEAN, FIGURES, ensure_directories
from build_chile_dataset import fetch_series, get_credentials

# Monthly CPI variation (MoM, %) series ids -- filled from discover_ipc_series.py.
# label -> F074.IPC... id
MAPPING: dict[str, str] = {
    "Cheia (IPC geral)": "F074.IPC.VAR.Z.EP23.C.M",
    "Nucleo (SAE)": "F074.IPCSAE.VAR.Z.EP23.Z.M",
    "Servicos": "F074.IPCS.VAR.Z.EP23.Z.M",
    "Bens": "F074.IPCB.VAR.Z.EP23.Z.M",
    "Energia": "F074.IPCE.VAR.Z.EP23.Z.M",
    "Alimentos": "F074.IPCA.VAR.Z.EP23.Z.M",
}
COLORS = {
    "Cheia (IPC geral)": "#111827",
    "Nucleo (SAE)": "#1E4E79",
    "Servicos": "#16A34A",
    "Bens": "#7C3AED",
    "Energia": "#B45309",
    "Alimentos": "#0891B2",
}
TARGET, BAND = 3.0, (2.0, 4.0)
START_PLOT = "2021-01-01"


def momentum(mom_pct: pd.Series) -> pd.DataFrame:
    """Monthly variation (MoM, %), its 3/6-month moving averages (%/mo), YoY (%).

    MoM/MM3M/MM6M stay on the monthly scale (left axis); YoY is the accumulated
    12-month change on the annual scale (right axis). Annualising single monthly
    prints is avoided -- it explodes for volatile cores (energy tariff months).
    """
    s = mom_pct.sort_index()
    f = 1.0 + s / 100.0  # gross monthly factors
    out = pd.DataFrame(index=s.index)
    out["MoM"] = s                                  # monthly variation, %
    out["MM3M"] = s.rolling(3).mean()               # 3-month moving average, %/mo
    out["MM6M"] = s.rolling(6).mean()               # 6-month moving average, %/mo
    out["YoY"] = (f.rolling(12).apply(np.prod, raw=True) - 1.0) * 100.0  # 12-month, %
    return out


def main() -> None:
    ensure_directories()
    creds = get_credentials()
    if not creds:
        raise SystemExit("No BCCh credentials (set BCCH_USER/BCCH_PASS or "
                         "data/raw/bcch_credentials.json).")
    user, password = creds
    if len(MAPPING) < 2:
        raise SystemExit("Fill MAPPING with the core series ids first "
                         "(run discover_ipc_series.py).")

    cores = {label: momentum(fetch_series(sid, user, password, first="2018-01-01"))
             for label, sid in MAPPING.items()}

    # Tidy CSV: every core x metric.
    wide = pd.concat({label: df for label, df in cores.items()}, axis=1)
    wide.index.name = "date"
    wide.to_csv(DATA_CLEAN / "chile_inflation_cores_bcch.csv")
    last = max(df["YoY"].dropna().index.max() for df in cores.values())

    labels = list(cores)
    ncol = 3
    nrow = int(np.ceil(len(labels) / ncol))
    fig, axes = plt.subplots(nrow, ncol, figsize=(5.6 * ncol, 3.4 * nrow),
                             sharex=True)
    axes = np.atleast_1d(axes).ravel()
    handles: list = []
    hlabels: list[str] = []
    for ax, label in zip(axes, labels):
        df = cores[label].loc[START_PLOT:]
        color = COLORS.get(label, "#1E4E79")
        # Left axis (monthly %): MoM bars + 3/6-month moving averages.
        ax.axhline(0, color="#9CA3AF", lw=0.6, zorder=1)
        b = ax.bar(df.index, df["MoM"], width=22, color=color, alpha=0.28,
                   label="MoM", zorder=2)
        l3, = ax.plot(df.index, df["MM3M"], color=color, lw=1.6, label="MM3M",
                      zorder=3)
        l6, = ax.plot(df.index, df["MM6M"], color=color, lw=1.2, ls=(0, (4, 2)),
                      label="MM6M", zorder=3)
        ax.tick_params(axis="y", labelsize=7, colors=color)
        ax.tick_params(axis="x", labelsize=8)
        # Right axis (annual %): YoY vs the 2-4% tolerance band.
        ax2 = ax.twinx()
        ax2.axhspan(BAND[0], BAND[1], color="#FDE68A", alpha=0.40, zorder=0)
        ax2.axhline(TARGET, color="#B45309", ls="--", lw=0.8, zorder=1)
        ly, = ax2.plot(df.index, df["YoY"], color="#111827", lw=1.9, label="YoY",
                       zorder=4)
        ax2.tick_params(axis="y", labelsize=7)
        yoy_last = df["YoY"].dropna().iloc[-1]
        inside = BAND[0] <= yoy_last <= BAND[1]
        ax.set_title(f"{label} -- YoY {yoy_last:.1f}% "
                     f"({'dentro' if inside else 'fora'} da banda)", fontsize=9.5)
        ax.grid(alpha=0.18)
        if not handles:
            handles = [b, l3, l6, ly]
            hlabels = ["MoM (mensal, esq.)", "MM3M (esq.)", "MM6M (esq.)",
                       "YoY (anual, dir.)"]
    for ax in axes[len(labels):]:
        ax.set_visible(False)
    fig.legend(handles, hlabels, loc="lower center", fontsize=9, ncol=4,
               bbox_to_anchor=(0.5, 0.03), frameon=False)
    fig.suptitle("IPC chileno por nucleo -- barras: MoM e linhas finas MM3M/MM6M "
                 "(eixo esq., %/mes)\n* linha preta: YoY (eixo dir., % a.a.) vs banda "
                 "de tolerancia 2-4%", fontsize=11.5)
    fig.text(0.01, 0.004, "Fonte: Banco Central de Chile -- IPC Analiticos empalmados "
             f"(cifras oficiais INE), base 2023=100, SieteRestWS. Ultima obs.: {last:%Y-%m}.",
             fontsize=7.5)
    fig.tight_layout(rect=(0, 0.06, 1, 0.93))
    fig.savefig(FIGURES / "ipc_cores_momentum.png", dpi=170)
    plt.close(fig)

    # ------------------------------------------------------------------ #
    # Figure 2: drivers -- YoY timeline + recent momentum (MM3M) ranking. #
    # ------------------------------------------------------------------ #
    timeline = [l for l in ("Cheia (IPC geral)", "Nucleo (SAE)", "Energia",
                            "Servicos") if l in cores]
    fig, (axL, axR) = plt.subplots(1, 2, figsize=(13.5, 5.0),
                                   gridspec_kw={"width_ratios": [2.1, 1.0]})
    for label in timeline:
        s = cores[label]["YoY"].loc[START_PLOT:].dropna()
        axL.plot(s.index, s.values, color=COLORS.get(label, "#1E4E79"), lw=2.0,
                 label=label)
    axL.axhspan(BAND[0], BAND[1], color="#FDE68A", alpha=0.30)
    axL.axhline(TARGET, color="#B45309", ls="--", lw=1.0)
    axL.set_title("Inflacao anual (YoY) por nucleo")
    axL.set_ylabel("% a.a.")
    axL.grid(alpha=0.25)
    axL.legend(fontsize=8)
    latest_mm3 = {l: float(cores[l]["MM3M"].dropna().iloc[-1]) for l in cores}
    order = sorted(latest_mm3, key=latest_mm3.get)
    axR.barh(range(len(order)), [latest_mm3[l] for l in order],
             color=[COLORS.get(l, "#1E4E79") for l in order])
    axR.set_yticks(range(len(order)))
    axR.set_yticklabels(order, fontsize=8)
    axR.axvline(0, color="#9CA3AF", lw=0.6)
    axR.set_title(f"Momentum recente -- MM3M (%/mes), {last:%Y-%m}")
    axR.set_xlabel("%/mes")
    axR.grid(alpha=0.25, axis="x")
    for i, l in enumerate(order):
        axR.text(latest_mm3[l], i, f" {latest_mm3[l]:.1f}", va="center", fontsize=8)
    fig.suptitle("Quem empurra a inflacao chilena: a energia lidera o momentum; "
                 "servicos e o persistente")
    fig.tight_layout(rect=(0, 0, 1, 0.94))
    fig.savefig(FIGURES / "ipc_cores_drivers_bcch.png", dpi=170)
    plt.close(fig)

    print(f"Wrote chile_inflation_cores_bcch.csv + ipc_cores_momentum.png "
          f"+ ipc_cores_drivers_bcch.png (through {last:%Y-%m}).")
    for label, df in cores.items():
        print(f"  {label:28s} YoY {df['YoY'].dropna().iloc[-1]:5.1f}%  "
              f"MM3M {df['MM3M'].dropna().iloc[-1]:5.1f}%")


if __name__ == "__main__":
    main()
\end{lstlisting}
\clearpage
\section{python/build\_labor\_activity.py}
\begin{lstlisting}[language=Python]
"""Chilean activity and labour-market analysis from public FRED data, with an
Atlanta-Fed-style labour-market spider chart.

FRED has no Chilean vacancies/hires/quits (JOLTS is US-only), so the spider uses
the available Chilean equivalents: unemployment, youth unemployment, employment
rate, participation, industrial production and the model's (Kalman) output gap --
each as a percentile of its own history, oriented so OUTWARD = stronger labour
market/activity, current versus pre-pandemic (2019).

Outputs: labor_activity_dashboard.png, labor_spider.png, data/clean/chile_labor.csv.
"""

from __future__ import annotations

import io
import urllib.request

import matplotlib

matplotlib.use("Agg")
import matplotlib.pyplot as plt
import numpy as np
import pandas as pd

from common import DATA_CLEAN, FIGURES, ensure_directories

FRED_CSV = "https://fred.stlouisfed.org/graph/fredgraph.csv?id={series_id}"
SERIES = {
    "unemployment": "LRHUTTTTCLM156S",   # unemployment rate, %
    "youth_unemp": "SLUEM1524ZSCHL",     # youth (15-24) unemployment, %
    "employment": "LREM64TTCLM156S",     # employment rate 15-64, %
    "participation": "LRACTTTTCLM156S",  # participation rate, %
    "industrial": "CHLPROINDMISMEI",     # industrial production index
}


def download_fred(series_id: str) -> pd.Series:
    with urllib.request.urlopen(FRED_CSV.format(series_id=series_id), timeout=60) as resp:
        raw = resp.read().decode("utf-8", "replace")
    frame = pd.read_csv(io.StringIO(raw))
    frame.columns = ["date", "value"]
    frame["date"] = pd.to_datetime(frame["date"], errors="coerce")
    frame["value"] = pd.to_numeric(frame["value"], errors="coerce")
    frame = frame.dropna()
    return pd.Series(frame["value"].to_numpy(), index=pd.DatetimeIndex(frame["date"])).sort_index()


def to_quarterly(series: pd.Series) -> pd.Series:
    q = series.resample("QS").mean()
    q.index = q.index.to_period("Q").astype(str)
    return q


def main() -> None:
    ensure_directories()
    raw = {name: download_fred(sid) for name, sid in SERIES.items()}
    quarterly = pd.DataFrame({name: to_quarterly(s) for name, s in raw.items()})
    quarterly["industrial_yoy"] = (quarterly["industrial"] /
                                   quarterly["industrial"].shift(4) - 1.0) * 100.0

    macro = pd.read_csv(DATA_CLEAN / "chile_macro_quarterly.csv").set_index("date")
    quarterly["output_gap"] = macro["output_gap"] * 100.0  # Kalman gap, %
    quarterly = quarterly.loc["2001Q1":]
    quarterly.index.name = "date"
    quarterly.to_csv(DATA_CLEAN / "chile_labor.csv")
    dates = quarterly.index.tolist()
    xt = list(range(0, len(dates), 8))

    # ------------------------------------------------------------------ #
    # Figure 1: activity + labour dashboard.                              #
    # ------------------------------------------------------------------ #
    fig, axes = plt.subplots(1, 3, figsize=(14, 4.6))
    a = axes[0]
    a.plot(range(len(dates)), quarterly["unemployment"], color="#B91C1C", linewidth=2, label="Desemprego")
    a.plot(range(len(dates)), quarterly["youth_unemp"], color="#F59E0B", linewidth=1.6,
           linestyle="--", label="Desemprego jovem")
    a.set_title("Desemprego (%)"); a.legend(fontsize=8); a.grid(alpha=0.25)
    b = axes[1]
    b.plot(range(len(dates)), quarterly["employment"], color="#1E4E79", linewidth=2, label="Taxa de emprego")
    b.plot(range(len(dates)), quarterly["participation"], color="#16A34A", linewidth=2, label="Participacao")
    b.set_title("Emprego e participacao (%)"); b.legend(fontsize=8); b.grid(alpha=0.25)
    c = axes[2]
    c.axhline(0, color="black", linewidth=0.7)
    c.plot(range(len(dates)), quarterly["output_gap"], color="#7C3AED", linewidth=2, label="Hiato (Kalman)")
    c.plot(range(len(dates)), quarterly["industrial_yoy"], color="#6B7280", linewidth=1.4,
           linestyle="--", label="Prod. industrial (YoY)")
    c.set_title("Atividade (%)"); c.legend(fontsize=8); c.grid(alpha=0.25)
    for ax in axes:
        ax.set_xticks(xt); ax.set_xticklabels([dates[t] for t in xt], rotation=70, fontsize=7)
    fig.suptitle("Chile: mercado de trabalho e atividade (FRED + hiato do modelo)")
    fig.text(0.01, 0.01, "Fonte: OECD via FRED; hiato pelo filtro de Kalman. Sem dados de vagas/"
             "admissoes para o Chile. Okun: desemprego move-se inversamente ao hiato.", fontsize=7.5)
    fig.tight_layout(rect=(0, 0.03, 1, 0.95))
    fig.savefig(FIGURES / "labor_activity_dashboard.png", dpi=170)
    plt.close(fig)

    # ------------------------------------------------------------------ #
    # Figure 2: Atlanta-Fed-style spider (percentiles, outward=stronger). #
    # ------------------------------------------------------------------ #
    # Orientation: higher percentile = stronger labour market.
    indicators = {
        "Emprego": ("employment", +1),
        "Participacao": ("participation", +1),
        "Atividade\n(hiato)": ("output_gap", +1),
        "Prod.\nindustrial": ("industrial_yoy", +1),
        "Baixo desemp.\njovem": ("youth_unemp", -1),
        "Baixo\ndesemprego": ("unemployment", -1),
    }

    def percentile(series: pd.Series, value: float, sign: int) -> float:
        s = series.dropna()
        pct = 100.0 * (s < value).mean()
        return pct if sign > 0 else 100.0 - pct

    prepan = "2019Q4"
    labels, now_vals, pre_vals = [], [], []
    for label, (col, sign) in indicators.items():
        s = quarterly[col]
        labels.append(label)
        now_vals.append(percentile(s, float(s.dropna().iloc[-1]), sign))
        pre_vals.append(percentile(s, float(s.loc[prepan]) if prepan in s.index else float(s.dropna().iloc[-1]), sign))

    angles = np.linspace(0, 2 * np.pi, len(labels), endpoint=False).tolist()
    angles += angles[:1]
    fig, ax = plt.subplots(figsize=(8.2, 7.2), subplot_kw=dict(polar=True))
    for vals, color, name in ((pre_vals, "#9CA3AF", "Pre-pandemia (2019)"),
                              (now_vals, "#1E4E79", "Atual (ultimo dado)")):
        v = vals + vals[:1]
        ax.plot(angles, v, color=color, linewidth=2.2, label=name)
        ax.fill(angles, v, color=color, alpha=0.18)
    ax.set_xticks(angles[:-1]); ax.set_xticklabels(labels, fontsize=9)
    ax.set_yticks([25, 50, 75]); ax.set_yticklabels(["25", "50", "75"], fontsize=7)
    ax.set_ylim(0, 100)
    ax.set_title("Spider do mercado de trabalho chileno (percentis vs historia)\n"
                 "outward = mais forte", fontsize=11)
    ax.legend(loc="upper right", bbox_to_anchor=(1.25, 1.1), fontsize=9)
    fig.text(0.01, 0.01, "Estilo FED Atlanta: cada indicador em percentil vs sua historia (2001+), "
             "orientado p/ fora=mercado mais forte. Sem JOLTS/vagas para o Chile.", fontsize=7.5)
    fig.tight_layout()
    fig.savefig(FIGURES / "labor_spider.png", dpi=170)
    plt.close(fig)

    print("Latest labour percentiles (outward=stronger):")
    for label, nv in zip(labels, now_vals):
        print(f"  {label.replace(chr(10),' ')}: {nv:.0f}")
    print(f"Unemployment last {quarterly['unemployment'].dropna().iloc[-1]:.1f}% | "
          f"employment {quarterly['employment'].dropna().iloc[-1]:.1f}% | "
          f"gap {quarterly['output_gap'].dropna().iloc[-1]:.1f}%")
    print("Wrote labor_activity_dashboard.png, labor_spider.png, chile_labor.csv")


if __name__ == "__main__":
    main()
\end{lstlisting}
\clearpage
\part*{Graficos}\addcontentsline{toc}{part}{Graficos}

\section{python/plot\_history.py}
\begin{lstlisting}[language=Python]
"""Phase 1 of the deep macro analysis: historical narrative + policy counterfactual.

Reads the Kalman-smoothed shocks and the Dynare historical shock decomposition
(produced by run_history.py) and answers two macro questions for Chile, 2001-2026:

1. Narrative: which structural shocks (demand, cost-push, monetary) drove the
   observed inflation, output gap and policy rate each quarter?
2. Counterfactual: feeding the same recovered demand and cost-push shocks through
   the model under a MORE HAWKISH Taylor rule, how different would inflation and
   output have been (especially in the 2021-2023 inflation surge)?

The reduced form  y_t = T i_{t-1} + R eps_t  is the exact first-order solution
(identical to Dynare's). Counterfactual = re-solve with new policy parameters and
re-run the same recovered shock sequence.
"""

from __future__ import annotations

import matplotlib

matplotlib.use("Agg")
import matplotlib.pyplot as plt
import numpy as np
import pandas as pd

from common import (
    DATA_CLEAN,
    DYNARE_OUTPUTS,
    FIGURES,
    TABLES,
    _state_coefficients,
    calibration_rho,
    complete_parameters,
    ensure_directories,
)

HIST = DYNARE_OUTPUTS / "history"
SHOCK_LABELS = {
    "e_x": ("Demand", "#1E4E79"),
    "e_pi": ("Cost-push", "#B45309"),
    "e_i": ("Monetary", "#16A34A"),
    "initial": ("Initial cond.", "#9CA3AF"),
}


def solution_matrices(params: dict) -> tuple[np.ndarray, np.ndarray]:
    a_x, a_pi, a_i = _state_coefficients(params)
    sigma, beta, kappa = params["sigma"], params["beta"], params["kappa"]
    rho_i, phi_pi, phi_x = params["rho_i"], params["phi_pi"], params["phi_x"]
    c_i = a_x - (1.0 - a_pi) / sigma
    structural = np.array(
        [
            [1.0, 0.0, -c_i],
            [-kappa, 1.0, -beta * a_pi],
            [-(1.0 - rho_i) * phi_x, -(1.0 - rho_i) * phi_pi, 1.0],
        ]
    )
    return np.array([a_x, a_pi, a_i]), np.linalg.inv(structural)


def simulate(params: dict, shocks: np.ndarray, i0: float = 0.0) -> np.ndarray:
    transition, impact = solution_matrices(params)
    n = len(shocks)
    out = np.zeros((n, 3))
    prev_i = i0
    for t in range(n):
        out[t] = transition * prev_i + impact @ shocks[t]
        prev_i = out[t, 2]
    return out


def main() -> None:
    ensure_directories()
    params = complete_parameters({"rho_i": calibration_rho()})

    macro = pd.read_csv(DATA_CLEAN / "chile_macro_quarterly.csv")
    dates = macro["date"].tolist()
    n = len(dates)
    decomp = pd.read_csv(HIST / "shock_decomp.csv")
    shocks_df = pd.read_csv(HIST / "smoothed_shocks.csv")

    mean_infl_q = float(macro["infl_q"].mean())
    mean_i_q = float(macro["i_q"].mean())
    target_q = (1.03) ** 0.25 - 1.0  # inflation observable centred on the 3% target

    # Annualised percentage-point scaling (linear, keeps the decomposition additive).
    scale = {"pi": 400.0, "i": 400.0, "x": 100.0}
    unit = {"pi": "p.p. anual", "i": "p.p. anual", "x": "% do PIB"}
    titlevar = {"pi": "Inflacao", "i": "Taxa de politica (TPM)", "x": "Hiato do produto"}

    xticks = list(range(0, n, 8))
    xlabels = [dates[t] for t in xticks]

    # ------------------------------------------------------------------ #
    # Figure 1: historical shock decomposition (the macro narrative).    #
    # ------------------------------------------------------------------ #
    fig, axes = plt.subplots(3, 1, figsize=(12, 10), sharex=True)
    comps = ["e_x", "e_pi", "e_i", "initial"]
    for ax, var in zip(axes, ["pi", "i", "x"]):
        sub = decomp[decomp["variable"] == var].sort_values("period").reset_index(drop=True)
        s = scale[var]
        pos_bottom = np.zeros(n)
        neg_bottom = np.zeros(n)
        xs = np.arange(n)
        for comp in comps:
            vals = sub[comp].to_numpy() * s
            pos = np.where(vals > 0, vals, 0.0)
            neg = np.where(vals < 0, vals, 0.0)
            label, color = SHOCK_LABELS[comp]
            ax.bar(xs, pos, bottom=pos_bottom, color=color, width=0.9,
                   label=label, edgecolor="none")
            ax.bar(xs, neg, bottom=neg_bottom, color=color, width=0.9, edgecolor="none")
            pos_bottom += pos
            neg_bottom += neg
        ax.plot(xs, sub["smoothed"].to_numpy() * s, color="black", linewidth=1.3,
                label="Observado (desvio)")
        ax.axhline(0, color="black", linewidth=0.6)
        ax.set_title(f"{titlevar[var]} -- contribuicao de cada choque")
        ax.set_ylabel(unit[var])
        ax.grid(alpha=0.2, axis="y")
    axes[0].legend(ncol=5, fontsize=8, loc="upper left")
    axes[-1].set_xticks(xticks)
    axes[-1].set_xticklabels(xlabels, rotation=70, fontsize=7)
    fig.suptitle("Chile: decomposicao historica de choques (modelo NK calibrado, suavizador de Kalman)")
    fig.text(0.01, 0.005, "Desvios da media amostral; escala anualizada linear. "
             "Fonte: BCCh; decomposicao do modelo, nao identificacao historica oficial.", fontsize=7)
    fig.tight_layout(rect=(0, 0.02, 1, 0.97))
    fig.savefig(FIGURES / "history_shock_decomposition.png", dpi=170)
    plt.close(fig)

    # ------------------------------------------------------------------ #
    # Figure 2: recovered structural shocks over time.                   #
    # ------------------------------------------------------------------ #
    fig, axes = plt.subplots(3, 1, figsize=(12, 7), sharex=True)
    for ax, comp in zip(axes, ["e_x", "e_pi", "e_i"]):
        label, color = SHOCK_LABELS[comp]
        ax.bar(np.arange(n), shocks_df[comp].to_numpy() * 100, color=color, width=0.9)
        ax.axhline(0, color="black", linewidth=0.6)
        ax.set_title(f"{label} shock ({comp})")
        ax.set_ylabel("p.p. trimestral")
        ax.grid(alpha=0.2, axis="y")
    axes[-1].set_xticks(xticks)
    axes[-1].set_xticklabels(xlabels, rotation=70, fontsize=7)
    fig.suptitle("Chile: choques estruturais recuperados pelo suavizador de Kalman")
    fig.tight_layout(rect=(0, 0, 1, 0.96))
    fig.savefig(FIGURES / "history_smoothed_shocks.png", dpi=170)
    plt.close(fig)

    # ------------------------------------------------------------------ #
    # Figure 3: policy counterfactual (hawkish rule).                    #
    # ------------------------------------------------------------------ #
    shocks = shocks_df[["e_x", "e_pi", "e_i"]].to_numpy()
    base = simulate(params, shocks)                      # reproduces the data
    hawk_params = complete_parameters({"rho_i": params["rho_i"], "phi_pi": 2.5})
    hawk = simulate(hawk_params, shocks)
    # Validation: base-sim vs observed (deviations), after initial transient.
    obs = pd.read_csv(DATA_CLEAN / "chile_observables.csv")[["x", "pi", "i"]].to_numpy()
    rmse = float(np.sqrt(np.mean((base[15:] - obs[15:]) ** 2)))

    def to_annual_infl(dev):
        return ((1.0 + target_q + dev) ** 4 - 1.0) * 100.0

    def to_annual_tpm(dev):
        return ((1.0 + mean_i_q + dev) ** 4 - 1.0) * 100.0

    start = max(0, n - 40)  # focus on the last ~10 years
    xs = np.arange(start, n)
    fig, axes = plt.subplots(1, 2, figsize=(13, 4.8))
    axes[0].plot(xs, [to_annual_infl(d) for d in base[start:, 1]], color="black",
                 linewidth=2, label="Atual (regra estimada, phipi=1,75)")
    axes[0].plot(xs, [to_annual_infl(d) for d in hawk[start:, 1]], color="#B91C1C",
                 linewidth=2, linestyle="--", label="Contrafactual (phipi=2,5)")
    axes[0].axhline(3.0, color="#B45309", linestyle=":", linewidth=1, label="Meta 3%")
    axes[0].set_title("Inflacao (% a.a.)")
    axes[1].plot(xs, 100 * base[start:, 0], color="black", linewidth=2, label="Atual")
    axes[1].plot(xs, 100 * hawk[start:, 0], color="#B91C1C", linewidth=2,
                 linestyle="--", label="Contrafactual (phipi=2,5)")
    axes[1].axhline(0, color="#6B7280", linewidth=0.7)
    axes[1].set_title("Hiato do produto (%)")
    for ax in axes:
        ax.set_xticks(list(range(start, n, 4)))
        ax.set_xticklabels([dates[t] for t in range(start, n, 4)], rotation=70, fontsize=7)
        ax.grid(alpha=0.25)
        ax.legend(fontsize=8)
    fig.suptitle("Contrafactual de politica: regra mais agressiva contra a inflacao "
                 f"(validacao base vs dados: RMSE={rmse:.2e})")
    fig.tight_layout(rect=(0, 0, 1, 0.95))
    fig.savefig(FIGURES / "history_counterfactual.png", dpi=170)
    plt.close(fig)

    # Tidy export of the decomposition (annualised pp) for the report.
    rows = []
    for var in ["pi", "i", "x"]:
        sub = decomp[decomp["variable"] == var].sort_values("period").reset_index(drop=True)
        for t in range(n):
            rows.append({
                "date": dates[t], "variable": var,
                "demand": sub["e_x"][t] * scale[var],
                "cost_push": sub["e_pi"][t] * scale[var],
                "monetary": sub["e_i"][t] * scale[var],
                "initial": sub["initial"][t] * scale[var],
                "observed": sub["smoothed"][t] * scale[var],
            })
    pd.DataFrame(rows).to_csv(TABLES / "history_shock_decomposition.csv", index=False)

    print(f"Counterfactual validation RMSE (base-sim vs data, t>=15): {rmse:.3e}")
    print("Wrote history decomposition, smoothed shocks and counterfactual figures.")


if __name__ == "__main__":
    main()
\end{lstlisting}
\clearpage
\section{python/plot\_ipc\_decomposition.py}
\begin{lstlisting}[language=Python]
"""Decomposition of Chilean IPC (CPI) inflation into structural drivers.

Two figures:

1. ipc_inflation_decomposition.png -- the observed IPC inflation (annualised %) and
   how much of each quarter's deviation from the sample mean (3.8%) came from
   DEMAND, COST-PUSH and MONETARY shocks, per the Kalman-smoothed baseline model.
   The 2008 and 2022 inflation surges are highlighted.

2. ipc_inflation_models.png -- the SAME IPC inflation read by TWO models: the
   forward-looking baseline and the hybrid NKPC (with inflation memory). The hybrid
   needs much smaller cost-push shocks, because intrinsic persistence carries part
   of the inflation -- i.e., the macro story of 2022 depends on the model.
"""

from __future__ import annotations

import matplotlib

matplotlib.use("Agg")
import matplotlib.pyplot as plt
import numpy as np
import pandas as pd

from common import DATA_CLEAN, DYNARE_OUTPUTS, FIGURES, ensure_directories

HIST = DYNARE_OUTPUTS / "history"
HIST_HYB = DYNARE_OUTPUTS / "history_hybrid"
SHOCKS = {
    "e_x": ("Demanda", "#1E4E79"),
    "e_pi": ("Custo (cost-push)", "#B45309"),
    "e_i": ("Politica monetaria", "#16A34A"),
    "initial": ("Condicao inicial", "#9CA3AF"),
}
ANN = 400.0  # linear annualisation of a quarterly inflation deviation (keeps additivity)


def main() -> None:
    ensure_directories()
    macro = pd.read_csv(DATA_CLEAN / "chile_macro_quarterly.csv")
    dates = macro["date"].tolist()
    n = len(dates)
    mean_infl_annual = ((1.0 + macro["infl_q"].mean()) ** 4 - 1.0) * 100.0
    target_annual = 3.0  # observable is centred on the target; bars stack from 3%
    ipc = macro["infl_annual_pct"].to_numpy()

    decomp = pd.read_csv(HIST / "shock_decomp.csv")
    pi = decomp[decomp["variable"] == "pi"].sort_values("period").reset_index(drop=True)

    xticks = list(range(0, n, 8))
    xlabels = [dates[t] for t in xticks]
    # Episode shading (2008 and 2022 surges).
    def idx(label):
        return dates.index(label) if label in dates else None
    episodes = [("2008Q3", "2009Q2"), ("2021Q4", "2023Q1")]

    # ------------------------------------------------------------------ #
    # Figure 1: IPC inflation decomposition (baseline).                  #
    # ------------------------------------------------------------------ #
    fig, ax = plt.subplots(figsize=(12.5, 6.2))
    xs = np.arange(n)
    pos = np.full(n, target_annual)
    neg = np.full(n, target_annual)
    for comp in ("e_x", "e_pi", "e_i", "initial"):
        vals = pi[comp].to_numpy() * ANN
        p = np.where(vals > 0, vals, 0.0)
        m = np.where(vals < 0, vals, 0.0)
        label, color = SHOCKS[comp]
        ax.bar(xs, p, bottom=pos, color=color, width=0.92, label=label, edgecolor="none")
        ax.bar(xs, m, bottom=neg, color=color, width=0.92, edgecolor="none")
        pos += p
        neg += m
    ax.plot(xs, ipc, color="black", linewidth=1.6, label="IPC observado (% a.a.)")
    ax.axhline(mean_infl_annual, color="#374151", linewidth=1.0, linestyle="-",
               label=f"Media 2001--2026 ({mean_infl_annual:.1f}%)")
    ax.axhline(3.0, color="#B45309", linewidth=1.0, linestyle="--", label="Meta 3%")
    for a, b in episodes:
        ia, ib = idx(a), idx(b)
        if ia is not None and ib is not None:
            ax.axvspan(ia, ib, color="#FCD34D", alpha=0.18)
    ax.set_title("Decomposicao da inflacao do IPC chileno (modelo baseline, suavizador de Kalman)")
    ax.set_ylabel("Inflacao anualizada (% a.a.)")
    ax.set_xticks(xticks)
    ax.set_xticklabels(xlabels, rotation=70, fontsize=7)
    ax.grid(alpha=0.2, axis="y")
    ax.legend(ncol=3, fontsize=8, loc="upper center")
    fig.text(0.01, 0.005, "Barras = contribuicao de cada choque ao desvio da media (anualizado, "
             "aproximacao linear). Faixas amarelas: surtos de 2008 e 2022. "
             "Decomposicao do modelo, nao identificacao oficial.", fontsize=7)
    fig.tight_layout(rect=(0, 0.03, 1, 1))
    fig.savefig(FIGURES / "ipc_inflation_decomposition.png", dpi=170)
    plt.close(fig)

    # ------------------------------------------------------------------ #
    # Figure 2: the same IPC inflation, two models.                      #
    # ------------------------------------------------------------------ #
    sh_base = pd.read_csv(HIST / "smoothed_shocks.csv")
    sh_hyb = pd.read_csv(HIST_HYB / "smoothed_shocks.csv")
    std_base = sh_base["e_pi"].std()
    std_hyb = sh_hyb["e_pi"].std()

    fig, (axA, axB) = plt.subplots(1, 2, figsize=(13.5, 5.0))

    # Left: cost-push shock series, baseline vs hybrid.
    axA.bar(xs - 0.2, sh_base["e_pi"].to_numpy() * 100, width=0.4, color="#B45309",
            label=f"Baseline (sigma={std_base*100:.2f} p.p.)")
    axA.bar(xs + 0.2, sh_hyb["e_pi"].to_numpy() * 100, width=0.4, color="#B91C1C",
            label=f"Hibrida (sigma={std_hyb*100:.2f} p.p.)")
    axA.axhline(0, color="black", linewidth=0.7)
    for a, b in episodes:
        ia, ib = idx(a), idx(b)
        if ia is not None and ib is not None:
            axA.axvspan(ia, ib, color="#FCD34D", alpha=0.18)
    axA.set_title("Choque de custo recuperado (e_pi)")
    axA.set_ylabel("p.p. trimestral")
    axA.set_xticks(xticks)
    axA.set_xticklabels(xlabels, rotation=70, fontsize=7)
    axA.grid(alpha=0.2, axis="y")
    axA.legend(fontsize=9)

    # Right: cost-push contribution to inflation, baseline vs hybrid.
    pi_hyb = pd.read_csv(HIST_HYB / "shock_decomp.csv")
    pi_hyb = pi_hyb[pi_hyb["variable"] == "pi"].sort_values("period").reset_index(drop=True)
    axB.plot(xs, pi["e_pi"].to_numpy() * ANN, color="#B45309", linewidth=2.2,
             label="Baseline")
    axB.plot(xs, pi_hyb["e_pi"].to_numpy() * ANN, color="#B91C1C", linewidth=2.2,
             linestyle="--", label="Hibrida")
    axB.axhline(0, color="black", linewidth=0.7)
    for a, b in episodes:
        ia, ib = idx(a), idx(b)
        if ia is not None and ib is not None:
            axB.axvspan(ia, ib, color="#FCD34D", alpha=0.18)
    axB.set_title("Contribuicao do custo a inflacao (% a.a.)")
    axB.set_ylabel("p.p. anual")
    axB.set_xticks(xticks)
    axB.set_xticklabels(xlabels, rotation=70, fontsize=7)
    axB.grid(alpha=0.2, axis="y")
    axB.legend(fontsize=9)

    reduction = 100.0 * (std_base - std_hyb) / std_base
    fig.suptitle("A mesma inflacao do IPC, dois modelos: com inercia (hibrida) os choques de custo "
                 f"caem {reduction:.0f}%", fontsize=12)
    fig.tight_layout(rect=(0, 0, 1, 0.95))
    fig.savefig(FIGURES / "ipc_inflation_models.png", dpi=170)
    plt.close(fig)

    print(f"Mean IPC inflation {mean_infl_annual:.2f}% a.a. | cost-push std baseline "
          f"{std_base*100:.2f} vs hybrid {std_hyb*100:.2f} p.p. (-{reduction:.0f}%)")
    print("Wrote ipc_inflation_decomposition.png and ipc_inflation_models.png")


if __name__ == "__main__":
    main()
\end{lstlisting}
\clearpage
\section{python/plot\_forecast\_outlook.py}
\begin{lstlisting}[language=Python]
"""A direct-answer 'outlook' figure: where is Chile's economy heading?

Reads the native Dynare unconditional forecast and draws a clean, presentation-
ready dashboard for the next 8 quarters (TPM, inflation, output gap), with the
1-year-ahead point clearly annotated. The 1-year horizon is highlighted because
the pseudo-out-of-sample test (evaluate_forecasts.py) is where the NK model
actually beats naive benchmarks for inflation and the gap.
"""

from __future__ import annotations

import matplotlib

matplotlib.use("Agg")
import matplotlib.pyplot as plt
import numpy as np
import pandas as pd

from common import DATA_CLEAN, DYNARE_OUTPUTS, FIGURES, ensure_directories
from hybrid_solution import hybrid_forecast, hybrid_params, solve_hybrid

FORECAST = DYNARE_OUTPUTS / "forecast" / "forecast_unconditional.csv"


def main() -> None:
    ensure_directories()
    macro = pd.read_csv(DATA_CLEAN / "chile_macro_quarterly.csv")
    fc = pd.read_csv(FORECAST)

    mean_infl_q = float(macro["infl_q"].mean())
    mean_i_q = float(macro["i_q"].mean())
    target_q = (1.03) ** 0.25 - 1.0  # inflation observable is centred on the 3% target
    mean_infl_annual = ((1.0 + mean_infl_q) ** 4 - 1.0) * 100.0
    mean_tpm_annual = ((1.0 + mean_i_q) ** 4 - 1.0) * 100.0

    def to_level(var, dev):
        dev = np.asarray(dev, dtype=float)
        if var == "x":
            return 100.0 * dev
        if var == "pi":
            return ((1.0 + target_q + dev) ** 4 - 1.0) * 100.0
        return ((1.0 + mean_i_q + dev) ** 4 - 1.0) * 100.0

    last_period = pd.Period(str(macro.iloc[-1]["date"]), freq="Q")
    horizon = int(fc["horizon"].max())
    future = [str(last_period + h) for h in range(1, horizon + 1)]

    # Hybrid (data-preferred) forecast from the current state, for robustness.
    hybrid_sol = solve_hybrid(hybrid_params(0.35))
    i_dev0 = float(macro["i_q"].iloc[-1] - mean_i_q)
    pi_dev0 = float(macro["infl_q"].iloc[-1] - target_q)
    hybrid_dev = np.array(
        [hybrid_forecast(hybrid_sol, i_dev0, pi_dev0, h) for h in range(1, horizon + 1)]
    )  # (horizon, 3) deviations of [x, pi, i]
    var_col = {"x": 0, "pi": 1, "i": 2}

    hist = macro.tail(8).copy()
    hist_dates = hist["date"].tolist()
    n_hist = len(hist_dates)
    all_dates = hist_dates + future
    x_hist = np.arange(n_hist)
    x_fc = np.arange(n_hist - 1, n_hist - 1 + horizon + 1)  # connect last obs to forecast

    panels = [
        ("i", hist["tpm_annual_pct"].to_numpy(), "Taxa de politica (TPM, % a.a.)", mean_tpm_annual, None),
        ("pi", hist["infl_annual_pct"].to_numpy(), "Inflacao (% a.a.)", mean_infl_annual, 3.0),
        ("x", 100.0 * hist["output_gap"].to_numpy(), "Hiato do produto (%)", 0.0, None),
    ]

    fig, axes = plt.subplots(1, 3, figsize=(13.5, 5.0))
    for ax, (var, hist_vals, title, neutral, target) in zip(axes, panels):
        sub = fc[fc["variable"] == var].sort_values("horizon")
        mean_lv = to_level(var, sub["mean"].to_numpy())
        inf_lv = to_level(var, sub["hpd_inf"].to_numpy())
        sup_lv = to_level(var, sub["hpd_sup"].to_numpy())
        last_obs = float(hist_vals[-1])

        ax.plot(x_hist, hist_vals, color="black", marker="o", markersize=3.5,
                linewidth=1.6, label="Observado")
        ax.plot(x_fc, np.concatenate([[last_obs], mean_lv]), color="#1E4E79",
                linewidth=2.6, label="Previsao (central)")
        ax.fill_between(x_fc, np.concatenate([[last_obs], inf_lv]),
                        np.concatenate([[last_obs], sup_lv]),
                        color="#1E4E79", alpha=0.15, label="Banda 90%")
        # Hybrid forecast (data-preferred model) for robustness.
        hyb_lv = to_level(var, hybrid_dev[:, var_col[var]])
        ax.plot(x_fc, np.concatenate([[last_obs], hyb_lv]), color="#B91C1C",
                linewidth=1.7, linestyle="--", label="Hibrida (preferida pelos dados)")
        # Highlight the 1-year-ahead (horizon 4) point.
        h1y = 4
        x1y = n_hist - 1 + h1y
        y1y = float(mean_lv[h1y - 1])
        ax.scatter([x1y], [y1y], color="#B91C1C", zorder=5, s=45)
        ax.annotate(f"1 ano:\n{y1y:.1f}%", (x1y, y1y), textcoords="offset points",
                    xytext=(8, 10), fontsize=9, color="#B91C1C", fontweight="bold")

        ax.axvline(n_hist - 1, color="#9CA3AF", linewidth=0.9, linestyle="-")
        if target is not None:
            ax.axhline(target, color="#B45309", linestyle="--", linewidth=1, label="Meta 3%")
        ax.axhline(neutral, color="#6B7280", linestyle=":", linewidth=1,
                   label="Neutro/media" if var != "x" else "Potencial")
        ax.set_title(title, fontsize=11)
        step = 2
        ticks = list(range(0, len(all_dates), step))
        ax.set_xticks(ticks)
        ax.set_xticklabels([all_dates[t] for t in ticks], rotation=70, fontsize=7)
        ax.grid(alpha=0.22)
        ax.legend(fontsize=7, loc="best")

    fig.suptitle("Previsao do modelo para o Chile -- proximos 8 trimestres (incondicional, Dynare)",
                 fontsize=13)
    fig.text(0.01, 0.005,
             "Linha vertical = ultimo dado (2026Q1). Ponto vermelho = horizonte de 1 ano. A linha "
             "tracejada e a NKPC hibrida (preferida pelos dados): segura mais a inflacao no curto "
             "prazo, mas converge para a mesma previsao de 1 ano. Nao e previsao oficial do BCCh.",
             fontsize=7.5)
    fig.tight_layout(rect=(0, 0.04, 1, 0.95))
    fig.savefig(FIGURES / "forecast_outlook.png", dpi=180)
    plt.close(fig)
    print("Wrote forecast_outlook.png")


if __name__ == "__main__":
    main()
\end{lstlisting}
\clearpage
\section{python/plot\_hybrid\_test.py}
\begin{lstlisting}[language=Python]
"""A two-panel 'what the hybrid is, and whether it pays off' figure for the slide.

Left  : inflation response to a cost-push shock, baseline vs hybrid -- the hybrid
        carries inflation forward (intrinsic persistence 0.05 -> 0.38).
Right : out-of-sample inflation RMSE, baseline vs hybrid, at 1 and 4 quarters --
        the inertia cuts the 1-quarter error by ~19% and ties at 1 year.
"""

from __future__ import annotations

import matplotlib

matplotlib.use("Agg")
import matplotlib.pyplot as plt
import numpy as np
import pandas as pd

from common import DYNARE_OUTPUTS, FIGURES, TABLES, ensure_directories

BASE = "#1E4E79"
HYB = "#B91C1C"


def main() -> None:
    ensure_directories()
    base_irf = pd.read_csv(DYNARE_OUTPUTS / "baseline" / "irfs.csv")
    hyb_irf = pd.read_csv(DYNARE_OUTPUTS / "hybrid_nkpc" / "irfs.csv")
    oos = pd.read_csv(TABLES / "forecast_oos_metrics.csv")

    fig, (axL, axR) = plt.subplots(1, 2, figsize=(12, 4.7))

    # ---- Left: inflation IRF to a cost-push shock (annualised p.p.) ----
    h = base_irf["horizon"].to_numpy()
    axL.plot(h, base_irf["pi_e_pi"].to_numpy() * 400, color=BASE, linewidth=2.4,
             label="Baseline (sem inercia)")
    axL.plot(h, hyb_irf["pi_e_pi"].to_numpy() * 400, color=HYB, linewidth=2.4,
             linestyle="--", label="Hibrida (com inercia)")
    axL.axhline(0, color="black", linewidth=0.7)
    axL.set_title("Como a inflacao reage a um choque de custo")
    axL.set_xlabel("Trimestres apos o choque")
    axL.set_ylabel("Resposta da inflacao (p.p. anual)")
    axL.grid(alpha=0.25)
    axL.legend(fontsize=9)
    axL.annotate("Persistencia da inflacao:\nbaseline 0,05  ->  hibrida 0,38",
                 xy=(7.5, axL.get_ylim()[1] * 0.55), fontsize=9,
                 bbox=dict(boxstyle="round", fc="#FEF3C7", ec="#B45309", alpha=0.9))

    # ---- Right: out-of-sample inflation RMSE, baseline vs hybrid ----
    pi = oos[(oos["variable"] == "pi") & oos["model"].isin(["NK", "NK hibrido"])]
    def rmse(model, horizon):
        return float(pi[(pi["model"] == model) & (pi["horizon"] == horizon)]["rmse"].iloc[0])

    groups = ["1 trimestre", "1 ano (4 trim.)"]
    base_vals = [rmse("NK", 1), rmse("NK", 4)]
    hyb_vals = [rmse("NK hibrido", 1), rmse("NK hibrido", 4)]
    xpos = np.arange(2)
    width = 0.36
    axR.bar(xpos - width / 2, base_vals, width, color=BASE, label="Baseline")
    axR.bar(xpos + width / 2, hyb_vals, width, color=HYB, label="Hibrida")
    improvement = 100.0 * (base_vals[0] - hyb_vals[0]) / base_vals[0]
    axR.annotate(f"-{improvement:.0f}%", xy=(0, hyb_vals[0]), xytext=(0, hyb_vals[0] + 0.4),
                 ha="center", fontsize=11, fontweight="bold", color=HYB,
                 arrowprops=dict(arrowstyle="->", color=HYB))
    axR.annotate("empate", xy=(1, max(base_vals[1], hyb_vals[1])),
                 xytext=(1, max(base_vals[1], hyb_vals[1]) + 0.4), ha="center",
                 fontsize=10, color="#374151")
    axR.set_xticks(xpos)
    axR.set_xticklabels(groups)
    axR.set_title("Erro de previsao da inflacao (fora-da-amostra)")
    axR.set_ylabel("RMSE (p.p.)")
    axR.set_ylim(0, max(base_vals) * 1.25)
    axR.grid(alpha=0.25, axis="y")
    axR.legend(fontsize=9)

    fig.suptitle("A NKPC hibrida: inercia de inflacao e ganho de previsao de curto prazo",
                 fontsize=12)
    fig.tight_layout(rect=(0, 0, 1, 0.95))
    fig.savefig(FIGURES / "hybrid_test.png", dpi=180)
    plt.close(fig)
    print(f"Wrote hybrid_test.png (1q improvement {improvement:.1f}%)")


if __name__ == "__main__":
    main()
\end{lstlisting}
\clearpage
\section{python/plot\_irfs.py}
\begin{lstlisting}[language=Python]
"""Create the three required IRF figures with matplotlib."""

from __future__ import annotations

import json
import matplotlib

matplotlib.use("Agg")
import matplotlib.pyplot as plt
import numpy as np
import pandas as pd

from common import DATA_CLEAN, FIGURES, TABLES, ensure_directories


def load_irfs() -> pd.DataFrame:
    path = TABLES / "irfs_long.csv"
    if not path.exists():
        from collect_outputs import main as collect_main

        collect_main()
    return pd.read_csv(path)


def source_note(frame: pd.DataFrame) -> str:
    sources = sorted(frame["source"].dropna().unique())
    if sources == ["dynare"]:
        return "Source: Dynare"
    return "Source includes clearly labeled synthetic model fallback"


def plot_kappa(irfs: pd.DataFrame) -> None:
    selected = irfs[
        irfs["scenario"].str.startswith("kappa_")
        & (irfs["variable"] == "pi")
        & (irfs["shock"] == "e_x")
    ]
    fig, ax = plt.subplots(figsize=(8.5, 5.2))
    for scenario, frame in selected.groupby("scenario"):
        label = scenario.removeprefix("kappa_").replace("p", ".")
        ax.plot(frame["horizon"], 100 * frame["response"], linewidth=2, label=label)
    ax.axhline(0, color="black", linewidth=0.8)
    ax.set(
        title="Inflation response to a demand shock: Phillips-slope (kappa) sensitivity",
        xlabel="Quarters",
        ylabel="Percentage points",
    )
    ax.legend(title="kappa")
    ax.grid(alpha=0.25)
    fig.text(0.01, 0.01, source_note(selected), fontsize=8)
    fig.tight_layout(rect=(0, 0.03, 1, 1))
    fig.savefig(FIGURES / "irf_kappa_comparison.png", dpi=180)
    plt.close(fig)


def plot_kappa_tradeoffs(irfs: pd.DataFrame) -> None:
    selected = irfs[irfs["scenario"].str.startswith("kappa_")]
    fig, axes = plt.subplots(1, 2, figsize=(12, 4.8), sharex=True)
    for scenario, frame in selected.groupby("scenario"):
        label = scenario.removeprefix("kappa_").replace("p", ".")
        for ax, shock, title in (
            (axes[0], "e_x", "Demand shock"),
            (axes[1], "e_pi", "Cost-push shock"),
        ):
            line = frame[(frame["variable"] == "pi") & (frame["shock"] == shock)]
            ax.plot(line["horizon"], 100 * line["response"], linewidth=2, label=label)
            ax.set_title(title)
            ax.set_xlabel("Quarters")
            ax.grid(alpha=0.25)
    axes[0].set_ylabel("Inflation response (percentage points)")
    axes[0].legend(title="kappa")
    for ax in axes:
        ax.axhline(0, color="black", linewidth=0.8)
    fig.suptitle("Phillips-curve slope: transmission and stabilization trade-offs")
    fig.text(0.01, 0.01, source_note(selected), fontsize=8)
    fig.tight_layout(rect=(0, 0.03, 1, 0.94))
    fig.savefig(FIGURES / "irf_kappa_tradeoffs.png", dpi=180)
    plt.close(fig)


def plot_phi_pi(irfs: pd.DataFrame) -> None:
    selected = irfs[
        irfs["scenario"].str.startswith("phi_pi_")
        & (irfs["variable"] == "pi")
        & (irfs["shock"] == "e_pi")
    ]
    scenarios = sorted(selected["scenario"].unique())
    colors = plt.cm.viridis(np.linspace(0.05, 0.95, len(scenarios)))
    fig, ax = plt.subplots(figsize=(8.8, 5.4))
    for color, scenario in zip(colors, scenarios):
        frame = selected[selected["scenario"] == scenario]
        label = scenario.removeprefix("phi_pi_").replace("p", ".")
        ax.plot(
            frame["horizon"],
            100 * frame["response"],
            color=color,
            linewidth=1.6,
            label=label,
        )
    ax.axhline(0, color="black", linewidth=0.8)
    ax.set(
        title="Inflation response to a cost-push shock: Taylor-rule sensitivity",
        xlabel="Quarters",
        ylabel="Percentage points",
    )
    ax.legend(title="phi_pi", ncol=2, fontsize=8)
    ax.grid(alpha=0.25)
    fig.text(0.01, 0.01, source_note(selected), fontsize=8)
    fig.tight_layout(rect=(0, 0.03, 1, 1))
    fig.savefig(FIGURES / "irf_phi_pi_comparison.png", dpi=180)
    plt.close(fig)


def plot_phi_pi_tradeoffs(irfs: pd.DataFrame) -> None:
    scenarios = {
        "phi_pi_1p3": ("1.3", "#5B8FF9"),
        "baseline": ("1.75 baseline", "#111827"),
        "phi_pi_2p2": ("2.2", "#F4664A"),
    }
    fig, axes = plt.subplots(1, 3, figsize=(14, 4.8), sharex=True)
    for scenario, (label, color) in scenarios.items():
        selected = irfs[
            (irfs["scenario"] == scenario) & (irfs["shock"] == "e_pi")
        ]
        for ax, variable, title in zip(
            axes,
            ("pi", "x", "i"),
            ("Inflation", "Output gap", "Policy rate"),
        ):
            line = selected[selected["variable"] == variable]
            ax.plot(
                line["horizon"],
                100 * line["response"],
                linewidth=2,
                label=label,
                color=color,
            )
            ax.axhline(0, color="black", linewidth=0.7)
            ax.set_title(title)
            ax.set_xlabel("Quarters")
            ax.grid(alpha=0.25)
    axes[0].set_ylabel("Response (percentage points)")
    axes[0].legend(title="phi_pi")
    fig.suptitle("Cost-push shock: inflation-output-interest trade-off")
    fig.text(0.01, 0.01, source_note(irfs[irfs["scenario"].isin(scenarios)]), fontsize=8)
    fig.tight_layout(rect=(0, 0.03, 1, 0.94))
    fig.savefig(FIGURES / "irf_phi_pi_tradeoffs.png", dpi=180)
    plt.close(fig)


def plot_rho_comparison(irfs: pd.DataFrame) -> None:
    scenarios = {
        "baseline": ("Estimated rho_i = 0.934", "#1E4E79"),
        "rho_calibrated_0p80": ("Calibrated rho_i = 0.80", "#FF6B35"),
    }
    fig, axes = plt.subplots(1, 3, figsize=(14, 4.8), sharex=True)
    for scenario, (label, color) in scenarios.items():
        selected = irfs[
            (irfs["scenario"] == scenario) & (irfs["shock"] == "e_i")
        ]
        for ax, variable, title in zip(
            axes,
            ("i", "x", "pi"),
            ("Policy rate", "Output gap", "Inflation"),
        ):
            line = selected[selected["variable"] == variable]
            ax.plot(
                line["horizon"],
                100 * line["response"],
                linewidth=2.2,
                label=label,
                color=color,
            )
            ax.axhline(0, color="black", linewidth=0.7)
            ax.set_title(title)
            ax.set_xlabel("Quarters")
            ax.grid(alpha=0.25)
    axes[0].set_ylabel("Response (percentage points)")
    axes[0].legend(fontsize=8)
    fig.suptitle("Interest-rate smoothing route: calibrated versus estimated rho_i")
    fig.text(0.01, 0.01, source_note(irfs[irfs["scenario"].isin(scenarios)]), fontsize=8)
    fig.tight_layout(rect=(0, 0.03, 1, 0.94))
    fig.savefig(FIGURES / "irf_rho_comparison.png", dpi=180)
    plt.close(fig)


def plot_baseline(irfs: pd.DataFrame) -> None:
    selected = irfs[irfs["scenario"] == "baseline"]
    variables = ("x", "pi", "i")
    shocks = ("e_x", "e_pi", "e_i")
    labels = {"x": "Output gap", "pi": "Inflation", "i": "Policy rate"}
    shock_labels = {
        "e_x": "Demand shock",
        "e_pi": "Cost-push shock",
        "e_i": "Monetary-policy shock",
    }
    fig, axes = plt.subplots(3, 3, figsize=(11.5, 8.5), sharex=True)
    for row, variable in enumerate(variables):
        for column, shock in enumerate(shocks):
            ax = axes[row, column]
            frame = selected[
                (selected["variable"] == variable) & (selected["shock"] == shock)
            ]
            ax.plot(frame["horizon"], 100 * frame["response"], linewidth=1.8)
            ax.axhline(0, color="black", linewidth=0.7)
            ax.grid(alpha=0.22)
            if row == 0:
                ax.set_title(shock_labels[shock])
            if column == 0:
                ax.set_ylabel(f"{labels[variable]}\npp")
            if row == 2:
                ax.set_xlabel("Quarters")
    fig.suptitle("Baseline impulse responses across all shocks", fontsize=14)
    fig.text(0.01, 0.01, source_note(selected), fontsize=8)
    fig.tight_layout(rect=(0, 0.03, 1, 0.96))
    fig.savefig(FIGURES / "irf_baseline_all_shocks.png", dpi=180)
    plt.close(fig)


def plot_data_overview() -> None:
    path = DATA_CLEAN / "chile_macro_quarterly.csv"
    if not path.exists():
        return
    df = pd.read_csv(path)
    metadata_path = DATA_CLEAN / "dataset_metadata.json"
    metadata = (
        json.loads(metadata_path.read_text(encoding="utf-8"))
        if metadata_path.exists()
        else {"is_synthetic": True}
    )
    is_synthetic = bool(metadata.get("is_synthetic", True))
    source = (
        "Banco Central de Chile (TPM, IPC, PIB)"
        if not is_synthetic
        else "Synthetic fallback for pipeline testing"
    )
    fig, axes = plt.subplots(3, 1, figsize=(9, 8.2), sharex=True)
    axes[0].plot(df["date"], df["tpm_annual_pct"], color="#1E4E79")
    axes[0].set_ylabel("TPM (% a.a.)")
    axes[1].plot(df["date"], df["infl_annual_pct"], color="#FF6B35")
    axes[1].axhline(3.0, color="gray", ls="--", lw=0.8)
    axes[1].set_ylabel("Inflacao IPC (% a.a.)")
    axes[2].plot(df["date"], 100.0 * df["output_gap"], color="#16a34a")
    axes[2].axhline(0.0, color="black", lw=0.7)
    axes[2].set_ylabel("Hiato (%)")
    for ax in axes:
        ax.grid(alpha=0.25)
    step = max(1, len(df) // 12)
    axes[2].set_xticks(df["date"][::step])
    axes[2].tick_params(axis="x", rotation=90, labelsize=7)
    title_kind = "official quarterly data" if not is_synthetic else "synthetic test data"
    axes[0].set_title(f"Chile: {title_kind} - policy rate, inflation, output gap")
    fig.text(
        0.01,
        0.005,
        f"Source: {source}; output gap = HP cycle of log real GDP.",
        fontsize=7,
    )
    fig.tight_layout(rect=(0, 0.03, 1, 1))
    fig.savefig(FIGURES / "data_overview.png", dpi=180)
    plt.close(fig)


def main() -> None:
    ensure_directories()
    plot_data_overview()
    irfs = load_irfs()
    plot_kappa(irfs)
    plot_kappa_tradeoffs(irfs)
    plot_phi_pi(irfs)
    plot_phi_pi_tradeoffs(irfs)
    plot_rho_comparison(irfs)
    plot_baseline(irfs)
    print(f"Wrote IRF figures to {FIGURES}")


if __name__ == "__main__":
    main()
\end{lstlisting}
\clearpage
\section{python/plot\_mcmc.py}
\begin{lstlisting}[language=Python]
"""Plot Bayesian posterior summaries and convergence diagnostics."""

from __future__ import annotations

import matplotlib.pyplot as plt
import numpy as np
import pandas as pd

from common import FIGURES, TABLES


LABELS = {
    "stderr e_x": r"$\sigma_{e_x}$",
    "stderr e_pi": r"$\sigma_{e_\pi}$",
    "stderr e_i": r"$\sigma_{e_i}$",
    "sigma": r"$\sigma$",
    "kappa": r"$\kappa$",
    "rho_i": r"$\rho_i$",
    "phi_pi": r"$\phi_\pi$",
    "phi_x": r"$\phi_x$",
}


def main() -> None:
    posterior = pd.read_csv(TABLES / "mcmc_posterior.csv")
    diagnostics = pd.read_csv(TABLES / "mcmc_diagnostics.csv")
    merged = posterior.merge(diagnostics, on="parameter", validate="one_to_one")

    y = np.arange(len(merged))
    means = merged["posterior_mean"].to_numpy()
    low = merged["hpd90_low"].to_numpy()
    high = merged["hpd90_high"].to_numpy()

    fig, axes = plt.subplots(1, 2, figsize=(12, 5.8))
    axes[0].errorbar(
        means,
        y,
        xerr=np.vstack([means - low, high - means]),
        fmt="o",
        color="#145DA0",
        ecolor="#72A0C1",
        capsize=3,
    )
    axes[0].set_yticks(y, [LABELS.get(p, p) for p in merged["parameter"]])
    axes[0].invert_yaxis()
    axes[0].set_title("Posterior mean and 90% credible interval")
    axes[0].set_xlabel("Parameter value")
    axes[0].grid(axis="x", alpha=0.25)

    colors = np.where(merged["rhat"] < 1.05, "#2E8B57", "#C0392B")
    axes[1].barh(y, merged["rhat"] - 1.0, left=1.0, color=colors)
    axes[1].axvline(1.05, color="#333333", linestyle="--", linewidth=1)
    axes[1].set_yticks(y, [LABELS.get(p, p) for p in merged["parameter"]])
    axes[1].invert_yaxis()
    axes[1].set_xlim(0.99, max(1.08, merged["rhat"].max() * 1.02))
    axes[1].set_title(r"Convergence diagnostic ($\hat R$)")
    axes[1].set_xlabel(r"$\hat R$; values below 1.05 are preferred")
    axes[1].grid(axis="x", alpha=0.25)

    fig.suptitle("Bayesian MCMC extension: posterior uncertainty and convergence")
    fig.text(
        0.5,
        0.01,
        "Two Dynare chains; 10,000 proposals each; first 30% discarded. "
        "Results are model- and prior-conditional.",
        ha="center",
        fontsize=9,
    )
    fig.tight_layout(rect=(0, 0.04, 1, 0.96))
    target = FIGURES / "mcmc_posterior_diagnostics.png"
    fig.savefig(target, dpi=180, bbox_inches="tight")
    plt.close(fig)
    print(f"Wrote {target}")


if __name__ == "__main__":
    main()
\end{lstlisting}
\clearpage
\part*{Tabelas, coleta e entregaveis}\addcontentsline{toc}{part}{Tabelas, coleta e entregaveis}

\section{python/collect\_outputs.py}
\begin{lstlisting}[language=Python]
"""Collect Dynare CSV exports and create explicit synthetic fallbacks."""

from __future__ import annotations

import math
import pandas as pd

from common import (
    DYNARE_OUTPUTS,
    TABLES,
    calibration_rho,
    complete_parameters,
    ensure_directories,
    fallback_irfs,
    fallback_moments,
)

VARIABLES = ("x", "pi", "i")
SHOCKS = ("e_x", "e_pi", "e_i")


def read_dynare_fevd(scenario: str) -> pd.DataFrame | None:
    path = DYNARE_OUTPUTS / scenario / "fevd.csv"
    if not path.exists():
        return None
    wide = pd.read_csv(path)
    records = []
    for _, row in wide.iterrows():
        values = [float(row[shock]) for shock in SHOCKS]
        scale = 100.0 if sum(values) > 1.5 else 1.0
        for shock, value in zip(SHOCKS, values):
            records.append(
                {
                    "scenario": scenario,
                    "variable": row["variable"],
                    "shock": shock,
                    "variance_share": value / scale,
                    "source": "dynare",
                }
            )
    return pd.DataFrame(records)


def dynare_irfs(scenario: str, path) -> pd.DataFrame:
    wide = pd.read_csv(path)
    records = []
    for variable in VARIABLES:
        for shock in SHOCKS:
            column = f"{variable}_{shock}"
            if column not in wide:
                continue
            for horizon, response in zip(wide["horizon"], wide[column]):
                records.append(
                    {
                        "scenario": scenario,
                        "source": "dynare",
                        "variable": variable,
                        "shock": shock,
                        "horizon": int(horizon),
                        "response": float(response),
                    }
                )
    return pd.DataFrame(records)


def manifest_with_baseline() -> pd.DataFrame:
    path = TABLES / "scenario_manifest.csv"
    manifest = pd.read_csv(path) if path.exists() else pd.DataFrame()
    baseline = {"scenario": "baseline", "scenario_type": "baseline"}
    baseline.update(complete_parameters({"rho_i": calibration_rho()}))
    return pd.concat([pd.DataFrame([baseline]), manifest], ignore_index=True)


def row_parameters(row: pd.Series) -> dict:
    keys = (
        "rstar_annual",
        "sigma",
        "kappa",
        "rho_i",
        "phi_pi",
        "phi_x",
        "rstar",
        "beta",
    )
    return {key: float(row[key]) for key in keys if key in row and pd.notna(row[key])}


def collect_irfs(manifest: pd.DataFrame) -> pd.DataFrame:
    pieces = []
    for _, row in manifest.iterrows():
        scenario = str(row["scenario"])
        irf_path = DYNARE_OUTPUTS / scenario / "irfs.csv"
        if irf_path.exists():
            try:
                frame = dynare_irfs(scenario, irf_path)
                if not frame.empty:
                    pieces.append(frame)
                    continue
            except Exception as exc:
                print(f"Warning: could not parse {irf_path}: {exc}")
        params = complete_parameters(row_parameters(row))
        pieces.append(fallback_irfs(params, scenario))
    result = pd.concat(pieces, ignore_index=True)
    result.to_csv(TABLES / "irfs_long.csv", index=False)
    return result


def collect_moments(manifest: pd.DataFrame) -> None:
    frames = []
    for _, row in manifest.iterrows():
        scenario = str(row["scenario"])
        path = DYNARE_OUTPUTS / scenario / "moments.csv"
        if path.exists():
            moment = pd.read_csv(path)
            moment.insert(0, "scenario", scenario)
            moment["source"] = "dynare"
            frames.append(moment)
        else:
            params = complete_parameters(row_parameters(row))
            moment = fallback_moments(params)
            moment.insert(0, "scenario", scenario)
            frames.append(moment)
    if not frames:
        baseline = manifest.loc[manifest["scenario"] == "baseline"].iloc[0]
        moment = fallback_moments(complete_parameters(row_parameters(baseline)))
        moment.insert(0, "scenario", "baseline")
        frames.append(moment)
    moments_all = pd.concat(frames, ignore_index=True)
    moments_all.to_csv(TABLES / "moments_all_scenarios.csv", index=False)
    moments_all[moments_all["scenario"] == "baseline"].to_csv(
        TABLES / "moments.csv", index=False
    )


def collect_fevd(irfs: pd.DataFrame, manifest: pd.DataFrame) -> None:
    frames = []
    for scenario in manifest["scenario"].astype(str):
        frame = read_dynare_fevd(scenario)
        if frame is not None:
            frames.append(frame)
            continue
        selected = irfs[irfs["scenario"] == scenario].copy()
        selected["squared_response"] = selected["response"] ** 2
        totals = (
            selected.groupby(["variable", "shock"], as_index=False)["squared_response"]
            .sum()
            .rename(columns={"squared_response": "shock_contribution"})
        )
        totals["variance_share"] = totals["shock_contribution"] / totals.groupby(
            "variable"
        )["shock_contribution"].transform("sum")
        totals["scenario"] = scenario
        totals["source"] = "synthetic_irf_share_fallback"
        frames.append(
            totals[["scenario", "variable", "shock", "variance_share", "source"]]
        )
    fevd_all = pd.concat(frames, ignore_index=True)
    fevd_all.to_csv(TABLES / "fevd_all_scenarios.csv", index=False)
    fevd_all[fevd_all["scenario"] == "baseline"].to_csv(
        TABLES / "fevd_summary.csv", index=False
    )

def collect_determinacy(manifest: pd.DataFrame) -> None:
    rows = []
    for _, scenario_row in manifest.iterrows():
        scenario = str(scenario_row["scenario"])
        stability_path = DYNARE_OUTPUTS / scenario / "stability.csv"
        if stability_path.exists():
            row = pd.read_csv(stability_path).iloc[0].to_dict()
            row["source"] = "dynare_check"
        else:
            row = {
                "scenario": scenario,
                "status": "not_run_or_failed",
                "n_forward": None,
                "n_unstable": None,
                "n_eigenvalues": None,
                "source": "no_dynare_diagnostic",
            }
        eigen_path = DYNARE_OUTPUTS / scenario / "eigenvalues.csv"
        if eigen_path.exists():
            eigenvalues = pd.read_csv(eigen_path)
            stable = eigenvalues[eigenvalues["modulus"] < 1.0]["modulus"]
            unstable = eigenvalues[eigenvalues["modulus"] > 1.0]["modulus"]
            dominant = float(stable.max()) if not stable.empty else math.nan
            row["dominant_stable_modulus"] = dominant
            row["convergence_half_life_quarters"] = (
                math.log(0.5) / math.log(dominant)
                if 0.0 < dominant < 1.0
                else math.nan
            )
            row["smallest_unstable_modulus"] = (
                float(unstable.min()) if not unstable.empty else math.nan
            )
        else:
            row["dominant_stable_modulus"] = math.nan
            row["convergence_half_life_quarters"] = math.nan
            row["smallest_unstable_modulus"] = math.nan
        row["scenario_type"] = scenario_row.get("scenario_type", "")
        row["phi_pi"] = scenario_row.get("phi_pi", None)
        rows.append(row)
    pd.DataFrame(rows).to_csv(TABLES / "scenario_determinacy.csv", index=False)


def main() -> None:
    ensure_directories()
    manifest = manifest_with_baseline()
    irfs = collect_irfs(manifest)
    collect_moments(manifest)
    collect_fevd(irfs, manifest)
    collect_determinacy(manifest)
    print(f"Collected outputs in {TABLES}")
    print(irfs.groupby("source")["scenario"].nunique().to_string())


if __name__ == "__main__":
    main()
\end{lstlisting}
\clearpage
\section{python/make\_tables.py}
\begin{lstlisting}[language=Python]
"""Create the parameter-target table and the neutral-rate conversion table.

The parameter-target table is the assignment's headline "parameter <-> target" deliverable.
It lists each calibrated value next to its empirical counterpart (when an estimate exists:
AR(1) for rho_i, IV/OLS Taylor rule for phi_pi/phi_x, NKPC for kappa, HP proxy for r*).
"""

from __future__ import annotations

import pandas as pd

from common import (
    BASELINE,
    SHOCK_STD,
    TABLES,
    calibration_rho,
    complete_parameters,
    ensure_directories,
    quarterly_rstar,
)


def _safe_read(name: str) -> pd.DataFrame | None:
    path = TABLES / name
    return pd.read_csv(path) if path.exists() else None


def empirical_estimates() -> dict[str, str]:
    emp: dict[str, str] = {}
    rhoi = _safe_read("rhoi_estimate.csv")
    if (
        rhoi is not None
        and not rhoi.empty
        and not bool(rhoi.iloc[0].get("is_synthetic", True))
    ):
        r = rhoi.iloc[0]
        emp["rho_i"] = f"{r['rho_i_estimate']:.3f} (AR1; HAC SE {r['rho_i_se_hac']:.3f})"
    taylor = _safe_read("taylor_rule_estimates.csv")
    if (
        taylor is not None
        and not taylor.empty
        and not bool(taylor.iloc[0].get("is_synthetic", True))
    ):
        iv = taylor[taylor["method"].str.startswith("IV_2SLS")]
        ols = taylor[taylor["method"] == "OLS_HAC"]
        if not iv.empty and not ols.empty:
            emp["phi_pi"] = (
                f"{iv.iloc[0]['phi_pi']:.2f} (IV-HAC); "
                f"{ols.iloc[0]['phi_pi']:.2f} (OLS-HAC)"
            )
            emp["phi_x"] = f"{iv.iloc[0]['phi_x']:.2f} (IV-HAC)"
    nkpc = _safe_read("nkpc_estimates.csv")
    if (
        nkpc is not None
        and not nkpc.empty
        and not bool(nkpc.iloc[0].get("is_synthetic", True))
    ):
        bw = nkpc[nkpc["method"] == "backward_OLS_HAC"]
        if not bw.empty:
            emp["kappa"] = f"{bw.iloc[0]['kappa']:.3f} (backward OLS, HAC)"
    rstar = _safe_read("rstar_estimates.csv")
    if (
        rstar is not None
        and not rstar.empty
        and not bool(rstar.iloc[0].get("is_synthetic", True))
    ):
        hp = rstar[rstar["method"] == "hp_trend_real_rate_end_of_sample"]
        if not hp.empty:
            emp["rstar_annual"] = f"{hp.iloc[0]['rstar_annual_pct']:.2f}% a.a. (HP proxy)"
    return emp


def main() -> None:
    ensure_directories()
    rho_i = calibration_rho()
    beta = complete_parameters()["beta"]
    emp = empirical_estimates()

    rows = [
        ("rstar_annual", BASELINE["rstar_annual"], emp.get("rstar_annual", ""),
         "Annual neutral real rate; scenarios 2/3/4%"),
        ("beta", round(beta, 6), "",
         "Discount factor = 1/(1+rstar_q), derived from r*"),
        ("sigma", BASELINE["sigma"], "",
         "Inverse intertemporal elasticity; log-utility benchmark"),
        ("kappa", BASELINE["kappa"], emp.get("kappa", ""),
         "Phillips-curve slope; sensitivity grid {0.07,0.10,0.13}"),
        ("rho_i", round(rho_i, 4), emp.get("rho_i", ""),
         "Interest-rate smoothing; AR(1) of the quarterly TPM"),
        ("phi_pi", BASELINE["phi_pi"], emp.get("phi_pi", ""),
         "Inflation response; grid [1.3,2.2]; determinacy checked numerically"),
        ("phi_x", BASELINE["phi_x"], emp.get("phi_x", ""),
         "Output-gap response (held fixed at 0.5)"),
        ("std_e_x", SHOCK_STD["e_x"], "",
         "Demand-shock std; illustrative model-based FEVD calibration"),
        ("std_e_pi", SHOCK_STD["e_pi"], "",
         "Cost-push-shock std; illustrative model-based FEVD calibration"),
        ("std_e_i", SHOCK_STD["e_i"], "",
         "Monetary-shock std; illustrative model-based FEVD calibration"),
    ]
    parameter_table = pd.DataFrame(
        rows,
        columns=["parameter", "calibrated_value", "empirical_estimate",
                 "target_or_rationale"],
    )
    parameter_table["country"] = "Chile"
    parameter_table.to_csv(TABLES / "parameter_targets.csv", index=False)

    rows_rstar = []
    for annual in (0.02, 0.03, 0.04):
        quarterly = quarterly_rstar(annual)
        rows_rstar.append(
            {
                "rstar_annual": annual,
                "rstar_annual_percent": annual * 100.0,
                "rstar_quarterly": quarterly,
                "rstar_quarterly_percent": quarterly * 100.0,
                "beta": 1.0 / (1.0 + quarterly),
            }
        )
    pd.DataFrame(rows_rstar).to_csv(TABLES / "rstar_beta_table.csv", index=False)
    print("Wrote parameter_targets.csv and rstar_beta_table.csv")
    print(parameter_table.to_string(index=False))


if __name__ == "__main__":
    main()
\end{lstlisting}
\clearpage
\section{python/build\_final\_html.py}
\begin{lstlisting}[language=Python]
"""Build the self-contained, evaluator-facing final HTML report.

The output embeds every public figure, the main result tables, the full source
code, and a detailed Portuguese narrative. No network connection is required
to read the resulting file.
"""

from __future__ import annotations

import base64
import html
import json
from datetime import date
from pathlib import Path

import pandas as pd


ROOT = Path(__file__).resolve().parents[1]
TABLES = ROOT / "outputs" / "tables"
FIGURES = ROOT / "outputs" / "figures"
OUTPUT = ROOT / "Entrega Final.html"


def esc(value: object) -> str:
    return html.escape(str(value), quote=True)


def fmt(value: object) -> object:
    if pd.isna(value):
        return ""
    if isinstance(value, float):
        if abs(value) >= 100:
            return f"{value:,.2f}"
        if abs(value) >= 1:
            return f"{value:.3f}"
        if value == 0:
            return "0"
        return f"{value:.6f}"
    return value


def csv_table(
    filename: str,
    *,
    columns: list[str] | None = None,
    query: str | None = None,
    max_rows: int | None = None,
    table_class: str = "",
) -> str:
    path = TABLES / filename
    if not path.exists():
        return f'<div class="warning">Tabela ausente: <code>{esc(filename)}</code>.</div>'
    df = pd.read_csv(path)
    if query:
        df = df.query(query)
    if columns:
        df = df[[column for column in columns if column in df.columns]]
    if max_rows:
        df = df.head(max_rows)
    df = df.map(fmt)
    return df.to_html(index=False, border=0, classes=f"data-table {table_class}".strip(), escape=True)


def image(filename: str, alt: str, caption: str, analysis: str = "") -> str:
    path = FIGURES / filename
    if not path.exists():
        return f'<div class="warning">Figura ausente: <code>{esc(filename)}</code>.</div>'
    encoded = base64.b64encode(path.read_bytes()).decode("ascii")
    note = f'<div class="figure-analysis">{analysis}</div>' if analysis else ""
    return f"""
    <figure>
      <img src="data:image/png;base64,{encoded}" alt="{esc(alt)}" loading="lazy">
      <figcaption><strong>{esc(caption)}</strong><br><span>Arquivo: outputs/figures/{esc(filename)}</span></figcaption>
      {note}
    </figure>
    """


def callout(kind: str, title: str, body: str) -> str:
    return f'<aside class="callout {esc(kind)}"><strong>{esc(title)}</strong><div>{body}</div></aside>'


def code_details(path: Path) -> str:
    relative = path.relative_to(ROOT).as_posix()
    try:
        source = path.read_text(encoding="utf-8")
    except UnicodeDecodeError:
        source = path.read_text(encoding="latin-1")
    language = {
        ".py": "python",
        ".m": "matlab",
        ".mod": "dynare",
        ".md": "markdown",
        ".txt": "text",
    }.get(path.suffix.lower(), "text")
    lines = source.count("\n") + 1
    return f"""
    <details class="code-file">
      <summary><code>{esc(relative)}</code><span>{lines} linhas</span></summary>
      <pre><code class="language-{language}">{html.escape(source)}</code></pre>
    </details>
    """


def section(section_id: str, kicker: str, title: str, body: str) -> str:
    return f"""
    <section id="{esc(section_id)}">
      <div class="section-kicker">{esc(kicker)}</div>
      <h2>{title}</h2>
      {body}
    </section>
    """


def read_metadata() -> dict:
    path = ROOT / "data" / "clean" / "dataset_metadata.json"
    return json.loads(path.read_text(encoding="utf-8")) if path.exists() else {}


def build() -> None:
    metadata = read_metadata()
    figure_count = len(list(FIGURES.glob("*.png")))
    scenario_manifest = pd.read_csv(TABLES / "scenario_manifest.csv")
    determinacy = pd.read_csv(TABLES / "scenario_determinacy.csv")
    sign_checks = pd.read_csv(TABLES / "irf_sign_checks.csv")
    n_generated_scenarios = len(scenario_manifest)
    n_models = len(determinacy)
    n_determinate = int(determinacy["status"].eq("determinate_bk_count").sum())
    n_signs = int(sign_checks["impact_sign_matches"].astype(str).str.lower().eq("true").sum())

    nav = [
        ("resumo", "Resumo"),
        ("auditoria", "Aderencia a aula"),
        ("chile", "Chile e dados"),
        ("glossario", "Conceitos e siglas"),
        ("modelo", "Modelo teorico"),
        ("calibracao", "Calibracao"),
        ("solucao", "Dynare e solucao"),
        ("momentos", "Momentos e correlacoes"),
        ("fevd", "FEVD"),
        ("irfs", "IRFs baseline"),
        ("kappa", "Sensibilidade kappa"),
        ("phi", "Sensibilidade phipi"),
        ("rho", "Persistencia rhoi"),
        ("econometria", "Econometria"),
        ("previsoes", "Previsoes"),
        ("macro", "Analise macro aprofundada"),
        ("politica", "Analise monetaria"),
        ("limitacoes", "Limitacoes"),
        ("reproducao", "Reproducao"),
        ("codigo", "Codigo completo"),
    ]

    parts: list[str] = []
    parts.append(
        section(
            "resumo",
            "Entrega final",
            "Politica monetaria no Chile em um modelo Novo-Keynesiano",
            f"""
            <p class="lead">Calibracao, solucao, diagnostico de determinacao, funcoes
            impulso-resposta, decomposicao de variancia, econometria e cenarios condicionais
            em um projeto reproduzivel com <strong>Python, GNU Octave e Dynare</strong>.</p>
            <div class="metric-grid">
              <div class="metric"><span>{metadata.get("observations", 101)}</span><small>trimestres de dados reais</small></div>
              <div class="metric"><span>{n_models}</span><small>modelos Dynare executados</small></div>
              <div class="metric"><span>{n_determinate}/{n_models}</span><small>cenarios principais determinados</small></div>
              <div class="metric"><span>{n_signs}/9</span><small>sinais de impacto corretos</small></div>
            </div>
            {callout("success", "Resultado central",
                "O baseline e determinado e seus sinais de impacto sao coerentes com a teoria. "
                "Ao mesmo tempo, o confronto com os dados mostra que o modelo pequeno subestima "
                "volatilidades e persistencias; por isso, os resultados sao um laboratorio de "
                "mecanismos, nao uma descricao completa ou uma previsao oficial do Chile.")}
            <div class="two-col">
              <div>
                <h3>O que foi feito</h3>
                <ul>
                  <li>Dados oficiais trimestrais do Banco Central de Chile, 2001Q1-2026Q1.</li>
                  <li>Modelo NK linear de tres equacoes resolvido no Dynare 7.1 via Octave 11.1.</li>
                  <li>Baseline e {n_generated_scenarios} cenarios gerados para as grades e alternativas.</li>
                  <li>Duas rotas para <code>rho_i</code> e duas calibracoes de FEVD.</li>
                  <li>Previsao incondicional e condicionamento por juro e por inflacao.</li>
                  <li>Regra de Taylor, NKPC, taxa neutra e moda posterior exploratorias.</li>
                </ul>
              </div>
              <div>
                <h3>Leitura correta</h3>
                <ul>
                  <li>Calibracao nao e estimacao.</li>
                  <li>FEVD calibrada nao e decomposicao historica.</li>
                  <li>AR(1) da TPM nao identifica sozinho a suavizacao estrutural.</li>
                  <li>IRFs podem apresentar <em>overshooting</em> sem invalidar o modelo.</li>
                  <li>Cenarios condicionais sao experimentos internos, nao recomendacoes.</li>
                </ul>
              </div>
            </div>
            """,
        )
    )

    parts.append(
        section(
            "auditoria",
            "Paginas 26-55 da Aula 5",
            "O que o roteiro pediu e como o projeto respondeu",
            f"""
            <p>A revisao foi feita a partir do PDF privado disponivel localmente. O texto abaixo e
            uma sintese original: nenhuma pagina ou passagem foi incorporada ao repositorio.</p>
            <div class="audit-grid">
              <article><b>Taxa neutra</b><span>2%, 3% e 4% a.a.; conversao trimestral e beta.</span></article>
              <article><b>Persistencia</b><span>AR(1) real e alternativa calibrada em 0,80.</span></article>
              <article><b>Phillips</b><span>kappa = 0,07; 0,10; 0,13 com IRFs e custos reais.</span></article>
              <article><b>Regra monetaria</b><span>phipi = 1,3,...,2,2 e mapa ampliado de determinacao.</span></article>
              <article><b>Choques e FEVD</b><span>Alvo didatico e alvo alternativo executados.</span></article>
              <article><b>Previsoes</b><span>Incondicional, juro condicionado e inflacao condicionada.</span></article>
            </div>
            {callout("info", "Quando o roteiro diz "isto ou aquilo"",
                "A entrega nao escolhe silenciosamente uma opcao. Ela testa as duas: rho estimado "
                "e calibrado; FEVD didatica e alternativa; condicionamento pelo juro e pela inflacao.")}
            <h3>Correcoes resultantes da auditoria</h3>
            <ol>
              <li>Os nove sinais de impacto estao corretos, mas algumas respostas cruzam zero no
              trimestre seguinte. Portanto, nao se exige monotonicidade das IRFs.</li>
              <li>Um phipi maior reduz o impacto da inflacao, mas eleva o movimento de produto e juro;
              a soma absoluta da inflacao pode subir por <em>undershooting</em>.</li>
              <li>O comando de previsao condicional de primeira ordem usa choques controlados nao
              antecipados antes de cada periodo. Isso nao e previsao perfeita deterministica.</li>
              <li>Momentos e FEVD precisam ser confrontados com dados e identificados como resultados
              internos da calibracao.</li>
            </ol>
            <p class="source-note">Documento de auditoria:
            <code>docs/assignment_review_pages_26_55.md</code>.</p>
            """,
        )
    )

    parts.append(
        section(
            "chile",
            "Contexto empirico",
            "Por que Chile, quais dados e quais transformacoes",
            f"""
            <p>O Chile combina um regime de metas de inflacao, uma taxa de politica claramente
            identificada e series oficiais acessiveis. A meta e interpretada como inflacao projetada
            de 3% em horizonte de dois anos. A variavel operacional e a <strong>TPM</strong>
            (Tasa de Politica Monetaria). Esse ambiente e adequado para ilustrar o nucleo de um
            modelo NK, embora o carater de economia pequena e aberta seja uma limitacao decisiva.</p>
            <div class="three-col">
              <article class="card"><h3>TPM</h3><p>Serie diaria, convertida em media trimestral e
              em taxa trimestral composta.</p><code>F022.TPM.TIN.D001.NO.Z.D</code></article>
              <article class="card"><h3>IPC</h3><p>Variacao mensal composta dentro do trimestre e
              anualizada para apresentacao.</p><code>F074.IPC.VAR.Z.Z.C.M</code></article>
              <article class="card"><h3>PIB real</h3><p>Log do PIB dessazonalizado; ciclo do filtro HP
              com lambda=1600 usado como hiato.</p><code>F032.PIB.FLU.R.CLP.EP18.Z.Z.1.T</code></article>
            </div>
            {image("data_overview.png", "TPM, inflacao e hiato do produto no Chile",
                "Figura 1 -- Dados trimestrais utilizados",
                "A amostra cobre 101 trimestres. A pandemia produz um hiato extremo em 2020Q2, "
                "enquanto o ciclo inflacionario de 2021-2023 e acompanhado por forte elevacao da TPM. "
                "Esses episodios explicam por que momentos de amostra e estimacoes sao sensiveis ao periodo.")}
            <h3>Proveniencia e unidades</h3>
            <table class="data-table">
              <thead><tr><th>Item</th><th>Valor</th></tr></thead>
              <tbody>
                <tr><td>Instituicao</td><td>{esc(metadata.get("institution", "BCCh"))}</td></tr>
                <tr><td>Periodo</td><td>{esc(metadata.get("period", ""))}</td></tr>
                <tr><td>Observacoes</td><td>{esc(metadata.get("observations", ""))}</td></tr>
                <tr><td>Acesso</td><td>{esc(metadata.get("access_date", ""))}</td></tr>
                <tr><td>Base sintetica?</td><td><strong>{esc(metadata.get("is_synthetic", ""))}</strong></td></tr>
                <tr><td>Transformacoes</td><td>{esc(metadata.get("transformations", ""))}</td></tr>
              </tbody>
            </table>
            {callout("warning", "Hiato do produto nao e observado",
                "O filtro HP separa tendencia e ciclo mecanicamente e sofre vies de ponta. O hiato "
                "resultante e uma proxy, nao uma medida estrutural inequivoca de capacidade ociosa.")}
            """,
        )
    )

    glossary = [
        ("NK", "Novo-Keynesiano: estrutura com expectativas racionais, rigidez nominal e politica monetaria."),
        ("DSGE", "Modelo dinamico estocastico de equilibrio geral; o modelo desta entrega e um nucleo DSGE minimo."),
        ("IS", "Equacao de demanda intertemporal que liga hiato ao juro real ex ante."),
        ("NKPC", "New Keynesian Phillips Curve: curva de Phillips prospectiva."),
        ("Calvo", "Mecanismo em que apenas uma fracao das firmas pode reajustar precos a cada periodo."),
        ("EIS", "Elasticidade intertemporal de substituicao; no modelo e inversamente relacionada a sigma."),
        ("TPM", "Tasa de Politica Monetaria do Banco Central de Chile."),
        ("BCCh", "Banco Central de Chile."),
        ("IPC", "Indice de Precos ao Consumidor."),
        ("PIB", "Produto Interno Bruto."),
        ("IRF", "Impulse Response Function: trajetoria apos um choque isolado."),
        ("FEVD", "Forecast Error Variance Decomposition: parcela da variancia atribuida a cada choque."),
        ("BK", "Condicoes de Blanchard-Kahn para existencia e unicidade da solucao racional."),
        ("MQO/OLS", "Minimos Quadrados Ordinarios / Ordinary Least Squares."),
        ("HAC", "Matriz de covariancia robusta a heterocedasticidade e autocorrelacao."),
        ("Newey-West", "Estimador HAC usado nos erros-padrao das regressoes."),
        ("VI/IV", "Variaveis instrumentais, usadas quando regressores podem ser endogenos."),
        ("2SLS", "Minimos Quadrados em Dois Estagios."),
        ("MAP", "Moda a posteriori: ponto de maior densidade posterior."),
        ("Laplace", "Aproximacao local gaussiana da incerteza em torno da moda."),
        ("MCMC", "Simulacao de Monte Carlo via cadeias de Markov para a distribuicao posterior."),
        ("HPD", "Intervalo de maior densidade posterior."),
        ("HP", "Filtro Hodrick-Prescott, aqui com lambda 1600 para dados trimestrais."),
        ("UIP", "Paridade descoberta de juros, ausente no modelo fechado."),
        ("SVAR", "Vetor autorregressivo estrutural, alternativa empirica para identificar choques."),
        ("p.b.", "Ponto-base: 0,01 ponto percentual."),
    ]
    glossary_html = "".join(
        f"<dt>{esc(term)}</dt><dd>{esc(definition)}</dd>" for term, definition in glossary
    )
    parts.append(
        section(
            "glossario",
            "Base conceitual",
            "Glossario de conceitos, parametros e siglas",
            f"""
            <dl class="glossary">{glossary_html}</dl>
            <h3>Parametros e objetos do modelo</h3>
            <div class="parameter-grid">
              <article><b>beta</b><span>Fator de desconto. Quanto maior, maior o peso do futuro.</span></article>
              <article><b>sigma</b><span>Inverso da EIS. Regula a sensibilidade da demanda ao juro real.</span></article>
              <article><b>kappa</b><span>Inclinacao da NKPC. Transmissao do hiato para inflacao.</span></article>
              <article><b>rho<sub>i</sub></b><span>Persistencia/suavizacao da taxa de politica.</span></article>
              <article><b>phi<sub>pi</sub></b><span>Resposta de longo prazo do juro a inflacao.</span></article>
              <article><b>phi<sub>x</sub></b><span>Resposta de longo prazo do juro ao hiato.</span></article>
              <article><b>r*</b><span>Taxa real neutra compativel com equilibrio no modelo.</span></article>
              <article><b>e<sub>x</sub></b><span>Choque de demanda ou preferencia.</span></article>
              <article><b>e<sub>pi</sub></b><span>Choque de custo/markup.</span></article>
              <article><b>e<sub>i</sub></b><span>Inovacao monetaria nao explicada pela regra.</span></article>
            </div>
            <h3>Dinamica e solucao</h3>
            <p><strong>Variavel predeterminada</strong> tem seu valor corrente herdado do passado;
            aqui, <code>i(-1)</code>. <strong>Variaveis de salto</strong> podem reagir imediatamente
            as noticias; aqui, <code>x</code> e <code>pi</code>. Um <strong>autovalor estavel</strong>
            tem modulo menor que um. A determinacao exige que o numero de raizes instaveis seja
            igual ao numero de variaveis prospectivas.</p>
            <p><strong>Momento</strong> e uma caracteristica da distribuicao, como media, variancia,
            desvio-padrao, correlacao ou autocorrelacao. <strong>Convergencia</strong> descreve o
            retorno a vizinhanca do estado estacionario. <strong>Overshooting</strong> ou
            <strong>undershooting</strong> ocorre quando a variavel cruza seu nivel de equilibrio
            antes de convergir.</p>
            """,
        )
    )

    parts.append(
        section(
            "modelo",
            "Estrutura teorica",
            "As tres equacoes e a transmissao monetaria",
            """
            <div class="equation">x<sub>t</sub> = E<sub>t</sub>x<sub>t+1</sub>
            ? (1/sigma)[i<sub>t</sub> ? E<sub>t</sub>pi<sub>t+1</sub> ? r*] + e<sub>x,t</sub></div>
            <div class="equation">pi<sub>t</sub> = betaE<sub>t</sub>pi<sub>t+1</sub>
            + kappax<sub>t</sub> + e<sub>pi,t</sub></div>
            <div class="equation">i<sub>t</sub> = rho<sub>i</sub>i<sub>t?1</sub>
            + (1?rho<sub>i</sub>)[r* + phi<sub>pi</sub>pi<sub>t</sub>
            + phi<sub>x</sub>x<sub>t</sub>] + e<sub>i,t</sub></div>
            <div class="three-col">
              <article class="card"><h3>Curva IS</h3><p>Se o juro real ex ante supera r*, adiar
              consumo fica relativamente atraente e a demanda corrente cai. Expectativas futuras
              afetam a decisao hoje.</p></article>
              <article class="card"><h3>Curva de Phillips</h3><p>Firmas que nao reajustam precos
              tornam a inflacao dependente do custo marginal, aproximado pelo hiato, e da inflacao
              esperada.</p></article>
              <article class="card"><h3>Regra de Taylor</h3><p>O banco central reage a inflacao e
              atividade, mas suaviza alteracoes da taxa. O principio de Taylor requer resposta
              nominal suficiente para elevar o juro real.</p></article>
            </div>
            <h3>Cadeias de transmissao</h3>
            <ul>
              <li><b>Demanda positiva:</b> x sobe -> pi sobe -> regra eleva i -> juro real freia x.</li>
              <li><b>Custo positivo:</b> pi sobe e x cai sob aperto -> surge o trade-off de estabilizacao.</li>
              <li><b>Aperto monetario:</b> i sobe -> juro real sobe -> x cai -> pi cai.</li>
            </ul>
            <h3>Estado estacionario e unidades</h3>
            <p>Na representacao do arquivo principal, <code>x=0</code>, <code>pi=0</code> e
            <code>i=rstar</code>. As IRFs sao convertidas para pontos percentuais. Para previsoes em
            nivel, inflacao recebe de volta a meta de 3% e o juro recebe o neutro nominal de 6,09%,
            composto a partir de r* real de 3% e meta de inflacao de 3%.</p>
            """,
        )
    )

    parts.append(
        section(
            "calibracao",
            "Parametro ? alvo",
            "Calibracao, taxa neutra, beta e persistencia",
            f"""
            <h3>Mapa de parametros</h3>
            {csv_table("parameter_targets.csv")}
            <h3>Taxa neutra anual, taxa trimestral e beta</h3>
            {csv_table("rstar_beta_table.csv")}
            <p>A conversao e <code>rstar_q = (1+rstar_annual)^(1/4)-1</code> e
            <code>beta = 1/(1+rstar_q)</code>. Elevar r* de 2% para 4% reduz beta de 0,99506 para
            0,99024. Esses sao cenarios; nao sao tres estimativas concorrentes da taxa neutra chilena.</p>
            <h3>AR(1) da TPM</h3>
            {csv_table("rhoi_estimate.csv", columns=[
                "rho_i_estimate", "rho_i_se_hac", "rho_i_ci95_low", "rho_i_ci95_high",
                "half_life_quarters", "r_squared", "observations", "sample_start", "sample_end"
            ])}
            {callout("warning", "Persistencia reduzida nao e parametro estrutural puro",
                "O valor 0,934 descreve a inercia da TPM observada. Mudancas de regime, movimentos "
                "do neutro e resposta sistematica a inflacao e atividade tambem afetam o AR(1).")}
            """,
        )
    )

    parts.append(
        section(
            "solucao",
            "Octave + Dynare",
            "Pre-processamento, estado estacionario e Blanchard-Kahn",
            f"""
            <p>O Dynare pre-processa o arquivo <code>.mod</code>, monta as equacoes estaticas e
            dinamicas, calcula o estado estacionario, lineariza o sistema e aplica uma decomposicao
            generalizada para obter a solucao racional de primeira ordem. O comando
            <code>check;</code> reporta os autovalores; <code>stoch_simul(order=1, irf=20);</code>
            calcula momentos, FEVD e IRFs.</p>
            <div class="metric-grid">
              <div class="metric"><span>0,656</span><small>raiz estavel dominante</small></div>
              <div class="metric"><span>1,040</span><small>menor raiz instavel</small></div>
              <div class="metric"><span>2</span><small>raizes instaveis</small></div>
              <div class="metric"><span>2</span><small>variaveis forward-looking</small></div>
            </div>
            {callout("success", "Condicao de Blanchard-Kahn satisfeita",
                "Ha exatamente duas raizes fora do circulo unitario para x e pi, as duas variaveis "
                "prospectivas. O baseline possui solucao local unica e estavel.")}
            {image("determinacy_map.png", "Mapa de determinacao por phi pi",
                "Figura 2 -- Determinacao na grade ampliada de phipi",
                "A fronteira numerica ocorre em torno de phipi=1 para esta parametrizacao. A grade "
                "principal de 1,3 a 2,2 e inteiramente determinada; o ganho marginal passa a ser "
                "avaliado pela dinamica e pelo custo em produto e juros.")}
            <h3>Diagnostico de todos os cenarios</h3>
            {csv_table("scenario_determinacy.csv", columns=[
                "scenario", "status", "n_forward", "n_unstable",
                "dominant_stable_modulus", "convergence_half_life_quarters",
                "smallest_unstable_modulus"
            ])}
            """,
        )
    )

    parts.append(
        section(
            "momentos",
            "Ajuste quantitativo",
            "Momentos teoricos, dados e correlacoes",
            f"""
            {image("moments_model_vs_data.png", "Momentos do modelo comparados aos dados",
                "Figura 3 -- Volatilidade e autocorrelacao",
                "O modelo subestima fortemente a volatilidade. Mesmo retirando 2020Q2-2021Q4, "
                "o hiato observado continua mais volatil. A TPM empirica tambem e bem mais "
                "persistente do que o juro endogeno do modelo.")}
            {csv_table("moments_model_vs_data.csv")}
            <h3>Correlacoes contemporaneas</h3>
            {csv_table("correlations_model_vs_data.csv")}
            <p>O baseline produz correlacao negativa entre juro e hiato porque um aperto contrai
            atividade. Nos dados, a correlacao incondicional e positiva: o banco central tambem
            sobe juros quando a economia e a inflacao estao fortes. Essa inversao ilustra por que
            correlacao observada nao identifica causalidade monetaria.</p>
            {callout("info", "O objetivo nao e maximizar ajuste",
                "A calibracao de choques prioriza uma FEVD transparente e IRFs plausiveis. Uma "
                "estimacao completa exigiria mais friccoes, observaveis, regimes e tratamento de outliers.")}
            """,
        )
    )

    parts.append(
        section(
            "fevd",
            "Decomposicao de variancia",
            "Duas metas de FEVD e o que elas significam",
            f"""
            {image("fevd_calibration_comparison.png", "FEVD calibrada em duas metas",
                "Figura 4 -- Meta versus FEVD alcancada",
                "Com tres desvios-padrao nao se ajustam livremente nove parcelas. A calibracao "
                "minimiza o erro conjunto. A matriz alternativa tem RMSE de 0,95 p.p.; a didatica, "
                "1,82 p.p. Ambas reproduzem a narrativa de inflacao dominada por custo e juro "
                "dominado por inovacao monetaria.")}
            <h3>Desvios-padrao resultantes</h3>
            {csv_table("shock_sigmas_comparison.csv")}
            <h3>Alvo e resultado, celula por celula</h3>
            {csv_table("fevd_calibration_comparison.csv")}
            <h3>FEVD baseline</h3>
            {csv_table("fevd_summary.csv")}
            <p>No baseline, o hiato e explicado em aproximadamente 45% por demanda e 45% por
            politica; a inflacao, em 80% pelo choque de custo; o juro, em 75% pelo choque monetario.
            Isso valida o procedimento de calibracao, mas nao demonstra que esses percentuais
            descrevem historicamente o Chile.</p>
            """,
        )
    )

    parts.append(
        section(
            "irfs",
            "Mecanismos dinamicos",
            "IRFs do baseline: todos os choques e todas as variaveis",
            f"""
            {image("irf_baseline_all_shocks.png", "Grade 3 por 3 de IRFs do baseline",
                "Figura 5 -- Respostas a demanda, custo e politica monetaria",
                "Todos os sinais de impacto correspondem a teoria. Demanda positiva eleva x, pi e i; "
                "custo positivo eleva pi e i e reduz x; aperto monetario eleva i e reduz x e pi. "
                "As reversoes de sinal no horizonte 1 em algumas respostas sao overshooting, nao erro.")}
            <h3>Checagem automatica dos sinais</h3>
            {csv_table("irf_sign_checks.csv")}
            <div class="three-col">
              <article class="card"><h3>Choque de demanda</h3><p>Impactos: x +0,400 p.p.,
              pi +0,021 p.p., i +0,016 p.p. A reacao da politica faz x e pi cruzarem zero no
              trimestre seguinte.</p></article>
              <article class="card"><h3>Choque de custo</h3><p>Impactos: pi +0,236 p.p.,
              x ?0,144 p.p., i +0,022 p.p. E o trade-off de estabilizacao: reduzir inflacao
              requer contracao real.</p></article>
              <article class="card"><h3>Choque monetario</h3><p>Impactos: i +0,048 p.p.,
              x ?0,307 p.p., pi ?0,088 p.p. A inflacao cai imediatamente na solucao prospectiva.</p></article>
            </div>
            <h3>Forma reduzida da politica de decisao</h3>
            {csv_table("policy_transition_coefficients.csv")}
            <p>Os coeficientes unitarios explicam como o estado <code>i(-1)</code> e cada inovacao
            movem as variaveis. A coluna calibrada multiplica cada resposta pelo desvio-padrao do
            choque e coincide com o impacto das IRFs.</p>
            """,
        )
    )

    parts.append(
        section(
            "kappa",
            "Rigidez de precos",
            "Sensibilidade a kappa = 0,07; 0,10; 0,13",
            f"""
            {image("irf_kappa_comparison.png", "Inflacao sob tres valores de kappa",
                "Figura 6 -- Resposta da inflacao a um choque de demanda",
                "Quanto maior kappa, mais rapidamente a atividade se converte em inflacao. O impacto "
                "sobe de 0,0145 para 0,0279 p.p. quando kappa passa de 0,07 para 0,13.")}
            {image("irf_kappa_tradeoffs.png", "Trade-offs de kappa",
                "Figura 7 -- Efeitos de kappa em diferentes choques",
                "Apos choque de custo, kappa maior reduz o impacto inflacionario e o custo real "
                "acumulado. Apos demanda ou politica, porem, precos mais flexiveis transmitem mais "
                "fortemente a atividade para inflacao.")}
            <ul>
              <li>Inflacao de impacto apos demanda: 0,0145 -> 0,0212 -> 0,0279 p.p.</li>
              <li>Inflacao de impacto apos custo: 0,2451 -> 0,2357 -> 0,2276 p.p.</li>
              <li>Custo real acumulado apos custo: 0,4627 -> 0,4187 -> 0,3851 p.p.</li>
              <li>Raiz estavel dominante: 0,6869 -> 0,6561 -> 0,6301.</li>
            </ul>
            <p>A leitura monetaria correta depende do choque. Uma Phillips mais inclinada facilita
            a desinflacao induzida pela contracao do hiato, mas tambem faz choques de demanda
            aparecerem mais rapidamente nos precos.</p>
            """,
        )
    )

    parts.append(
        section(
            "phi",
            "Principio de Taylor",
            "Sensibilidade a phipi de 1,3 a 2,2",
            f"""
            {image("irf_phi_pi_comparison.png", "Inflacao sob diferentes phi pi",
                "Figura 8 -- Resposta da inflacao ao choque de custo",
                "A reacao mais forte reduz o impacto inicial, mas produz maior reversao posterior. "
                "Por isso, o pico e a soma absoluta contam historias diferentes.")}
            {image("irf_phi_pi_tradeoffs.png", "Trade-offs de phi pi",
                "Figura 9 -- Inflacao, produto, juro e convergencia",
                "Entre 1,3 e 2,2, o impacto da inflacao cai de 0,2423 para 0,2303 p.p. e a "
                "convergencia acelera pouco. Em troca, os movimentos acumulados de juro e hiato "
                "sobem de 0,0501 para 0,0796 e de 0,3519 para 0,4734 p.p.")}
            {callout("warning", "Mais reacao nao melhora toda metrica",
                "A soma absoluta da inflacao cresce ligeiramente, de 0,3107 para 0,3167 p.p., "
                "devido ao undershooting. A conclusao robusta e menor impacto inicial com maior "
                "custo real e financeiro, nao estabilizacao uniforme.")}
            <p>Todos os valores da grade principal satisfazem BK. O mapa ampliado mostra
            indeterminacao abaixo do limiar numerico em torno de um. Esse limiar depende de toda a
            parametrizacao, nao apenas de uma regra verbal sobre phipi.</p>
            """,
        )
    )

    parts.append(
        section(
            "rho",
            "Gradualismo",
            "rhoi estimado em 0,934 versus calibrado em 0,80",
            f"""
            {image("irf_rho_comparison.png", "Comparacao de IRFs por rho i",
                "Figura 10 -- Persistencia estimada e calibrada",
                "A maior persistencia estende os efeitos do choque monetario. No baseline, os "
                "efeitos acumulados absolutos sobre inflacao e hiato sao aproximadamente quatro "
                "e tres vezes os do cenario com rho=0,80.")}
            <table class="data-table">
              <thead><tr><th>Resposta ao choque monetario</th><th>rho=0,934</th><th>rho=0,80</th></tr></thead>
              <tbody>
                <tr><td>Juro, soma absoluta</td><td>0,1388 p.p.</td><td>0,0976 p.p.</td></tr>
                <tr><td>Inflacao, soma absoluta</td><td>0,2555 p.p.</td><td>0,0592 p.p.</td></tr>
                <tr><td>Hiato, soma absoluta</td><td>0,8911 p.p.</td><td>0,2777 p.p.</td></tr>
                <tr><td>Meia-vida da raiz estavel</td><td>1,645 trim.</td><td>1,107 trim.</td></tr>
              </tbody>
            </table>
            <p>Esta e uma das analises "ou/ou" executadas dos dois modos. O valor estimado e mais
            aderente a persistencia bruta da TPM, enquanto 0,80 separa melhor uma calibracao
            estrutural convencional. A diferenca entre eles e economicamente grande.</p>
            """,
        )
    )

    parts.append(
        section(
            "econometria",
            "Evidencia complementar",
            "Regra de Taylor, NKPC, r* e moda posterior",
            f"""
            <h3>Regra de Taylor reduzida</h3>
            {csv_table("taylor_rule_estimates.csv")}
            <p>MQO-HAC sugere phipi=0,83; VI/2SLS-HAC, 3,05. O primeiro estagio da inflacao tem
            F=5,82, sinal de instrumento potencialmente fraco. A divergencia nao autoriza escolher
            mecanicamente a estimativa mais compativel com a teoria; ela revela identificacao fragil.</p>
            <h3>Curva de Phillips</h3>
            {csv_table("nkpc_estimates.csv")}
            <p>A especificacao backward encontra kappa~=0,098, proximo da calibracao. A versao forward
            apresenta sinal negativo, erro-padrao elevado e primeiro estagio fraco. Isso nao valida
            empiricamente a NKPC prospectiva; apenas mostra que a grade escolhida e compativel com
            uma relacao reduzida backward.</p>
            <h3>Referencias empiricas para r*</h3>
            {csv_table("rstar_estimates.csv")}
            <p>A tendencia HP do juro real e uma proxy descritiva, sujeita a vies de ponta. Uma
            estimativa estrutural de r* exigiria espaco de estados e dinamica de produto potencial.</p>
            <h3>Moda posterior do DSGE</h3>
            {csv_table("bayesian_estimates.csv")}
            {callout("warning", "MAP nao e posterior completo",
                "Como mh_replic=0, nao ha cadeias MCMC. Os intervalos sao aproximacoes locais de "
                "Laplace em torno da moda e podem ser pouco informativos fora dessa vizinhanca.")}
            """,
        )
    )

    parts.append(
        section(
            "previsoes",
            "Exercicios prospectivos",
            "Previsao incondicional e duas formas de condicionamento",
            f"""
            <p>A previsao e gerada pelo proprio Dynare, como nas paginas 48-55 da aula. O comando
            <code>estimation(datafile=..., mode_compute=0, forecast=8)</code> mantem a calibracao
            fixa, roda o filtro/suavizador de Kalman sobre o historico e projeta oito trimestres a
            frente, devolvendo <code>oo_.forecast.Mean</code> e as bandas <code>HPDinf/HPDsup</code>.
            Os observaveis sao desvios da media amostral; aqui as series voltam ao nivel anualizado
            (TPM media 4,09% a.a.; inflacao media 3,79% a.a.; meta oficial de 3%).</p>
            {image("forecast_fanchart.png", "Fan chart da previsao incondicional",
                "Figura 11 -- Cenario incondicional (Dynare) e bandas HPD",
                "Partindo do estado suavizado do ultimo trimestre, o modelo projeta inflacao de "
                "3,30% no primeiro horizonte convergindo para 3,76% no oitavo, e a TPM caindo de "
                "4,36% para ~4,11% (a media amostral). As bandas sao as faixas de credibilidade do "
                "proprio Dynare, dadas as variancias dos choques calibrados.")}
            {image("conditional_scenarios.png", "Cenarios condicionais para juro e inflacao",
                "Figura 12 -- Caminhos condicionados (conditional_forecast)",
                "As duas formas de condicionamento da aula: controlar a TPM (manter em 4,5% por 2 "
                "trimestres; apertar para ~5,1% por 4) e controlar a inflacao (forcar a meta de 3% "
                "por 4 trimestres). Os caminhos estao ancorados no estado atual via a forma reduzida "
                "exata, identica a solucao do Dynare.")}
            {image("conditional_policy_shocks.png", "Choques monetarios implicitos",
                "Figura 13 -- Inovacoes e_i necessarias para cumprir as restricoes",
                "Os choques implicitos tornam explicito o custo de impor cada caminho: o aperto de "
                "~5,1% exige a maior inovacao inicial; forcar a inflacao a meta de 3% exige um "
                "aperto pontual mais modesto.")}
            <h3>Resumo numerico dos horizontes 1, 4 e 8</h3>
            {csv_table("forecast_summary.csv", columns=[
                "horizon", "date", "scenario", "variable", "value_level",
                "p05_level", "p95_level", "implied_e_i_q"
            ])}
            {callout("danger", "Nao interpretar como previsao oficial",
                "O modelo nao carrega separadamente o ultimo x e a ultima pi observados (o unico "
                "estado e a taxa passada), nao contem cambio, cobre ou condicoes externas e nao "
                "incorpora julgamento do BCCh. Por nao ter inercia de inflacao, a projecao reverte "
                "rapido a media -- e um laboratorio mecanico, nao previsao perfeita.")}
            """,
        )
    )

    parts.append(
        section(
            "macro",
            "Extensoes e validacao",
            "Historia macro, evidencia empirica e mecanismos ausentes no baseline",
            f"""
            <p class="lead">Este bloco amplia o trabalho obrigatorio em tres direcoes: usa o
            suavizador para organizar a historia macro; confronta o DSGE com benchmarks empiricos;
            e reespecifica o modelo para inercia inflacionaria e economia pequena e aberta.</p>

            <h3>1. Decomposicao historica e contrafactual</h3>
            {image("history_shock_decomposition.png", "Decomposicao historica dos observaveis",
                "Figura 14 -- Contribuicoes historicas dos tres choques",
                "Como ha tres observaveis, tres choques e nenhum erro de medida, a reconstrucao e "
                "quase exata. Isso valida o codigo, mas nao identifica causalmente a historia.")}
            {image("history_smoothed_shocks.png", "Choques recuperados pelo suavizador",
                "Figura 15 -- Inovacoes estruturais suavizadas",
                "Crises exigem inovacoes muito maiores que os desvios-padrao ilustrativos do "
                "baseline, pois poucos choques absorvem toda a volatilidade observada.")}
            {image("history_counterfactual.png", "Contrafactual com reacao mais agressiva",
                "Figura 16 -- Mesmos choques sob phi_pi = 2,5",
                "A regra hawkish muda inflacao, atividade e juros. E um exercicio condicional "
                "sujeito a critica de Lucas, nao uma avaliacao causal do BCCh.")}

            <h3>2. SVAR, r* e desempenho preditivo</h3>
            {image("svar_vs_dsge_irf.png", "IRFs monetarias do SVAR e DSGE",
                "Figura 17 -- Evidencia reduzida versus restricoes estruturais",
                "O SVAR passa o teste de brancura, mas rejeita normalidade e apresenta activity "
                "puzzle inicial. A divergencia informa sobre identificacao e especificacao.")}
            {csv_table("svar_diagnostics.csv")}
            {image("rstar_time_varying.png", "Proxy estatistica de r-star",
                "Figura 18 -- Componente local do juro real ex ante",
                "A estimativa termina perto de 0,85% a.a. A amplitude da serie revela incerteza e "
                "contaminacao ciclica; nao e uma estimacao Laubach-Williams estrutural.")}
            {csv_table("rstar_time_varying_summary.csv")}
            {image("forecast_oos_comparison.png", "Erros de previsao fora da amostra",
                "Figura 19 -- NK versus AR(1), passeio aleatorio e VAR",
                "O NK e fraco a um trimestre, mas ganha do passeio aleatorio para inflacao, hiato "
                "e marginalmente juros em quatro trimestres. O teste usa dados revisados.")}
            {csv_table("forecast_oos_metrics.csv", columns=[
                "horizon", "variable", "model", "rmse", "relative_rmse_vs_random_walk"
            ])}

            <h3>Nucleos da inflacao e mercado de trabalho (FRED)</h3>
            {image("ipc_cores_panel.png", "IPC por nucleo",
                "Inflacao por nucleo vs a banda de 2-4%",
                "Energia muito volatil; nucleo e servicos persistentes e acima da banda no surto "
                "recente. Componentes OECD/FRED ate ~2023.")}
            {image("ipc_cores_drivers.png", "Quem mais empurrou a inflacao",
                "Drivers do surto de 2021-23 por nucleo",
                "A energia liderou (pico de 24% a.a.), a frente da cheia (14%), servicos (12%) e "
                "nucleo (11%). Energia = choque de custo; nucleo/servicos = inercia da hibrida.")}
            {image("labor_activity_dashboard.png", "Mercado de trabalho e atividade",
                "Desemprego, emprego, participacao e o hiato do modelo",
                "Emprego e participacao em maximas; desemprego ainda 8,7%; hiato (Kalman) ~=0. Sem "
                "dados de vagas/admissoes para o Chile no FRED. Okun: desemprego inverso ao hiato.")}
            {image("labor_spider.png", "Spider do mercado de trabalho",
                "Spider estilo FED Atlanta (percentis vs historia)",
                "Forte em emprego (92o percentil) e participacao (76o); fraco em desemprego (34o) e "
                "atividade (32o); atual vs pre-pandemia.")}

            <h3>3. Economia aberta e inercia inflacionaria</h3>
            {image("open_economy_data.png", "Cambio e cobre usados na extensao",
                "Figura 20 -- Canais externos relevantes para o Chile",
                "Depreciacao do peso e ciclo do cobre ancoram persistencias e volatilidades. Os "
                "demais coeficientes estruturais continuam sendo calibracoes ilustrativas.")}
            {image("open_economy_irfs.png", "IRFs do modelo de economia aberta",
                "Figura 21 -- Depreciacao, cobre e transmissao monetaria",
                "Depreciacao eleva inflacao por pass-through; cobre positivo expande demanda e "
                "pressiona precos e juros. O modelo satisfaz Blanchard-Kahn.")}
            {image("open_vs_closed_monetary_irf.png", "Modelo aberto versus fechado",
                "Figura 22 -- O setor externo altera a transmissao",
                "O cambio cria uma margem adicional mesmo para choques monetarios, mostrando o "
                "custo analitico de tratar o Chile como economia fechada.")}
            {image("hybrid_nkpc_irfs.png", "IRFs com Phillips hibrida",
                "Figura 23 -- Indexacao e persistencia inflacionaria",
                "Com gamma_pi=0,35, a autocorrelacao modelada da inflacao sobe de cerca de 0,05 "
                "para 0,38. O ajuste melhora, mas a indexacao e calibrada.")}

            <h3>4. Posterior completa por MCMC</h3>
            {image("mcmc_posterior_diagnostics.png", "Posteriores e R-hat do MCMC",
                "Figura 24 -- Incerteza parametrica e convergencia",
                "Duas cadeias de 10.000 propostas, com 30% descartado. R-hat fica abaixo de 1,05, "
                "mas ESS modesto e aceitacao alta recomendam cautela.")}
            {csv_table("mcmc_posterior.csv")}
            {csv_table("mcmc_diagnostics.csv")}

            <h3>5. Comparacao por evidencia marginal</h3>
            {image("bayesian_model_comparison.png", "Evidencia marginal dos modelos de Phillips",
                "Figura 25 -- NKPC prospectiva versus NKPC hibrida",
                "Com os mesmos 101 trimestres, observaveis e priors comuns, a aproximacao de "
                "Laplace favorece a NKPC hibrida por 12,02 log-pontos. O fator de Bayes aproximado "
                "e cerca de 166 mil, mesmo apos a penalizacao do parametro gamma_pi.")}
            {csv_table("bayesian_model_comparison.csv")}
            {callout("info", "Como ler a comparacao",
                "A evidencia marginal integra ajuste e complexidade. O resultado e formal, mas "
                "local ao modo posterior e condicionado aos priors. Nao e bridge sampling nem "
                "uma estimativa de evidencia baseada em cadeias completas para os dois modelos.")}

            {callout("warning", "Hierarquia das conclusoes",
                "IRFs do DSGE sao mecanismos condicionais; SVAR depende da identificacao; r* e "
                "proxy estatistica; a decomposicao historica e exata por construcao; modelo aberto "
                "e NKPC hibrida sao calibracoes; MCMC e posterior condicionada ao modelo e priors.")}
            """,
        )
    )

    parts.append(
        section(
            "politica",
            "Sintese economica",
            "Revisao aprofundada da analise monetaria",
            """
            <h3>1. Determinacao e credibilidade</h3>
            <p>A determinacao requer que expectativas privadas sejam ancoradas por uma regra
            suficientemente ativa. No baseline, phipi=1,75 coloca o sistema acima da fronteira BK.
            Entretanto, a evidencia reduzida para o Chile e ampla: MQO, VI e MAP nao entregam um
            mesmo numero. A conclusao correta e que a calibracao representa um regime ativo
            plausivel, nao uma estimativa pontual incontroversa da reacao do BCCh.</p>

            <h3>2. Gradualismo e horizonte de transmissao</h3>
            <p>O AR(1) elevado faz os movimentos da TPM durarem e amplia o efeito acumulado de uma
            inovacao monetaria. Isso e consistente com decisoes graduais e comunicacao por
            trajetorias, mas o AR(1) tambem absorve mudancas lentas do neutro e persistencia dos
            proprios fundamentos. A comparacao com rho=0,80 mostra que a escolha altera
            substancialmente a transmissao.</p>

            <h3>3. Choques de custo e o trade-off</h3>
            <p>Quando a inflacao sobe por custo, elevar juros reduz demanda, mas nao remove
            diretamente o choque. A inflacao cai ao custo de hiato negativo. Um phipi mais alto reduz
            o impacto inicial e acelera pouco a convergencia, mas aprofunda a contracao. O
            undershooting mostra que "mais agressivo" nao equivale automaticamente a menor
            variabilidade total.</p>

            <h3>4. Rigidez nominal e razao de sacrificio</h3>
            <p>Com kappa maior, uma dada contracao do hiato produz mais desinflacao; por isso, o custo
            real do choque de custo cai. Mas a mesma inclinacao faz expansoes de demanda aparecerem
            mais rapidamente nos precos. A estrutura da inflacao determina tanto a eficacia quanto
            o custo da politica.</p>

            <h3>5. Taxa neutra e postura da politica</h3>
            <p>A postura nao deve ser avaliada apenas pelo nivel nominal da TPM. E necessario
            compara-la com inflacao esperada e r*. No exercicio prospectivo, r*=3% e meta de 3%
            implicam neutro nominal de 6,09%. Essa hipotese faz 4,5% parecer expansionista. Como r*
            e incerto e variavel no tempo, a conclusao e condicional a calibracao, nao avaliacao
            historica definitiva.</p>

            <h3>6. FEVD, narrativa e identificacao</h3>
            <p>A FEVD foi desenhada para uma narrativa transparente: inflacao majoritariamente por
            custo, taxa por inovacao monetaria e hiato por demanda e politica. Ela organiza o
            experimento, mas nao identifica choques chilenos. Para uma afirmacao historica seriam
            necessarias restricoes de identificacao, dados adicionais e comparacao com SVAR ou
            DSGE estimado.</p>

            <h3>7. Implicacao para o Chile</h3>
            <p>O nucleo NK ajuda a pensar expectativas, gradualismo e trade-offs, mas omite canais
            centrais de uma economia pequena e aberta. Cambio, preco do cobre, premio de risco,
            inflacao importada e demanda externa podem alterar tanto a origem quanto a transmissao
            da inflacao. A recomendacao metodologica e usar este modelo como benchmark e nao como
            sistema completo para prescricao.</p>
            """,
        )
    )

    parts.append(
        section(
            "limitacoes",
            "Escopo e honestidade",
            "O que a entrega nao demonstra",
            """
            <div class="two-col">
              <div>
                <h3>Limitacoes economicas</h3>
                <ul>
                  <li>Economia fechada: sem cambio, UIP, cobre, termos de troca ou risco-pais.</li>
                  <li>Sem indexacao, habitos, capital, salarios rigidos ou friccoes financeiras.</li>
                  <li>r* constante e hiato obtido por filtro estatistico.</li>
                  <li>Uma unica regra linear para regimes monetarios distintos.</li>
                </ul>
              </div>
              <div>
                <h3>Limitacoes empiricas</h3>
                <ul>
                  <li>Covid-19 e um outlier grande para amostra e filtro.</li>
                  <li>Instrumentos fracos na NKPC forward e na inflacao da regra de Taylor.</li>
                  <li>MAP sem MCMC e intervalos apenas locais.</li>
                  <li>Choques calibrados, nao identificados por likelihood no baseline.</li>
                </ul>
              </div>
            </div>
            <h3>Extensoes prioritarias</h3>
            <ol>
              <li>Modelo NK de economia aberta com cambio, inflacao importada e premio de risco.</li>
              <li>r* e produto potencial em espaco de estados.</li>
              <li>Estimacao bayesiana completa com MCMC, erros de medida e tratamento da pandemia.</li>
              <li>Identificacao externa ou SVAR para comparar choques monetarios e de oferta.</li>
              <li>Avaliacao fora da amostra das previsoes e comparacao com projecoes do BCCh.</li>
            </ol>
            """,
        )
    )

    parts.append(
        section(
            "reproducao",
            "Projeto reproduzivel",
            "Arquitetura, comandos e verificacoes",
            f"""
            <div class="pipeline">
              <span>BCCh API</span><b>-></b><span>dados trimestrais</span><b>-></b>
              <span>estimacoes</span><b>-></b><span>calibracao</span><b>-></b>
              <span>modelos .mod</span><b>-></b><span>Octave + Dynare</span><b>-></b>
              <span>tabelas, figuras e HTML</span>
            </div>
            <h3>Execucao completa a partir da raiz</h3>
            <pre><code>python -m pip install -r requirements.txt
python python/build_chile_dataset.py
python python/estimate_rhoi_chile.py
python python/estimate_taylor_rule.py
python python/estimate_nkpc.py
python python/estimate_rstar.py
python python/calibrate_shocks.py
python python/make_tables.py
python python/generate_dynare_models.py
python python/run_dynare_batch.py --all --timeout 300
python python/run_bayesian.py
python python/collect_outputs.py
python python/analyze_determinacy.py
python python/analyze_model_results.py
python python/run_forecast.py
python python/forecast_model.py
python python/plot_irfs.py
python python/run_history.py
python python/plot_history.py
python python/analyze_svar.py
python python/estimate_time_varying_rstar.py
python python/evaluate_forecasts.py
python python/build_open_economy_dataset.py
python python/generate_macro_extension_models.py
python python/run_macro_extensions.py
python python/analyze_macro_extensions.py
python python/plot_mcmc.py
python python/compare_bayesian_models.py --reuse-baseline
python python/build_final_html.py</code></pre>
            <h3>Tratamento do Windows e OneDrive</h3>
            <p>O runner copia os modelos para uma pasta temporaria com caminho ASCII antes de
            iniciar o Octave. Isso evita travamentos do pre-processador do Dynare ao criar e
            renomear pastas em caminhos sincronizados, com espacos e acentos. Resultados CSV sao
            copiados de volta para <code>outputs/dynare/</code>.</p>
            <h3>Checklist</h3>
            <ul class="checklist">
              <li>Dados marcados como reais em <code>dataset_metadata.json</code>.</li>
              <li>19 modelos com saidas Dynare e diagnostico BK.</li>
              <li>{figure_count} figuras publicas incorporadas neste arquivo.</li>
              <li>Tabelas carregadas diretamente dos CSVs produzidos pelo pipeline.</li>
              <li>Codigo integral incorporado abaixo.</li>
              <li>PDF privado e credenciais fora do versionamento.</li>
            </ul>
            """,
        )
    )

    python_files = sorted((ROOT / "python").glob("*.py"))
    dynare_core = sorted(
        path for path in (ROOT / "dynare").iterdir()
        if path.is_file() and path.suffix.lower() in {".m", ".mod"}
    )
    generated_mods = sorted((ROOT / "dynare" / "generated").glob("*.mod"))
    code_blocks = [
        "<h3>Python</h3>",
        *(code_details(path) for path in python_files),
        "<h3>Octave e Dynare: arquivos centrais</h3>",
        *(code_details(path) for path in dynare_core),
        "<h3>Dynare: cenarios gerados</h3>",
        *(code_details(path) for path in generated_mods),
        "<h3>Configuracao</h3>",
        code_details(ROOT / "requirements.txt"),
    ]
    parts.append(
        section(
            "codigo",
            "Apendice tecnico",
            "Codigo-fonte integral usado na entrega",
            """
            <p>Cada bloco e expansivel. Os arquivos sao incorporados como texto no momento da
            geracao do HTML; portanto, este apendice representa exatamente a versao usada para
            produzir as tabelas e figuras desta entrega.</p>
            <div class="code-toolbar">
              <button type="button" onclick="toggleAll(true)">Expandir todos</button>
              <button type="button" onclick="toggleAll(false)">Recolher todos</button>
              <input id="codeSearch" type="search" placeholder="Filtrar arquivos..." oninput="filterCode(this.value)">
            </div>
            """
            + "\n".join(code_blocks),
        )
    )

    nav_html = "".join(f'<a href="#{sid}">{esc(label)}</a>' for sid, label in nav)
    body = "\n".join(parts)
    styles = """
    :root{--ink:#132238;--muted:#5d6b7d;--paper:#f5f1e8;--card:#fffdf8;--navy:#0d3558;
    --blue:#1b668f;--teal:#19756f;--gold:#bd7d16;--red:#a53c35;--line:#d8d1c3;
    --shadow:0 12px 36px rgba(20,35,52,.10)}
    *{box-sizing:border-box}html{scroll-behavior:smooth}body{margin:0;background:var(--paper);
    color:var(--ink);font:16px/1.67 "Segoe UI",Arial,sans-serif}
    header{min-height:76vh;padding:8rem max(6vw,2rem) 5rem;color:white;position:relative;overflow:hidden;
    background:linear-gradient(125deg,#071d31 0%,#0d3f63 52%,#16756f 100%)}
    header:after{content:"";position:absolute;right:-12vw;top:-12vw;width:46vw;height:46vw;
    border:1px solid rgba(255,255,255,.2);border-radius:50%;box-shadow:0 0 0 5vw rgba(255,255,255,.025),
    0 0 0 11vw rgba(255,255,255,.02)}.eyebrow{letter-spacing:.18em;text-transform:uppercase;
    font-weight:700;color:#a6e4dc}h1{font-family:Georgia,serif;font-size:clamp(3rem,7vw,6.8rem);
    max-width:1100px;line-height:1.02;margin:.4rem 0 1.5rem}.subtitle{font-size:1.3rem;max-width:780px;
    color:#dceaf2}.cover-meta{display:flex;gap:1rem;flex-wrap:wrap;margin-top:2rem}
    .cover-meta span{border:1px solid rgba(255,255,255,.35);padding:.45rem .8rem;border-radius:999px}
    nav{position:sticky;top:0;z-index:20;background:rgba(255,253,248,.96);backdrop-filter:blur(12px);
    border-bottom:1px solid var(--line);display:flex;gap:.2rem;overflow-x:auto;padding:.55rem 2vw;
    box-shadow:0 3px 15px rgba(0,0,0,.05)}nav a{white-space:nowrap;color:var(--navy);
    text-decoration:none;padding:.4rem .65rem;border-radius:6px;font-size:.86rem;font-weight:650}
    nav a:hover{background:#e4eee9}main{max-width:1180px;margin:auto;padding:1rem 2rem 6rem}
    section{padding:5rem 0;border-bottom:1px solid var(--line);scroll-margin-top:4rem}
    .section-kicker{text-transform:uppercase;letter-spacing:.16em;color:var(--teal);font-weight:800;
    font-size:.78rem}h2{font:700 clamp(2.2rem,4vw,4rem)/1.08 Georgia,serif;margin:.35rem 0 2rem;
    color:var(--navy)}h3{color:var(--navy);margin-top:2rem}p{max-width:920px}.lead{font-size:1.3rem;
    line-height:1.65}.metric-grid{display:grid;grid-template-columns:repeat(4,1fr);gap:1rem;margin:2rem 0}
    .metric{background:var(--navy);color:white;padding:1.4rem;border-radius:10px;box-shadow:var(--shadow)}
    .metric span{display:block;font:700 2.3rem Georgia,serif;color:#8ed7cc}.metric small{display:block}
    .two-col,.three-col{display:grid;grid-template-columns:repeat(2,1fr);gap:1.2rem;margin:1.5rem 0}
    .three-col{grid-template-columns:repeat(3,1fr)}.card,.audit-grid article,.parameter-grid article{
    background:var(--card);border:1px solid var(--line);border-radius:10px;padding:1.2rem;box-shadow:0 4px 18px rgba(0,0,0,.04)}
    .audit-grid,.parameter-grid{display:grid;grid-template-columns:repeat(3,1fr);gap:.8rem;margin:1.5rem 0}
    .audit-grid b,.parameter-grid b{display:block;color:var(--navy);font-size:1.05rem}
    .audit-grid span,.parameter-grid span{display:block;color:var(--muted);margin-top:.35rem}
    .callout{padding:1.2rem 1.4rem;border-left:5px solid;margin:1.5rem 0;background:white;
    border-radius:4px 10px 10px 4px;box-shadow:0 5px 20px rgba(0,0,0,.05)}
    .callout strong{display:block;font-size:1.05rem;margin-bottom:.25rem}.callout.success{border-color:var(--teal)}
    .callout.info{border-color:var(--blue)}.callout.warning{border-color:var(--gold)}
    .callout.danger{border-color:var(--red)}figure{margin:2rem 0;background:var(--card);padding:1rem;
    border:1px solid var(--line);border-radius:12px;box-shadow:var(--shadow)}figure img{display:block;width:100%;
    height:auto;border-radius:7px}figcaption{padding:1rem .4rem .5rem;color:var(--navy)}
    figcaption span,.source-note{color:var(--muted);font-size:.88rem}.figure-analysis{background:#edf4f1;
    padding:1rem;border-radius:7px;margin:.5rem .3rem;color:#24443f}.data-table{width:100%;border-collapse:collapse;
    display:block;overflow-x:auto;background:var(--card);margin:1rem 0 2rem;font-size:.88rem}
    .data-table th{background:var(--navy);color:white;text-align:left}.data-table th,.data-table td{
    border:1px solid #ddd6c9;padding:.55rem .65rem;white-space:nowrap}.data-table tr:nth-child(even){background:#f1eee7}
    .equation{font:1.18rem Georgia,serif;text-align:center;background:var(--card);border:1px solid var(--line);
    border-left:5px solid var(--teal);padding:1.2rem;margin:.8rem 0;overflow-x:auto}.glossary{display:grid;
    grid-template-columns:minmax(80px,150px) 1fr;gap:.2rem 1rem}.glossary dt{font-weight:800;color:var(--navy);
    padding:.45rem}.glossary dd{margin:0;padding:.45rem;border-bottom:1px dotted var(--line)}
    code{font-family:Consolas,"Courier New",monospace;background:#e9e5dd;padding:.08rem .25rem;border-radius:3px}
    pre{background:#081521;color:#d7e4ec;padding:1.2rem;border-radius:8px;overflow:auto;max-height:70vh;
    line-height:1.45;font-size:.82rem}pre code{background:none;padding:0}.pipeline{display:flex;flex-wrap:wrap;
    align-items:center;gap:.5rem;margin:1.5rem 0}.pipeline span{background:var(--navy);color:white;padding:.6rem .8rem;
    border-radius:6px}.pipeline b{color:var(--gold)}.checklist{list-style:none;padding:0}.checklist li:before{
    content:"?";color:var(--teal);font-weight:900;margin-right:.6rem}.code-file{background:var(--card);
    border:1px solid var(--line);border-radius:7px;margin:.55rem 0}.code-file summary{cursor:pointer;padding:.8rem 1rem;
    display:flex;justify-content:space-between;gap:1rem}.code-file summary span{color:var(--muted);font-size:.85rem}
    .code-file pre{margin:0;border-radius:0 0 7px 7px}.code-toolbar{position:sticky;top:3.8rem;z-index:10;
    display:flex;gap:.5rem;flex-wrap:wrap;background:var(--paper);padding:.7rem 0}.code-toolbar button,
    .code-toolbar input{border:1px solid var(--line);background:white;padding:.55rem .75rem;border-radius:6px}
    .code-toolbar button{cursor:pointer;background:var(--navy);color:white}.warning{padding:1rem;background:#fff0d7;
    border-left:4px solid var(--gold)}footer{background:#071d31;color:#dceaf2;padding:3rem max(5vw,2rem)}
    footer a{color:#8ed7cc}@media(max-width:850px){.metric-grid,.audit-grid,.parameter-grid,.three-col,.two-col{
    grid-template-columns:1fr 1fr}header{min-height:auto;padding-top:5rem}}@media(max-width:560px){
    main{padding:1rem}.metric-grid,.audit-grid,.parameter-grid,.three-col,.two-col{grid-template-columns:1fr}
    h1{font-size:3rem}.glossary{grid-template-columns:1fr}.glossary dd{padding-top:0}}
    @media print{nav,.code-toolbar{display:none}header{min-height:auto;padding:3rem;background:#123!important;
    -webkit-print-color-adjust:exact}section{break-inside:auto}figure,.card,.callout{break-inside:avoid}
    details:not([open])>*:not(summary){display:block}.code-file pre{max-height:none}}
    """

    script = """
    function toggleAll(openState){
      document.querySelectorAll('.code-file').forEach(function(d){d.open=openState;});
    }
    function filterCode(term){
      term=term.toLowerCase();
      document.querySelectorAll('.code-file').forEach(function(d){
        d.style.display=d.querySelector('summary').innerText.toLowerCase().includes(term)?'block':'none';
      });
    }
    const observer=new IntersectionObserver(entries=>{
      entries.forEach(entry=>{
        if(entry.isIntersecting){
          document.querySelectorAll('nav a').forEach(a=>a.classList.remove('active'));
          const link=document.querySelector('nav a[href="#'+entry.target.id+'"]');
          if(link) link.classList.add('active');
        }
      });
    },{rootMargin:'-20% 0px -70% 0px'});
    document.querySelectorAll('main section').forEach(s=>observer.observe(s));
    """

    document = f"""<!doctype html>
<html lang="pt-BR">
<head>
  <meta charset="utf-8">
  <meta name="viewport" content="width=device-width,initial-scale=1">
  <meta name="description" content="Entrega final de politica monetaria: modelo Novo-Keynesiano calibrado para o Chile.">
  <title>Entrega Final | Politica Monetaria no Chile</title>
  <style>{styles}</style>
</head>
<body>
<header>
  <div class="eyebrow">FGV EESP * Politica Monetaria * 2026</div>
  <h1>Entrega Final</h1>
  <p class="subtitle">Modelo Novo-Keynesiano calibrado para o Chile: dados, teoria,
  solucao, evidencia, politica monetaria e codigo reproduzivel.</p>
  <div class="cover-meta">
    <span>Chile</span><span>2001Q1-2026Q1</span><span>Octave 11.1</span>
    <span>Dynare 7.1</span><span>Python</span><span>Gerado em {date.today().strftime("%d/%m/%Y")}</span>
  </div>
</header>
<nav aria-label="Navegacao do relatorio">{nav_html}</nav>
<main>{body}</main>
<footer>
  <strong>Entrega Final -- Politica Monetaria no Chile</strong>
  <p>Dados: Banco Central de Chile. Resultados do modelo: pipeline publico deste projeto.
  Este documento e autocontido e pode ser aberto sem conexao com a internet.</p>
  <p>Referencias oficiais:
  <a href="https://www.bcentral.cl/web/banco-central/areas/politica-monetaria">marco de politica monetaria</a> *
  <a href="https://si3.bcentral.cl/Siete">Base de Datos Estadisticos</a> *
  <a href="https://www.dynare.org/manual/">manual do Dynare</a>.</p>
</footer>
<script>{script}</script>
</body>
</html>"""
    document = "\n".join(line.rstrip() for line in document.splitlines()) + "\n"
    OUTPUT.write_text(document, encoding="utf-8")
    print(f"Wrote {OUTPUT} ({OUTPUT.stat().st_size / 1024 / 1024:.2f} MiB)")
    print(f"Embedded {len(list(FIGURES.glob('*.png')))} PNG figures")
    print(f"Embedded {len(python_files) + len(dynare_core) + len(generated_mods) + 1} source files")


if __name__ == "__main__":
    build()
\end{lstlisting}
\clearpage
\section{python/build\_codigo\_final.py}
\begin{lstlisting}[language=Python]
"""Generate "Codigo Final": ONE PDF document with ALL project code.

Reads every python/*.py and dynare/*.mod, transliterates non-ASCII characters in
the *code* to ASCII (pi, ->, --, accents stripped) so pdflatex + listings never
chokes, and writes entrega_aula5/Codigo_Final.tex. Prose/titles keep their
accents via inputenc utf8. Build with: entrega_aula5/build.sh Codigo_Final

The file order is grouped to read like the pipeline; any file not listed is
appended automatically, so the document always contains *all* the code.
"""

from __future__ import annotations

import datetime as dt
import unicodedata
from pathlib import Path

HERE = Path(__file__).resolve().parent          # python/
ROOT = HERE.parent
MOD = ROOT / "dynare"
OUT = ROOT / "entrega_aula5" / "Codigo_Final.tex"

# Non-ASCII -> ASCII so the code listings are pure ASCII.
GREEK = {
    "pi": "pi", "kappa": "kappa", "beta": "beta", "sigma": "sigma", "phi": "phi", "rho": "rho",
    "alpha": "alpha", "theta": "theta", "lambda": "lambda", "gamma": "gamma", "delta": "delta",
    "epsilon": "epsilon", "eta": "eta", "mu": "mu", "nu": "nu", "tau": "tau", "omega": "omega",
    "Phi": "Phi", "Sigma": "Sigma", "Pi": "Pi", "Delta": "Delta", "Omega": "Omega",
}
SYM = {
    "--": "--", "-": "-", "x": "x", "->": "->", "<-": "<-", "^2": "^2", "^3": "^3",
    "~=": "~=", "<=": "<=", ">=": ">=", "!=": "!=", "*": "*", "...": "...", "-": "-",
    "'": "'", "'": "'", """: '"', """: '"', "deg": "deg", "+/-": "+/-", "inf": "inf",
    "sqrt": "sqrt", "sum": "sum", "in": "in", "d": "d", "==": "==", "Delta": "Delta",
    " ": " ", " ": " ", "": "",
}


def translit(text: str) -> str:
    out = []
    for ch in text:
        if ord(ch) < 128:
            out.append(ch)
        elif ch in GREEK:
            out.append(GREEK[ch])
        elif ch in SYM:
            out.append(SYM[ch])
        else:
            d = unicodedata.normalize("NFKD", ch).encode("ascii", "ignore").decode()
            out.append(d if d else "?")
    return "".join(out)


# Grouped, pipeline-ordered. Remaining files are appended in a final group.
GROUPS: list[tuple[str, list[str]]] = [
    ("Infraestrutura compartilhada", ["common.py"]),
    ("Dados (BCCh / FRED / filtro de Kalman)", [
        "build_chile_dataset.py", "reprocess_dataset.py",
        "build_open_economy_dataset.py", "discover_ipc_series.py"]),
    ("Geracao dos modelos Dynare", [
        "generate_dynare_models.py", "generate_macro_extension_models.py"]),
    ("Modelos Dynare (.mod)", [
        "nk_chile_base.mod", "nk_chile_estim.mod", "nk_chile_hybrid.mod",
        "nk_chile_hybrid_estim.mod", "nk_chile_forecast.mod",
        "nk_chile_history.mod", "nk_chile_history_hybrid.mod",
        "nk_chile_mcmc.mod", "nk_chile_open.mod"]),
    ("Solucao e calibracao", [
        "hybrid_solution.py", "calibrate_shocks.py", "analyze_determinacy.py"]),
    ("Estimacao dos parametros", [
        "estimate_nkpc.py", "estimate_taylor_rule.py", "estimate_rhoi_chile.py",
        "estimate_rstar.py", "estimate_time_varying_rstar.py"]),
    ("Execucao (Octave/Dynare + Bayesiano)", [
        "run_dynare_batch.py", "run_bayesian.py", "run_mcmc.py",
        "run_forecast.py", "run_history.py", "run_macro_extensions.py"]),
    ("Previsao e avaliacao", [
        "forecast_model.py", "evaluate_forecasts.py"]),
    ("Analise (SVAR, Bayes, resultados)", [
        "analyze_svar.py", "compare_bayesian_models.py",
        "analyze_model_results.py", "analyze_macro_extensions.py"]),
    ("Nucleos da inflacao e mercado de trabalho", [
        "build_inflation_cores.py", "build_inflation_cores_bcch.py",
        "build_labor_activity.py"]),
    ("Graficos", [
        "plot_history.py", "plot_ipc_decomposition.py", "plot_forecast_outlook.py",
        "plot_hybrid_test.py", "plot_irfs.py", "plot_mcmc.py"]),
    ("Tabelas, coleta e entregaveis", [
        "collect_outputs.py", "make_tables.py", "build_final_html.py",
        "build_codigo_final.py"]),
]


def resolve(name: str) -> Path:
    return (MOD / name) if name.endswith(".mod") else (HERE / name)


def latex_escape_title(name: str) -> str:
    return name.replace("\\", r"\textbackslash{}").replace("_", r"\_") \
               .replace("&", r"\&").replace("%", r"\%").replace("#", r"\#")


PREAMBLE = r"""\documentclass[9pt]{extarticle}
\usepackage[a4paper,margin=1.8cm]{geometry}
\usepackage[utf8]{inputenc}
\usepackage[T1]{fontenc}
\usepackage{lmodern}
\usepackage{xcolor}
\usepackage{listings}
\usepackage[hidelinks]{hyperref}
\definecolor{cmt}{RGB}{96,139,84}
\definecolor{kw}{RGB}{0,0,170}
\definecolor{str}{RGB}{163,21,21}
\definecolor{ln}{RGB}{150,150,150}
\lstset{
  basicstyle=\ttfamily\scriptsize,
  breaklines=true,breakatwhitespace=false,
  columns=fullflexible,keepspaces=true,
  showstringspaces=false,upquote=true,tabsize=4,
  commentstyle=\color{cmt},keywordstyle=\color{kw},stringstyle=\color{str},
  numbers=left,numberstyle=\tiny\color{ln},numbersep=6pt,xleftmargin=2.2em,
  frame=lines,framesep=4pt,
}
\title{\textbf{C\'odigo Final}\\[2mm]\large Modelo Novo-Keynesiano de Pol\'itica
Monet\'aria para o Chile\\\normalsize Todo o c\'odigo (Python + Dynare) em um
\'unico documento}
\author{Daniel Colli e Jo\~ao Zangrandi}
\date{REPLACE_DATE}
\begin{document}
\maketitle
{\small Este documento re\'une \textbf{todo} o c\'odigo do projeto: TARGET_NFILES
arquivos (Python e Dynare \texttt{.mod}), agrupados na ordem do \emph{pipeline}.
Caracteres n\~ao-ASCII do c\'odigo (letras gregas, acentos) foram transliterados
para ASCII na listagem para garantir a compila\c{c}\~ao; a l\'ogica \'e id\^entica
\`a dos arquivos-fonte do reposit\'orio.\par}
\tableofcontents
\clearpage
"""


def main() -> None:
    used: list[str] = []
    for _, names in GROUPS:
        used.extend(names)
    # Append any file not explicitly grouped.
    extra = sorted(p.name for p in HERE.glob("*.py") if p.name not in used)
    extra += sorted(p.name for p in MOD.glob("*.mod") if p.name not in used)
    groups = GROUPS + ([("Outros", extra)] if extra else [])

    nfiles = sum(len(n) for _, n in groups)
    parts = [PREAMBLE.replace("REPLACE_DATE", dt.date.today().strftime("%d/%m/%Y"))
             .replace("TARGET_NFILES", str(nfiles))]

    for title, names in groups:
        parts.append(f"\\part*{{{title}}}\\addcontentsline{{toc}}{{part}}{{{title}}}\n")
        for name in names:
            path = resolve(name)
            if not path.exists():
                continue
            code = translit(path.read_text(encoding="utf-8", errors="replace"))
            code = code.replace("\t", "    ").rstrip("\n")
            lang = "Python" if name.endswith(".py") else "{}"
            rel = ("dynare/" if name.endswith(".mod") else "python/") + name
            parts.append(f"\\section{{{latex_escape_title(rel)}}}")
            parts.append(f"\\begin{{lstlisting}}[language={lang}]")
            parts.append(code)
            parts.append("\\end{lstlisting}\n\\clearpage")

    parts.append("\\end{document}\n")
    OUT.write_text("\n".join(parts), encoding="utf-8")
    print(f"Wrote {OUT} ({nfiles} files, {OUT.stat().st_size//1024} KB).")


if __name__ == "__main__":
    main()
\end{lstlisting}
\clearpage
\end{document}
